% Production d’écrits utilitaires : rédiger un rapport (rédaction de l’introduction, du corps et de la conclusion) — Français 1re Tech — Cours signé OnBuch+
% Programme officiel MINESEC. Compiler avec : tectonic N11-production-ecrits-utilitaires-rediger-rapport-redaction.tex
\newcommand{\DOCMATIERE}{Français}
\newcommand{\DOCNIVEAU}{1re Tech}
\newcommand{\DOCMODULE}{Français – classe de Première technique}
\newcommand{\DOCLECON}{N11}
\newcommand{\DOCTITRE}{Production d’écrits utilitaires : rédiger un rapport (rédaction de l’introduction, du corps et de la conclusion)}
% OnBuch+ — préambule « cours signés » ENRICHI (compilé avec tectonic / XeLaTeX)
% Système visuel « Pop » d'OnBuch+ : fond crème (jamais blanc), bordures encre
% épaisses, ombres « dures » décalées, titres Archivo Black en capitales.
% Version MATHÉMATIQUES : graphes (pgfplots/tikz), boîtes pédagogiques
% (définition, propriété, méthode, attention, le savais-tu, exercice résolu,
% exercice, TP, bilan), page de garde et signature OnBuch+ ancrée en bas.
%
% Macros attendues dans main.tex AVANT \input{preamble} :
%   \DOCMATIERE  \DOCNIVEAU  \DOCMODULE  \DOCLECON (numéro)  \DOCTITRE
\documentclass[11pt]{article}

\usepackage{fontspec}
\usepackage[french]{babel}
\usepackage[a4paper,margin=2cm,top=3.4cm,bottom=2.8cm,headheight=1.4cm,footskip=1.3cm]{geometry}
\usepackage{amsmath,amssymb,amsfonts}
\usepackage{mathtools} % \xmapsto, \coloneqq... souvent utilisés par les modèles
\usepackage{cancel} % simplifications barrées (bilans de forces)
% \abs{z} (module, valeur absolue) et \norm{u} (norme), utilisés par les modèles
\providecommand{\abs}{}\let\abs\relax\DeclarePairedDelimiter{\abs}{\lvert}{\rvert}
\providecommand{\norm}{}\let\norm\relax\DeclarePairedDelimiter{\norm}{\lVert}{\rVert}
\usepackage[table]{xcolor} % [table] : fournit aussi \rowcolors (alternance de lignes)
\usepackage{pagecolor}
\usepackage{tikz}
\usetikzlibrary{arrows.meta,positioning,calc,shapes.geometric,shapes.misc,shapes.arrows,decorations.pathmorphing,decorations.pathreplacing,decorations.markings,patterns,shadows,mindmap,trees,fit,backgrounds,matrix,intersections,angles,quotes}
\usepackage{pgfplots}
\usepackage[version=4]{mhchem} % équations nucléaires \ce{^{235}_{92}U}
\usepackage[european,siunitx]{circuitikz} % schémas électriques
\pgfplotsset{compat=1.18}
\usepgfplotslibrary{fillbetween,groupplots} % aires entre courbes (calcul intégral)
\usepackage{enumitem}
\usepackage{fancyhdr}
% [most] charge toutes les bibliothèques tcolorbox (listings, minted, raster,
% documentation...) et gonfle inutilement la mémoire TeX sur un document long
% (~300 boîtes) : on ne charge que ce qui est réellement utilisé.
\usepackage[skins,breakable]{tcolorbox}
\usepackage{graphicx}
\usepackage{colortbl}
\usepackage{booktabs}
\usepackage{tabularx}
\usepackage{multirow}
\usepackage{diagbox} % case d'en-tête diagonale des tableaux à double entrée
% Colonnes extensibles centrée / gauche / droite (C, L, R), souvent utilisées par les modèles
\newcolumntype{C}{>{\centering\arraybackslash}X}
\newcolumntype{L}{>{\raggedright\arraybackslash}X}
\newcolumntype{R}{>{\raggedleft\arraybackslash}X}
\usepackage{array}
\usepackage{multicol}
\usepackage{siunitx}
% parse-numbers=false : accepte aussi \SI{2,5 \times 10^{-3}}{...}
\sisetup{locale=FR,output-decimal-marker={,},group-minimum-digits=4,parse-numbers=false}
\DeclareSIUnit\ohm{\ensuremath{\symup{\Omega}}} % U+2126 absent de la police
\DeclareSIUnit\division{div} % réglages d'oscilloscope (ms/div, V/div)
\DeclareSIUnit\tr{tr} % tours (vitesse de rotation, ex. \SI{1500}{\tr\per\minute})
\DeclareSIUnit\molar{M} % concentration molaire (ex. \SI{3}{\molar}, \SI{1}{\micro\molar})
\DeclareSIUnit\an{an} % année (le modèle confond parfois avec \year, un compteur TeX) — ex. \SI{5}{\kilo\meter\per\an}
\DeclareSIUnit\FCFA{FCFA} % franc CFA (ex. \SI{500}{\FCFA\per\kilo\gram})
\DeclareSIUnit\degreeBrix{\textdegree Brix} % teneur en sucre d'un sirop/jus (réfractométrie)
\DeclareSIUnit\month{mois} % le modèle utilise parfois \month, un compteur TeX, comme unité de durée
\DeclareSIUnit\microliter{\micro\liter} % le modèle écrit parfois \microliter en un seul token
\DeclareSIUnit\million{millions} % dénombrement (ex. \SI{15}{\million\per\mL} spermatozoïdes)
\usepackage{qrcode}
% \degree : commande fréquemment utilisée par les modèles pour un angle ou une
% température en dehors de \SI{}{} (non fournie par les paquets chargés).
\providecommand{\degree}{^{\circ}}
% \qed (fin de démonstration), utilisé par les modèles sans amsthm
\providecommand{\qed}{\hfill$\square$}
% \barre : double barre des valeurs interdites dans un tableau de variations (tkz-tab)
\providecommand{\barre}{\ensuremath{\|}}
% environnement proof (amsthm non chargé) utilisé par les modèles
\ifdefined\proof\else\newenvironment{proof}[1][Démonstration]{\par\noindent\textit{#1.}\ }{\qed\par}\fi
\usepackage{lastpage}
\usepackage{hyperref}
\hypersetup{hidelinks}

% ============================================================
% Palette « Pop » OnBuch+ — mêmes hex que la classe P (plus_ui.dart)
% ============================================================
\definecolor{poporange}{HTML}{F59321}
\definecolor{popdark}{HTML}{E07A0C}
\definecolor{popcream}{HTML}{FFF6E8}
\definecolor{popink}{HTML}{1C1714}
\definecolor{poppurple}{HTML}{7A5AE0}
\definecolor{popgreen}{HTML}{2E9E6A}
\definecolor{poppink}{HTML}{E0457B}
\definecolor{popblue}{HTML}{2F7DE1}
\definecolor{popgold}{HTML}{C98A12}
\definecolor{popmuted}{HTML}{8A7F76}
\definecolor{popline}{HTML}{EADFD2}
\definecolor{popgray}{HTML}{8A7F76}
\colorlet{popgrey}{popgray}
% Alias tolérants (le modèle nomme parfois les couleurs différemment)
\colorlet{popwhite}{white}
\colorlet{popblack}{popink}
\colorlet{popred}{poppink}
\colorlet{popyellow}{popgold}
% Teintes claires (fonds de tableaux/encadrés) — le modèle nomme parfois
% les couleurs avec un suffixe L (ex. popblueL) sans les avoir définies.
\colorlet{poporangeL}{poporange!20}
\colorlet{popdarkL}{popdark!20}
\colorlet{poppurpleL}{poppurple!20}
\colorlet{popgreenL}{popgreen!20}
\colorlet{poppinkL}{poppink!20}
\colorlet{popblueL}{popblue!20}
\colorlet{popgoldL}{popgold!20}
\colorlet{popredL}{popred!20}
\colorlet{popgrayL}{popgray!20}
% Teintes claires pour fonds de tableaux / zones
\colorlet{popblueL}{popblue!10!white}
\colorlet{poporangeL}{poporange!14!white}
\colorlet{popgreenL}{popgreen!12!white}
\colorlet{poppurpleL}{poppurple!12!white}
\colorlet{poppinkL}{poppink!12!white}
\colorlet{popgoldL}{popgold!14!white}

\pagecolor{popcream}

% ============================================================
% Typographie
% ============================================================
\defaultfontfeatures{Ligatures=TeX}
\setmainfont{PlusJakartaSans-Medium.ttf}[Path=../../../fonts/,BoldFont=PlusJakartaSans-Bold.ttf]
% Maths en sans-serif (Fira Math) avec chiffres et lettres droites de Plus
% Jakarta, pour que les formules se fondent dans le texte.
\usepackage[math-style=ISO,bold-style=ISO]{unicode-math}
\setmathfont{FiraMath-Regular.otf}[Path=../../../fonts/]
\setmathfont{PlusJakartaSans-Medium.ttf}[Path=../../../fonts/,range={up/num,up/{Latin,latin},\mathup}]
\usepackage{icomma} % virgule décimale sans espace en mode math
% Caractères mathématiques Unicode absents des polices de texte (Plus Jakarta,
% Archivo Black) : rendus par leur équivalent mathématique, en texte comme en
% formule — sinon ils s'affichent en carré vide (ex. « Arithmétique dans ℤ »).
\usepackage{newunicodechar}
\newunicodechar{ℝ}{\ensuremath{\mathbb{R}}}
\newunicodechar{ℤ}{\ensuremath{\mathbb{Z}}}
\newunicodechar{ℂ}{\ensuremath{\mathbb{C}}}
\newunicodechar{ℕ}{\ensuremath{\mathbb{N}}}
\newunicodechar{ℚ}{\ensuremath{\mathbb{Q}}}
\newunicodechar{𝔻}{\ensuremath{\mathbb{D}}}
\newunicodechar{α}{\ensuremath{\alpha}}
\newunicodechar{β}{\ensuremath{\beta}}
\newunicodechar{γ}{\ensuremath{\gamma}}
\newunicodechar{δ}{\ensuremath{\delta}}
\newunicodechar{ε}{\ensuremath{\varepsilon}}
\newunicodechar{θ}{\ensuremath{\theta}}
\newunicodechar{λ}{\ensuremath{\lambda}}
\newunicodechar{μ}{\ensuremath{\mu}}
\newunicodechar{σ}{\ensuremath{\sigma}}
\newunicodechar{τ}{\ensuremath{\tau}}
\newunicodechar{φ}{\ensuremath{\varphi}}
\newunicodechar{ψ}{\ensuremath{\psi}}
\newunicodechar{ω}{\ensuremath{\omega}}
\newunicodechar{Σ}{\ensuremath{\Sigma}}
% majuscules produites par \MakeUppercase dans les titres (α-aminés -> Α)
\newunicodechar{Α}{\ensuremath{\alpha}}
\newunicodechar{Β}{\ensuremath{\beta}}
\newunicodechar{Γ}{\ensuremath{\Gamma}}
\newunicodechar{Δ}{\ensuremath{\Delta}}
\newunicodechar{Ε}{\ensuremath{\varepsilon}}
\newunicodechar{Θ}{\ensuremath{\Theta}}
\newunicodechar{Λ}{\ensuremath{\Lambda}}
\newunicodechar{Μ}{\ensuremath{\mu}}
\newunicodechar{Τ}{\ensuremath{\tau}}
\newunicodechar{Φ}{\ensuremath{\Phi}}
\newunicodechar{Ψ}{\ensuremath{\Psi}}
\newunicodechar{Ω}{\ensuremath{\Omega}}
\newunicodechar{↦}{\ensuremath{\mapsto}}
\newunicodechar{∈}{\ensuremath{\in}}
\newunicodechar{∉}{\ensuremath{\notin}}
\newunicodechar{∧}{\ensuremath{\wedge}}
\newunicodechar{⇔}{\ensuremath{\Leftrightarrow}}
\newunicodechar{⟺}{\ensuremath{\Longleftrightarrow}}
\newunicodechar{⇒}{\ensuremath{\Rightarrow}}
\newunicodechar{⟹}{\ensuremath{\Longrightarrow}}
\newunicodechar{∘}{\ensuremath{\circ}}
\newunicodechar{∪}{\ensuremath{\cup}}
\newunicodechar{∩}{\ensuremath{\cap}}
\newunicodechar{≡}{\ensuremath{\equiv}}
\newunicodechar{⊂}{\ensuremath{\subset}}
\newunicodechar{⊥}{\ensuremath{\perp}}
\newunicodechar{∃}{\ensuremath{\exists}}
\newunicodechar{∀}{\ensuremath{\forall}}
\newunicodechar{∛}{\ensuremath{\sqrt[3]{\,}}}
\newunicodechar{∜}{\ensuremath{\sqrt[4]{\,}}}
\newunicodechar{⃗}{\ensuremath{{}^{\to}}}
\newunicodechar{ⁿ}{\ifmmode^{n}\else\textsuperscript{\ensuremath{n}}\fi}
\newunicodechar{ˣ}{\ifmmode^{x}\else\textsuperscript{\ensuremath{x}}\fi}
\newunicodechar{ᵇ}{\ifmmode^{b}\else\textsuperscript{\ensuremath{b}}\fi}
\newunicodechar{ʳ}{\ifmmode^{r}\else\textsuperscript{\ensuremath{r}}\fi}
\newunicodechar{ᵘ}{\ifmmode^{u}\else\textsuperscript{\ensuremath{u}}\fi}
\newunicodechar{ʸ}{\ifmmode^{y}\else\textsuperscript{\ensuremath{y}}\fi}
\newunicodechar{ᵃ}{\ifmmode^{a}\else\textsuperscript{\ensuremath{a}}\fi}
\newunicodechar{ᵉ}{\ifmmode^{e}\else\textsuperscript{\ensuremath{e}}\fi}
\newunicodechar{ᵏ}{\ifmmode^{k}\else\textsuperscript{\ensuremath{k}}\fi}
\newunicodechar{ᵖ}{\ifmmode^{p}\else\textsuperscript{\ensuremath{p}}\fi}
\newunicodechar{ᴺ}{\ifmmode^{N}\else\textsuperscript{\ensuremath{N}}\fi}
\newunicodechar{ᶜ}{\ifmmode^{c}\else\textsuperscript{\ensuremath{c}}\fi}
\newunicodechar{ʰ}{\ifmmode^{h}\else\textsuperscript{\ensuremath{h}}\fi}
\newunicodechar{ᵠ}{\ifmmode^{q}\else\textsuperscript{\ensuremath{q}}\fi}
\newunicodechar{ₙ}{\ifmmode_{n}\else\textsubscript{\ensuremath{n}}\fi}
\newunicodechar{ₐ}{\ifmmode_{a}\else\textsubscript{\ensuremath{a}}\fi}
\newunicodechar{ᵢ}{\ifmmode_{i}\else\textsubscript{\ensuremath{i}}\fi}
\newunicodechar{ᵣ}{\ifmmode_{r}\else\textsubscript{\ensuremath{r}}\fi}
\newunicodechar{ₖ}{\ifmmode_{k}\else\textsubscript{\ensuremath{k}}\fi}
\newunicodechar{ᵦ}{\ifmmode_{\beta}\else\textsubscript{\ensuremath{\beta}}\fi}
\newunicodechar{ₓ}{\ifmmode_{x}\else\textsubscript{\ensuremath{x}}\fi}
\newfontfamily\popdisplay{ArchivoBlack-Regular.ttf}[Path=../../../fonts/]
\newfontfamily\poplogo{SpaceGrotesk-Bold.ttf}[Path=../../../fonts/]
\newfontfamily\popsemi{PlusJakartaSans-SemiBold.ttf}[Path=../../../fonts/]
\newfontfamily\popxbold{PlusJakartaSans-ExtraBold.ttf}[Path=../../../fonts/]
\newfontfamily\popxboldls{PlusJakartaSans-ExtraBold.ttf}[Path=../../../fonts/,LetterSpace=3.0]

\setlist[enumerate]{leftmargin=1.7em,itemsep=5pt,topsep=4pt}
\setlist[itemize]{leftmargin=1.7em,itemsep=4pt,topsep=4pt,label={\color{poporange}\textbullet}}
\setlength{\parindent}{0pt}
\setlength{\parskip}{5pt}
\renewcommand{\arraystretch}{1.25}

% pgfplots : style Pop par défaut
\pgfplotsset{
  every axis/.append style={axis line style={popink,line width=1pt},tick style={popink},
    grid style={popline,line width=0.5pt},label style={font=\small},tick label style={font=\footnotesize}},
  popcurve/.style={poporange,line width=1.8pt,no markers,smooth},
}
% Même style disponible dans un \draw TikZ simple (hors pgfplots)
\tikzset{popcurve/.style={poporange,line width=1.8pt}}
\tikzset{popgrid/.style={popline,very thin}}
\colorlet{popcurve}{poporange} % le nom du style est parfois utilisé comme couleur

% ============================================================
% Pill / squiggle
% ============================================================
\newcommand{\poppill}[3]{%
  \tikz[baseline=-0.55ex]\node[draw=popink,line width=1.2pt,rounded corners=2.6pt,%
    fill=#2,inner xsep=6.5pt,inner ysep=3pt]{%
    \popxboldls\fontsize{7}{7}\selectfont\color{#3}\MakeUppercase{#1}};%
}
\newcommand{\popsquiggle}[1]{%
  \tikz[baseline=-0.4ex]{
    \draw[#1,line width=2.6pt,line cap=round]
      plot [smooth] coordinates {(0,0.09) (0.34,0.01) (0.68,0.21) (1.02,0.01) (1.36,0.10)};
  }%
}
% Mot-clé surligné (marqueur orange sous le texte)
\newcommand{\cle}[1]{{\popxbold\color{popdark}#1}}

% ============================================================
% En-tête / pied
% ============================================================
\pagestyle{fancy}
\fancyhf{}
\renewcommand{\headrulewidth}{0pt}
\renewcommand{\footrulewidth}{0pt}
\lhead{\raisebox{-0.15ex}{\poplogo\fontsize{13}{13}\selectfont\color{popink}OnBuch\color{poporange}+}}
\rhead{\poppill{\DOCMATIERE\ \textbullet\ \DOCNIVEAU\ \textbullet\ Leçon \DOCLECON}{white}{popink}}
\lfoot{\popxbold\fontsize{7.6}{7.6}\selectfont\color{popmuted}\MakeUppercase{Cours signé OnBuch+}\ \textbullet\ {\popsemi\DOCTITRE}}
\rfoot{\tikz[baseline=-0.6ex]\node[rounded corners=8.5pt,fill=popink,minimum height=17pt,minimum width=17pt,inner xsep=5pt,inner sep=0]{\popxbold\fontsize{8}{8}\selectfont\color{popcream}\thepage\,/\,\pageref*{LastPage}};}

% ============================================================
% Titres
% ============================================================
\newcommand{\chaptitre}[2]{%
  \begin{tcolorbox}[enhanced,colback=poporange,colframe=popink,boxrule=2.6pt,arc=11pt,
      drop shadow=popink,left=15pt,right=15pt,top=12pt,bottom=11pt,before skip=3pt,after skip=18pt]
    \poppill{#2}{popink}{popcream}\par\vspace{9pt}
    {\raggedright\hyphenpenalty=10000\exhyphenpenalty=10000\popdisplay\fontsize{22}{24}\selectfont\color{popcream}\MakeUppercase{#1}\par}
  \end{tcolorbox}
}
% Section numérotée : pastille orange avec le numéro + titre Archivo
\newcounter{popsec}
\newcommand{\coursec}[1]{%
  \stepcounter{popsec}\setcounter{popsub}{0}%
  \par\vspace{16pt}\noindent
  \begin{minipage}{\linewidth}
  \tikz[baseline=-0.7ex]\node[draw=popink,line width=1.6pt,fill=poporange,rounded corners=4pt,minimum size=22pt,inner sep=0]{\popdisplay\fontsize{12}{12}\selectfont\color{popcream}\thepopsec};%
  \hspace{8pt}{\popdisplay\fontsize{15}{17}\selectfont\color{popink}\MakeUppercase{#1}}\par
  \vspace{4pt}\noindent{\color{poporange}\rule{36pt}{3.6pt}}
  \end{minipage}\par\nopagebreak\vspace{8pt}
}
\newcounter{popsub}
\newcommand{\courssub}[1]{%
  \stepcounter{popsub}%
  \par\vspace{9pt}\noindent{\popxbold\fontsize{11.5}{13}\selectfont\color{popink}{\color{poporange}\thepopsec.\thepopsub}\hspace{6pt}#1}\par\nopagebreak\vspace{4pt}
}

% ============================================================
% Boîtes pédagogiques — même recette : fond blanc, bordure épaisse colorée,
% ombre dure encre, badge plein. Titre optionnel [#1].
% ============================================================
\tcbset{popbase/.style={breakable,enhanced,colback=white,boxrule=2.3pt,arc=9pt,
  shadow={3pt}{-3pt}{0pt}{popink},left=14pt,right=14pt,top=11pt,bottom=11pt,before skip=11pt,after skip=15pt,
  fontupper=\popsemi\fontsize{10.6}{14.5}\selectfont\color{popink},
  fontlower=\popsemi\fontsize{10.6}{14.5}\selectfont\color{popink}}}
% Coche dessinée (le glyphe ✓ n'existe pas dans les polices du document)
\newcommand{\popcheck}[1]{\tikz[baseline=-0.5ex]\draw[#1,line width=1.6pt,line cap=round,line join=round](-0.13,0.0)--(-0.04,-0.09)--(0.13,0.1);}
\providecommand{\coche}{\popcheck{popgreen}}
\providecommand{\resultat}[1]{\par\smallskip\noindent\fcolorbox{popgreen}{popgreenL}{\parbox{\dimexpr\linewidth-2\fboxsep-2\fboxrule\relax}{\textbf{Résultat.} #1}}\par\smallskip}
\newcommand{\popboxhead}[3]{\poppill{#1}{#2}{white}\ifx\relax#3\relax\else\hspace{6pt}{\popxbold\fontsize{10.6}{12}\selectfont\color{popink}#3}\fi\par\vspace{7pt}}

\newtcolorbox{aretenir}[1][]{popbase,colframe=popgreen,before upper={\popboxhead{À retenir}{popgreen}{#1}}}
\newtcolorbox{exemplebox}[1][]{popbase,colframe=poporange,before upper={\popboxhead{Exemple}{poporange}{#1}}}
\newtcolorbox{definition}[1][]{popbase,colframe=popblue,before upper={\popboxhead{Définition}{popblue}{#1}}}
\newtcolorbox{propriete}[1][]{popbase,colframe=poppurple,before upper={\popboxhead{Propriété}{poppurple}{#1}}}
\newtcolorbox{methode}[1][]{popbase,colframe=popgold,before upper={\popboxhead{Méthode}{popgold}{#1}}}
\newtcolorbox{attention}[1][]{popbase,colframe=poppink,before upper={\popboxhead{Attention}{poppink}{#1}}}
\newtcolorbox{experience}[1][]{popbase,colframe=popblue,colback=popblueL,before upper={\popboxhead{Expérience}{popblue}{#1}}}
% « Le savais-tu ? » : carte violette pleine (contexte, culture, Cameroun)
\newtcolorbox{savaistu}[1][]{popbase,colback=poppurple,colframe=popink,
  fontupper=\popsemi\fontsize{10.4}{14.2}\selectfont\color{white},
  before upper={\poppill{Le savais-tu\,?}{white}{poppurple}\ifx\relax#1\relax\else\hspace{6pt}{\popxbold\color{white}#1}\fi\par\vspace{7pt}}}
% Exercice résolu : énoncé en haut, \tcblower puis la solution sur fond crème
\newcounter{popexo}\newcounter{popexr}
\newtcolorbox{exoresolu}[1][]{popbase,colframe=poporange,colbacklower=popcream,
  segmentation style={popink,line width=1.2pt,dashed},
  before upper={\stepcounter{popexr}\popboxhead{Exercice résolu \thepopexr}{poporange}{#1}},
  before lower={\poppill{Solution}{popink}{popcream}\par\vspace{6pt}}}
% Exercice (non résolu en place) — la correction est dans la partie Corrigés
\newtcolorbox{exercice}[1][]{popbase,colframe=popink,
  before upper={\stepcounter{popexo}\popboxhead{Exercice \thepopexo}{popink}{#1}}}
% Corrigé
\newtcolorbox{corrige}[1][]{popbase,colframe=popgreen,colback=popgreenL,
  before upper={\popboxhead{Corrigé}{popgreen}{#1}}}
% Objectifs / prérequis / bilan
\newtcolorbox{objectifs}{popbase,colframe=popink,colback=poporangeL,
  before upper={\popboxhead{Objectifs de la leçon}{popink}{}}}
\newtcolorbox{prerequis}{popbase,colframe=popink,colback=popblueL,
  before upper={\popboxhead{Prérequis}{popblue}{}}}
\newtcolorbox{bilan}{popbase,colframe=popink,colback=white,boxrule=2.8pt,
  before upper={\popboxhead{Fiche bilan}{poporange}{L'essentiel à connaître pour l'examen}}}
% Figure encadrée avec légende
\newtcolorbox{popfigure}[1][]{enhanced,breakable=false,colback=white,colframe=popline,boxrule=1.4pt,arc=8pt,
  left=10pt,right=10pt,top=10pt,bottom=8pt,before skip=10pt,after skip=12pt,halign=center,
  lower separated=false,halign lower=center,
  fontlower=\popsemi\fontsize{9}{11.5}\selectfont\color{popmuted},
  #1}
\newcommand{\legende}[1]{\par\vspace{6pt}{\popsemi\fontsize{9}{11.5}\selectfont\color{popmuted}#1}}

% ============================================================
% Signature OnBuch+
% ============================================================
\newcommand{\poplogoinline}[1]{{\poplogo\fontsize{#1}{#1}\selectfont\color{popink}OnBuch\color{poporange}+}}
\newcommand{\signatureonbuch}{%
  \begin{tcolorbox}[enhanced,colback=popink,colframe=popink,boxrule=2.2pt,arc=10pt,
      drop shadow=poporange,left=14pt,right=14pt,top=10pt,bottom=10pt,before skip=0pt,after skip=0pt]
    \tikz[baseline=-0.7ex]\node[circle,fill=popgreen,draw=popcream,line width=1.3pt,minimum size=22pt,inner sep=0]{\popcheck{white}};%
    \hspace{9pt}\begin{minipage}[c]{0.62\linewidth}
      {\popxbold\fontsize{12}{13}\selectfont\color{popcream}Signé\ \textbullet\ {\poplogo OnBuch\color{poporange}+}}\par\vspace{2pt}
      {\raggedright\popsemi\fontsize{8}{10.5}\selectfont\color{popline}\MakeUppercase{Équipe pédagogique OnBuch+}\\\MakeUppercase{Conforme au programme officiel MINESEC}\par}
    \end{minipage}\hfill
    {\poplogo\fontsize{22}{22}\selectfont\color{popcream}OnBuch\color{poporange}+}
  \end{tcolorbox}%
}

% ============================================================
% Page de garde
% #1 titre · #2 matière · #3 niveau · #4 module · #5 n° leçon · #6 durée ·
% #7 description / objectifs (texte ou itemize)
% ============================================================
\newcommand{\pagegarde}[7]{%
  \thispagestyle{empty}
  \newgeometry{a4paper,margin=1.7cm,top=1.6cm,bottom=1.4cm}%
  \begin{tikzpicture}[remember picture,overlay]
    \fill[poporange!18] ([xshift=-3.2cm,yshift=-3.6cm]current page.north east) circle (4.6cm);
    \fill[poppurple!14] ([xshift=2.4cm,yshift=7.6cm]current page.south west) circle (3.8cm);
  \end{tikzpicture}%
  {\poplogo\fontsize{30}{30}\selectfont\color{popink}OnBuch\color{poporange}+}\hfill
  \raisebox{0.6ex}{\poppill{Cours signé \textbullet\ Premium}{popink}{popcream}}\par
  \vspace{1pt}\popsquiggle{poppurple}\par
  \vspace{8pt}
  \begin{tcolorbox}[enhanced,colback=poporange,colframe=popink,boxrule=2.8pt,arc=13pt,
      drop shadow=popink,left=18pt,right=18pt,top=12pt,bottom=13pt,after skip=10pt]
    \poppill{#2 \textbullet\ #3}{popink}{popcream}\hspace{5pt}\poppill{#4}{white}{popink}\par\vspace{8pt}
    {\popdisplay\fontsize{14}{14}\selectfont\color{popink}LEÇON #5}\par\vspace{4pt}
    {\raggedright\hyphenpenalty=10000\exhyphenpenalty=10000\popdisplay\fontsize{25}{27}\selectfont\color{popcream}\MakeUppercase{#1}\par}
    \vspace{8pt}
    {\popxbold\fontsize{9.5}{9.5}\selectfont\color{popink}\MakeUppercase{Durée conseillée : #6}}
  \end{tcolorbox}
  \begin{tcolorbox}[enhanced,colback=white,colframe=popink,boxrule=2.2pt,arc=10pt,
      drop shadow=popink,left=16pt,right=16pt,top=10pt,bottom=10pt,after skip=8pt]
    \poppill{Dans cette leçon}{poppurple}{white}\par\vspace{5pt}
    \popsemi\fontsize{9.3}{12.6}\selectfont\color{popink}#7
  \end{tcolorbox}
  \noindent\begin{minipage}[t]{0.487\linewidth}
    \begin{tcolorbox}[enhanced,colback=popgreen,colframe=popink,boxrule=2.2pt,arc=10pt,
        drop shadow=popink,left=11pt,right=11pt,top=10pt,bottom=10pt,height=3.1cm,valign=center]
      \tcbox[colback=white,colframe=popink,boxrule=1.3pt,arc=5pt,boxsep=0pt,nobeforeafter,
        left=3pt,right=3pt,top=3pt,bottom=3pt,valign=center]{\qrcode[height=1.9cm]{https://play.google.com/store/apps/details?id=com.luvvix.onbuch&hl=fr}}%
      \hspace{8pt}\begin{minipage}[c]{0.5\linewidth}
        {\popdisplay\fontsize{11}{12.5}\selectfont\color{white}\MakeUppercase{Télécharge OnBuch}}\par\vspace{4pt}
        {\raggedright\popsemi\fontsize{8.4}{11}\selectfont\color{white}Gratuit \textbullet\ Google Play\\Tuteur IA, annales, concours, résultats\par}
      \end{minipage}
    \end{tcolorbox}
  \end{minipage}\hfill
  \begin{minipage}[t]{0.487\linewidth}
    \begin{tcolorbox}[enhanced,colback=poppurple,colframe=popink,boxrule=2.2pt,arc=10pt,
        drop shadow=popink,left=11pt,right=11pt,top=10pt,bottom=10pt,height=3.1cm,valign=center]
      \tikz[baseline=-0.7ex]\node[circle,fill=white,draw=popink,line width=1.3pt,minimum size=1.6cm,inner sep=0]{\popdisplay\fontsize{24}{24}\selectfont\color{poppurple}+};%
      \hspace{8pt}\begin{minipage}[c]{0.6\linewidth}
        {\popdisplay\fontsize{11}{12.5}\selectfont\color{white}\MakeUppercase{Passe à OnBuch+}}\par\vspace{4pt}
        {\raggedright\popsemi\fontsize{8.4}{11}\selectfont\color{white}Tous les cours signés, Tuteur IA illimité, fiches PDF — onglet OnBuch+ de l'app\par}
      \end{minipage}
    \end{tcolorbox}
  \end{minipage}
  \vspace{8pt}
  \popcontenu
  \vfill
  \signatureonbuch
  \restoregeometry
  \newpage
}

% Rangée « Ce cours contient » (page de garde)
\newcommand{\poptile}[3]{% #1 couleur #2 gros texte #3 légende
  \begin{tcolorbox}[enhanced,colback=#1,colframe=popink,boxrule=2pt,arc=9pt,shadow={2.5pt}{-2.5pt}{0pt}{popink},
      left=6pt,right=6pt,top=9pt,bottom=9pt,halign=center,valign=center,height=1.75cm,width=\linewidth,nobeforeafter]
    {\popdisplay\fontsize{12.5}{13}\selectfont\color{white}\MakeUppercase{#2}}\par\vspace{3pt}
    {\centering\popsemi\fontsize{7.8}{9.5}\selectfont\color{white}#3\par}
  \end{tcolorbox}}
\newcommand{\popcontenu}{%
  \par\noindent{\popxboldls\fontsize{7.5}{7.5}\selectfont\color{popmuted}CE COURS SIGNÉ CONTIENT}\par\vspace{7pt}
  \noindent\begin{minipage}[t]{0.235\linewidth}\poptile{popblue}{Cours}{complet, illustré, pas à pas}\end{minipage}\hfill
  \begin{minipage}[t]{0.235\linewidth}\poptile{poporange}{Méthodes}{et exercices résolus}\end{minipage}\hfill
  \begin{minipage}[t]{0.235\linewidth}\poptile{poppink}{Exercices}{type examen + corrigés}\end{minipage}\hfill
  \begin{minipage}[t]{0.235\linewidth}\poptile{popgreen}{Fiche bilan}{l'essentiel à retenir}\end{minipage}\par}

% Sommaire
\newtcolorbox{sommaire}{enhanced,colback=white,colframe=popink,boxrule=2.2pt,arc=10pt,
  drop shadow=popink,left=18pt,right=18pt,top=13pt,bottom=13pt,before skip=6pt,after skip=20pt,
  before upper={\poppill{Sommaire}{poppurple}{white}\par\vspace{9pt}\popsemi\fontsize{10.6}{14.5}\selectfont\color{popink}}}

% Signature de fin de cours — poussée tout en bas de la dernière page
\newcommand{\findecours}{%
  \par\vfill\nopagebreak
  \begin{center}\popsquiggle{poporange}\end{center}\vspace{4pt}
  \signatureonbuch
}
\AtBeginDocument{\def\llbracket{\mathopen{[\mkern-3.5mu[}}\def\rrbracket{\mathclose{]\mkern-3.5mu]}}} % intervalles d'entiers (glyphe ⟦ absent de la police)
% Glyphes absents de Fira Math : redéfinis à partir de glyphes disponibles
\AtBeginDocument{%
  \def\setminus{\mathbin{\backslash}}\let\smallsetminus\setminus
  \def\vdots{\mathord{\vbox{\baselineskip3.2pt\lineskiplimit0pt\kern4pt\hbox{.}\hbox{.}\hbox{.}}}}%
  \def\ddots{\mathinner{\mkern1mu\raise6.5pt\hbox{.}\mkern2mu\raise3.6pt\hbox{.}\mkern2mu\raise0.7pt\hbox{.}\mkern1mu}}%
  \def\triangle{\mathord{\scalebox{0.9}{$\symup{\Delta}$}}}%
  \def\nabla{\mathord{\raisebox{\depth}{\scalebox{1}[-1]{$\symup{\Delta}$}}}}%
  \def\checkmark{\popcheck{popdark}}%
  \def\star{\mathbin{\ast}}\def\bigstar{\mathord{\ast}}%
  \def\diamond{\mathbin{\scalebox{0.6}{$\Diamond$}}}%
  \def\bigcup{\mathop{\mathchoice{\raisebox{-0.3em}{\scalebox{1.7}{$\cup$}}}{\scalebox{1.25}{$\cup$}}{\cup}{\cup}}\displaylimits}%
  \def\bigcap{\mathop{\mathchoice{\raisebox{-0.3em}{\scalebox{1.7}{$\cap$}}}{\scalebox{1.25}{$\cap$}}{\cap}{\cap}}\displaylimits}%
  \def\lesseqqgtr{\mathrel{\lessgtr}}\def\bigoplus{\mathop{\oplus}\displaylimits}%
}
\providecommand{\ssi}{si et seulement si} % abréviation fréquente
\AtBeginDocument{\providecommand{\wideparen}{\overparen}} % arc de cercle

\begin{document}

\pagegarde{Production d’écrits utilitaires : rédiger un rapport (rédaction de l’introduction, du corps et de la conclusion)}{Français}{1re Tech}{Français – classe de Première technique}{N11}{23 séances (fiche de progression harmonisée)}{Cette leçon outille l'élève de Première technique pour rédiger un rapport complet : préparation (structures externe et interne), rédaction de l'introduction, du corps et de la conclusion, puis amélioration du texte. Elle s'appuie sur l'exploitation de textes fonctionnels et iconiques et mobilise les tonalités satirique et tragique comme ressources langagières.
\vspace{4pt}
\begin{multicols}{2}\small
\begin{itemize}
\item À la fin de cette leçon, je sais définir le rapport et distinguer ses types (rapport de stage, d'activité, d'incident, de visite)
\item À la fin de cette leçon, je sais identifier la structure externe d'un rapport (en-tête, références, objet, date, signature)
\item À la fin de cette leçon, je sais organiser la structure interne d'un rapport (introduction, corps, conclusion)
\item À la fin de cette leçon, je sais analyser une situation concrète et élaborer le plan d'un rapport
\item À la fin de cette leçon, je sais rédiger l'introduction et la conclusion d'un rapport dans un style objectif
\item À la fin de cette leçon, je sais rédiger le corps d'un rapport en présentant les faits, leur analyse et des recommandations
\end{itemize}\end{multicols}}

\begin{sommaire}
\begin{multicols}{2}
\begin{enumerate}
\item Le rapport : définition, types et fonctions
\item La structure externe du rapport
\item La structure interne : introduction, corps, conclusion
\item Exploiter les documents : textes fonctionnels et iconiques
\item Les tonalités satirique et tragique
\item De la situation au plan : préparer et rédiger le rapport
\item Compte rendu, remédiation et évaluation
\item Activité d'intégration
\item Méthodes et exercices résolus
\item Exercices
\item Corrigés des exercices
\item Fiche bilan
\end{enumerate}
\end{multicols}
\end{sommaire}

\begin{objectifs}
\begin{itemize}
\item À la fin de cette leçon, je sais définir le rapport et distinguer ses types (rapport de stage, d'activité, d'incident, de visite)
\item À la fin de cette leçon, je sais identifier la structure externe d'un rapport (en-tête, références, objet, date, signature)
\item À la fin de cette leçon, je sais organiser la structure interne d'un rapport (introduction, corps, conclusion)
\item À la fin de cette leçon, je sais analyser une situation concrète et élaborer le plan d'un rapport
\item À la fin de cette leçon, je sais rédiger l'introduction et la conclusion d'un rapport dans un style objectif
\item À la fin de cette leçon, je sais rédiger le corps d'un rapport en présentant les faits, leur analyse et des recommandations
\item À la fin de cette leçon, je sais exploiter des textes fonctionnels et iconiques (tableaux, photographies, organigrammes) pour alimenter un rapport
\item À la fin de cette leçon, je sais reconnaître et employer les tonalités satirique et tragique dans un texte
\item À la fin de cette leçon, je sais confronter ma production à celle d'autrui, la corriger et l'améliorer
\end{itemize}
\end{objectifs}

\begin{prerequis}
\begin{itemize}
\item Le compte rendu : définition et rédaction (classe de Seconde)
\item Les types de textes et les textes fonctionnels (début d'année de Première)
\item Le plan et les connecteurs logiques (classes de collège)
\item Les temps du récit et la concordance des temps
\item La lettre administrative : en-tête, objet, formules de politesse
\end{itemize}
\end{prerequis}

\newpage

\coursec{Le rapport : définition, types et fonctions}

Le préfet du Mfoundi a demandé au proviseur d'un lycée technique de Yaoundé un rapport sur l'état des ateliers après les pluies diluviennes d'octobre. Le délégué de classe, chargé de collecter les informations, a rédigé un long texte désordonné, mélangeant plaintes, rumeurs et faits : le document a été rejeté. Qu'est-ce qui distingue un simple récit d'un rapport professionnel ? Comment organiser les faits observés, les analyser et formuler des recommandations ? Cette leçon nous apprend à rédiger, de l'introduction à la conclusion, un écrit utilitaire conforme aux exigences de la vie administrative et professionnelle camerounaise.

\courssub{Qu'est-ce qu'un rapport ?}

\textbf{Intuition.} Imaginez que vous êtes délégué de classe. La porte de la salle informatique est cassée, trois ordinateurs sont hors d'usage. Vous pouvez en parler à vos camarades : c'est une conversation. Vous pouvez écrire dans votre cahier ce qui s'est passé : c'est un récit personnel. Mais si le proviseur vous demande un document écrit qui décrit les dégâts, en explique les causes et propose des solutions, vous devez produire un \cle{rapport}. Ce qui change, c'est le \textbf{destinataire} (une autorité), le \textbf{but} (informer pour permettre une décision) et la \textbf{forme} (un texte ordonné, objectif, daté et signé).

\begin{definition}[Le rapport]
Le \cle{rapport} est un écrit fonctionnel (ou utilitaire) qui présente de manière \textbf{objective et ordonnée} des faits constatés au cours d'une mission, d'une activité, d'une visite ou d'un incident, leur \textbf{analyse} et éventuellement des \textbf{recommandations}, à destination d'une autorité (chef d'établissement, chef d'entreprise, préfet, ministre, encadreur).
\end{definition}

Un rapport repose donc sur trois piliers qu'il faut retenir ensemble :

\begin{itemize}
\item \textbf{Des faits observés} : ce que l'on a vu, mesuré, entendu, vérifié. Rien qui ne soit vérifiable.
\item \textbf{Une analyse} : les causes, les conséquences, l'importance des faits.
\item \textbf{Des recommandations} (souvent, pas toujours) : ce que l'auteur propose à l'autorité de décider.
\end{itemize}

\begin{aretenir}[Les trois questions du rapport]
Tout rapport répond, dans l'ordre, à trois questions : \textbf{Que s'est-il passé ?} (les faits) — \textbf{Pourquoi et avec quelles conséquences ?} (l'analyse) — \textbf{Que faut-il faire ?} (les recommandations). Un texte qui ne répond qu'à la première question n'est pas encore un rapport : c'est un simple récit.
\end{aretenir}

\courssub{Rapport, compte rendu, procès-verbal, lettre administrative : ne pas confondre}

Dans la vie administrative camerounaise, plusieurs écrits utilitaires se ressemblent. Il est essentiel de les distinguer, car l'examinateur du Probatoire technique comme le chef de service attendent le bon document pour la bonne situation.

\begin{itemize}
\item Le \cle{compte rendu} relate fidèlement et chronologiquement ce qui s'est passé ou ce qui a été dit (une réunion, une visite, un cours). Il \textbf{relate} sans analyser ni proposer.
\item Le \cle{rapport} relate aussi, mais il \textbf{analyse} les faits et \textbf{propose} des solutions. Il engage davantage la réflexion de son auteur.
\item Le \cle{procès-verbal} (PV) est un document à valeur juridique qui constate officiellement des faits ou des décisions (PV de réunion, PV d'examen, PV de constat d'huissier). Il est signé par les personnes concernées et peut servir de preuve.
\item La \cle{lettre administrative} sert à \textbf{demander}, \textbf{informer} ou \textbf{répondre} (demande de stage, réclamation, convocation). Elle ne rend pas compte d'une mission.
\end{itemize}

\begin{methode}[Distinguer un rapport d'un compte rendu]
Pour savoir si un texte est un rapport ou un compte rendu, appliquez ces étapes :
\begin{enumerate}
\item Repérez la \textbf{fonction dominante} du texte : s'il se contente de relater fidèlement des faits ou des paroles, c'est un compte rendu ; s'il analyse et propose, c'est un rapport.
\item Cherchez la présence d'une \textbf{analyse} (causes, conséquences, interprétation des faits) : absente du compte rendu, présente dans le rapport.
\item Cherchez des \textbf{recommandations ou suggestions} adressées à l'autorité : c'est la marque la plus sûre du rapport.
\item Vérifiez le \textbf{degré d'engagement de l'auteur} : le rapporteur prend position de manière argumentée ; le rédacteur du compte rendu reste un simple témoin.
\end{enumerate}
Retenez la formule : \textbf{le compte rendu relate fidèlement, le rapport analyse et propose.}
\end{methode}

\begin{popfigure}
\begin{center}
\begin{tabularx}{\linewidth}{|l|X|X|X|}
\hline
\rowcolor{popblueL} \textbf{Critère} & \textbf{Rapport} & \textbf{Compte rendu} & \textbf{Procès-verbal} \\
\hline
\textbf{Destinataire} & Une autorité (chef, préfet, ministre) & Un supérieur ou un groupe & Autorité administrative ou judiciaire \\
\hline
\textbf{Objet} & Rendre compte d'une mission, analyser, proposer & Relater fidèlement un événement ou des paroles & Constater officiellement des faits ou décisions \\
\hline
\textbf{Ton} & Objectif, analytique, argumenté & Objectif, neutre, descriptif & Solennel, précis, juridique \\
\hline
\textbf{Structure} & Introduction, corps (faits, analyse, recommandations), conclusion & Ordre chronologique simple & Mentions légales, constats, signatures \\
\hline
\textbf{Valeur} & Document de travail et d'aide à la décision & Document d'information & Document faisant preuve \\
\hline
\end{tabularx}
\end{center}
\legende{Tableau comparatif des trois écrits utilitaires proches du rapport : critères de distinction.}
\end{popfigure}

\courssub{Les types de rapports}

Selon la mission accomplie, on distingue plusieurs types de rapports. Tous partagent la même structure générale, mais diffèrent par leur objet.

\begin{itemize}
\item \textbf{Le rapport de stage.} Rédigé par un élève ou un étudiant à l'issue d'une période d'immersion en entreprise, il présente la structure d'accueil, les tâches effectuées, les acquis de la formation et des suggestions. C'est le rapport que vous rédigerez vous-même au cours de votre formation technique.
\item \textbf{Le rapport d'activité.} Il fait le bilan d'une activité menée sur une période donnée : rapport d'activité d'un club, d'une coopérative, d'un atelier de maintenance.
\item \textbf{Le rapport d'incident.} Il décrit un événement anormal (accident, panne, dégradation, bagarre), en établit les circonstances et les responsabilités éventuelles, et propose des mesures pour éviter qu'il ne se reproduise.
\item \textbf{Le rapport de visite ou d'inspection.} Rédigé après la visite d'un lieu (chantier, école, marché, usine), il décrit l'état des lieux, signale les insuffisances et formule des recommandations.
\item \textbf{Le rapport d'enquête.} Il présente les résultats d'une investigation menée sur une question précise (enquête sur l'absentéisme, sur la qualité d'un service) à partir de données recueillies.
\end{itemize}

\begin{popfigure}
\begin{center}
\begin{tikzpicture}[
  grow=right,
  level 1/.style={sibling distance=2.1cm, level distance=3.2cm},
  level 2/.style={sibling distance=1.1cm, level distance=5.2cm, font=\small},
  racine/.style={rectangle, rounded corners, draw=popdark, fill=popblue, text=white, font=\bfseries, align=center, inner sep=5pt},
  type/.style={rectangle, rounded corners, draw=popdark, fill=popblueL, align=center, inner sep=4pt, font=\bfseries\small},
  ex/.style={rectangle, rounded corners, draw=popline, fill=popcream, align=left, inner sep=3pt}
]
\node[racine, align=center]{Types de\\rapports}
  child { node[type, align=center]{Rapport\\d'enquête} child { node[ex, align=center]{Enquête sur l'absentéisme\\au lycée technique de Bafoussam} } }
  child { node[type, align=center]{Rapport de visite\\ou d'inspection} child { node[ex, align=center]{Inspection des ateliers\\après les pluies (Mfoundi)} } }
  child { node[type, align=center]{Rapport\\d'incident} child { node[ex, align=center]{Panne du transformateur\\d'un atelier de Douala} } }
  child { node[type, align=center]{Rapport\\d'activité} child { node[ex, align=center]{Bilan annuel du club\\informatique du lycée} } }
  child { node[type, align=center]{Rapport\\de stage} child { node[ex, align=center]{Stage d'un élève de Première F3\\dans une entreprise de Yaoundé} } };
\end{tikzpicture}
\end{center}
\legende{Classification des principaux types de rapports, avec un exemple camerounais pour chacun.}
\end{popfigure}

\begin{savaistu}[Le rapport de stage, un document décisif]
Au Cameroun, tout stagiaire en entreprise doit remettre un rapport de stage à la fin de sa période d'immersion. Ce document est souvent \textbf{exigé pour la validation du diplôme technique} (CAP, Probatoire et Baccalauréat techniques, BTS). Il est déposé auprès de l'établissement de formation et parfois soutenu devant un jury. Savoir rédiger un rapport n'est donc pas un exercice scolaire abstrait : c'est une compétence qui conditionne votre réussite professionnelle.
\end{savaistu}

\courssub{Les qualités d'un bon rapport}

Un rapport efficace possède cinq qualités inséparables. Reprenons la situation du préfet du Mfoundi : c'est précisément parce que le texte du délégué les violait toutes qu'il a été rejeté.

\begin{propriete}[Les cinq qualités du rapport]
\begin{enumerate}
\item \textbf{L'objectivité} : rapporter uniquement des faits vérifiables, sans rumeur, sans opinion personnelle, sans jugement de valeur.
\item \textbf{L'exactitude} : donner des informations précises (dates, lieux, nombres, noms), que l'on a soi-même vérifiées.
\item \textbf{La clarté} : employer un vocabulaire simple et précis, des phrases bien construites, un ton neutre et administratif.
\item \textbf{La concision} : aller à l'essentiel, sans détails inutiles ni répétitions.
\item \textbf{L'ordre logique} : présenter les faits dans un plan cohérent (chronologique ou thématique), avec des paragraphes et des titres.
\end{enumerate}
\end{propriete}

\begin{attention}[Erreur fréquente à éviter]
L'erreur la plus grave dans un rapport est d'y introduire des \textbf{rumeurs}, des \textbf{opinions personnelles} ou des \textbf{jugements de valeur}. Écrire « on dit que le gardien vend le carburant du groupe électrogène » ou « le menuisier de l'atelier est incompétent » détruit la crédibilité du document et peut exposer son auteur à des poursuites. On n'écrit que ce que l'on a \textbf{constaté} : « le réservoir du groupe électrogène était vide le 12 octobre, alors que la fiche de suivi indique un remplissage le 10 octobre ». Le fait remplace l'accusation.
\end{attention}

\courssub{Destinataires et contextes d'usage au Cameroun}

Le rapport circule dans toute la vie administrative et professionnelle camerounaise. Selon le contexte, le destinataire change :

\begin{itemize}
\item le \textbf{proviseur ou le principal} reçoit les rapports des délégués, des surveillants généraux, des chefs d'atelier ;
\item l'\textbf{encadreur de stage} et le chef d'établissement reçoivent le rapport de stage de l'élève ;
\item le \textbf{chef d'entreprise} reçoit les rapports d'activité et d'incident de ses employés ;
\item le \textbf{préfet} reçoit les rapports des chefs d'établissement et des services déconcentrés de son département ;
\item le \textbf{ministre} reçoit les rapports d'inspection et d'enquête de ses services.
\end{itemize}

Dans tous les cas, la logique est la même : une autorité a \textbf{commandé} ou \textbf{attend} un document pour \textbf{prendre une décision}. Le rapport est donc un outil de décision, pas un simple exercice d'écriture.

\courssub{Exploitation du texte fonctionnel n\textsuperscript{o} 2 : un modèle de rapport}

Observons un modèle pour y repérer les caractéristiques du genre.

\begin{exemplebox}[Rapport du délégué de classe au proviseur]
\textbf{Rapport sur la dégradation de la salle informatique du lycée de Nkongsamba}\par
\emph{De :} MBAPPE Arnaud, délégué de la classe de Première F3.\par
\emph{À :} Monsieur le Proviseur du lycée de Nkongsamba.\par
\emph{Date :} le 15 octobre 2025.\par
\emph{Objet :} dégradation de la salle informatique.\par
\textbf{Introduction.} À la suite de votre demande en date du 13 octobre 2025, j'ai procédé, le 14 octobre, à la visite de la salle informatique afin de constater l'état des lieux.\par
\textbf{Faits constatés.} La porte principale est descellée ; trois ordinateurs sur vingt sont hors d'usage ; le climatiseur ne fonctionne plus ; l'humidité a endommagé deux prises électriques.\par
\textbf{Analyse.} Ces dégradations résultent des pluies abondantes d'octobre et de l'absence de gouttière au-dessus de la porte. Elles empêchent la tenue normale des travaux pratiques.\par
\textbf{Recommandations.} Je suggère la réparation urgente de la porte et des prises, le dépannage du climatiseur et l'installation d'une gouttière.\par
\textbf{Conclusion.} Je reste à votre disposition pour tout complément d'information.\par
\emph{Signature.}
\end{exemplebox}

Ce modèle illustre les marques du genre rapport : des \textbf{mentions d'en-tête} (auteur, destinataire, date, objet), une \textbf{introduction} qui rappelle la mission, un \textbf{corps} qui sépare nettement les faits de l'analyse et des recommandations, une \textbf{conclusion} brève et courtoise, et une \textbf{signature}. Le ton est neutre, les faits sont chiffrés (« trois ordinateurs sur vingt »), aucune rumeur n'apparaît.

\begin{exoresolu}[Identifier les qualités du rapport dans un extrait]
\textbf{Énoncé.} Relisez l'extrait suivant : « La porte principale est descellée ; trois ordinateurs sur vingt sont hors d'usage. » Indiquez deux qualités du rapport que cette phrase illustre et justifiez votre réponse.
\tcblower
\textbf{Solution.} Cette phrase illustre d'abord \textbf{l'objectivité} : l'auteur se limite à des faits qu'il a lui-même constatés lors de sa visite, sans accuser personne ni rapporter de rumeur. Elle illustre ensuite \textbf{l'exactitude} : les informations sont précises et chiffrées (« trois ordinateurs sur vingt », « la porte principale »), ce qui permet au proviseur de vérifier et de chiffrer les réparations. On peut ajouter la \textbf{concision} : une seule phrase suffit à rendre compte de deux constats essentiels.
\end{exoresolu}

\begin{exercice}[À vous de juger]
Le délégué de la classe de Première F4 écrit dans son rapport : « Tout le monde sait que c'est le gardien qui a laissé la porte ouverte, et d'ailleurs il ne fait jamais rien. » Relevez les deux fautes commises contre les qualités du rapport, puis réécrivez la phrase de manière conforme.
\end{exercice}

\begin{corrige}[Exercice : À vous de juger]
Première faute : la \textbf{rumeur} (« tout le monde sait que... ») remplace le fait vérifié. Deuxième faute : le \textbf{jugement de valeur} (« il ne fait jamais rien ») est une opinion personnelle, injurieuse de surcroît. Réécriture conforme : « Lors de ma visite du 14 octobre à 16 heures, la porte de la salle informatique était ouverte et sans surveillance ; le registre de présence du gardien ne mentionne aucune ronde après 14 heures. » Le rapporteur constate, il n'accuse pas : c'est à l'autorité de tirer les conclusions.
\end{corrige}

\begin{aretenir}[L'essentiel de la section]
Le rapport est un écrit utilitaire qui présente objectivement des faits constatés, leur analyse et des recommandations à une autorité. Il se distingue du compte rendu (qui relate sans analyser) et du procès-verbal (qui a valeur de preuve). Il existe plusieurs types de rapports (stage, activité, incident, visite, enquête), tous soumis aux mêmes exigences : objectivité, exactitude, clarté, concision et ordre logique. Maîtriser ces notions est le premier pas avant d'étudier, dans la prochaine section, la structure externe du rapport.
\end{aretenir}

\coursec{La structure externe du rapport}

Dans la section précédente, nous avons défini le rapport comme un texte fonctionnel qui rend compte, de manière objective et fidèle, d'une mission, d'une visite, d'une enquête ou d'une situation observée. Nous avons distingué ses types (rapport de stage, rapport de visite, rapport d'incident, rapport d'activité) et ses fonctions (informer, témoigner, permettre une décision). Mais avant même de rédiger le contenu, il faut savoir une chose essentielle : un rapport n'est pas une simple feuille couverte de phrases. C'est un document administratif, et comme tout document administratif, il obéit à une présentation codifiée. C'est cette présentation que nous appelons la \cle{structure externe}.

\courssub{Qu'est-ce que la structure externe ?}

Imaginez deux documents posés sur le bureau du proviseur. Le premier est une feuille blanche remplie de texte, sans nom, sans date, sans signature. Le second porte en haut à gauche l'en-tête du lycée, un numéro de référence, un objet clairement annoncé, la date, et se termine par une signature accompagnée d'une fonction. Lequel des deux sera enregistré, classé et pris au sérieux ? Le second, évidemment. Le premier est un texte perdu : on ne sait ni qui l'a écrit, ni quand, ni pour qui, ni à quoi il sert.

\begin{definition}[Structure externe]
La \cle{structure externe} d'un rapport est l'ensemble des éléments qui encadrent le texte proprement dit et qui permettent de l'\textbf{identifier}, de le \textbf{classer} et de le \textbf{faire circuler} : l'en-tête, les références, le destinataire, l'objet, la date et le lieu de rédaction, les pièces jointes, la signature et la fonction du rédacteur. Elle s'oppose à la \emph{structure interne} (introduction, corps, conclusion), qui concerne le contenu lui-même.
\end{definition}

\begin{aretenir}[Pourquoi la structure externe est-elle obligatoire ?]
Un rapport sans structure externe complète est un document \textbf{sans valeur administrative} : il ne peut pas être enregistré au bureau d'ordre, ni classé dans un dossier, ni cité comme preuve. La forme n'est pas un détail de politesse : elle garantit l'authenticité et la traçabilité du document.
\end{aretenir}

\courssub{L'en-tête : identifier la structure ou l'auteur}

L'\cle{en-tête} occupe la partie supérieure de la première page. Il indique \textbf{qui écrit} : soit une institution (lycée, entreprise, ministère), soit une personne à titre individuel.

Dans un rapport administratif camerounais, l'en-tête d'un établissement public suit une présentation traditionnelle en deux colonnes :

\begin{itemize}
\item à \textbf{gauche}, les mentions républicaines : \emph{République du Cameroun}, suivie de la devise \emph{Paix -- Travail -- Patrie}, puis le ministère de tutelle (par exemple \emph{Ministère des Enseignements Secondaires}) ;
\item à \textbf{droite}, parfois, la version anglaise : \emph{Republic of Cameroon -- Peace -- Work -- Fatherland} ;
\item en dessous, le nom de la structure émettrice : \emph{Lycée Technique de\ldots}, éventuellement le service ou l'atelier concerné.
\end{itemize}

Pour un rapport de stage ou un rapport individuel d'élève, l'en-tête peut être plus simple : nom et prénom du rédacteur, classe ou fonction, nom de l'établissement. L'essentiel est que le lecteur sache immédiatement \textbf{d'où vient le document}.

\begin{exemplebox}[Exemple d'en-tête]
\begin{center}
\begin{tabular}{p{0.45\linewidth} p{0.45\linewidth}}
RÉPUBLIQUE DU CAMEROUN & REPUBLIC OF CAMEROON\\
\emph{Paix -- Travail -- Patrie} & \emph{Peace -- Work -- Fatherland}\\
Ministère des Enseignements Secondaires & \\
Lycée Technique de Bafoussam & \\
\end{tabular}
\end{center}
\end{exemplebox}

\courssub{Les références et le destinataire}

Sous l'en-tête, à gauche, figurent les \cle{références} : le \textbf{numéro du rapport}, composé selon un code propre à la structure. Ce numéro permet d'enregistrer le document au bureau d'ordre et de le retrouver dans les archives. Il est souvent suivi de la mention du destinataire.

Le \cle{destinataire} est la personne ou l'autorité à qui le rapport est adressé. On l'introduit par la formule \emph{« À l'attention de\ldots »} suivie du titre et du nom : \emph{À l'attention de Monsieur le Proviseur}, \emph{À l'attention de Madame la Directrice des Ressources Humaines}. On ne s'adresse jamais à une fonction sans la nommer correctement, et on n'oublie jamais la civilité.

\begin{exemplebox}[Exemple chiffré de références]
\textbf{Rapport N°12/LT/2026} du 15 mars 2026, à l'attention de Monsieur le Proviseur.

Dans ce numéro : \emph{12} est le numéro d'ordre du rapport dans l'année, \emph{LT} désigne le Lycée Technique, \emph{2026} l'année d'émission. Ce codage permet de classer et de retrouver le document sans ambiguïté.
\end{exemplebox}

\begin{savaistu}[Le bureau d'ordre]
Dans l'administration camerounaise, tout document entrant ou sortant transite par le \textbf{bureau d'ordre}, qui lui attribue un numéro d'enregistrement et une date d'arrivée. Un document \textbf{sans référence ne peut pas être enregistré} : il est comme invisible. C'est pourquoi les chefs de service refusent de traiter toute correspondance dépourvue de numéro. Rédiger une référence, ce n'est pas faire de la décoration : c'est donner une existence officielle à son texte.
\end{savaistu}

\courssub{L'objet du rapport : annoncer le contenu en une ligne}

L'\cle{objet} est la formule brève qui annonce le thème du rapport. Il se place après les références, souvent précédé du mot \emph{« Objet : »} et parfois souligné ou mis en gras. C'est la première chose que le destinataire lit pour savoir de quoi il retourne.

\begin{methode}[Rédiger un objet correct]
\begin{enumerate}
\item Repérer le \textbf{thème exact} du rapport (ce dont on rend compte : un stage, une visite, un incident, un état des lieux).
\item Formuler une \textbf{phrase nominale courte}, c'est-à-dire un groupe nominal \textbf{sans verbe conjugué}.
\item Commencer par un mot-clé : \emph{Rapport sur\ldots}, \emph{Compte rendu de\ldots}, \emph{État de\ldots}
\item Vérifier que l'objet tient en \textbf{une seule ligne} et ne contient ni détail superflu ni jugement personnel.
\end{enumerate}
\end{methode}

\begin{exemplebox}[Objets bien rédigés]
\begin{itemize}
\item \emph{Objet : Rapport sur l'état des ateliers de menuiserie.}
\item \emph{Objet : Compte rendu de la visite de l'entreprise SOCATER du 10 mars 2026.}
\item \emph{Objet : Rapport de stage effectué à la SODECOTON du 2 au 27 février 2026.}
\end{itemize}
\end{exemplebox}

\begin{attention}[Pièges de l'objet]
\begin{itemize}
\item Ne pas écrire une phrase complète avec verbe conjugué : \emph{« J'ai visité les ateliers et j'ai constaté que\ldots »} n'est pas un objet, c'est déjà le corps du rapport.
\item Ne pas être vague : \emph{« Objet : Rapport »} n'informe sur rien.
\item Ne pas y glisser son opinion : \emph{« Rapport sur le désastreux état des ateliers »} est subjectif ; le jugement, s'il est fondé, appartient au corps du texte.
\end{itemize}
\end{attention}

\courssub{La date et le lieu de rédaction}

Tout rapport indique \textbf{où} et \textbf{quand} il a été rédigé. La formule traditionnelle se place à droite, avant le corps du texte ou en fin de document : \emph{« Fait à Bafoussam, le 15 mars 2026 »}. La date est capitale : c'est elle qui permet de situer les faits rapportés dans le temps, de vérifier des délais, et de trancher en cas de contestation. Un rapport non daté ne prouve rien.

\courssub{Les pièces jointes}

Un rapport peut être accompagné de \cle{pièces jointes} : photographies des lieux visités, tableaux de mesures ou de chiffres, schémas, copies de documents, annexes. Elles servent de \textbf{preuves matérielles} aux affirmations du corps du rapport. On les annonce clairement, en bas de page ou après la signature, par la mention \emph{« Pièces jointes : »} suivie de leur liste et de leur nombre : \emph{Pièces jointes : 3 photographies, 1 tableau récapitulatif}. Chaque pièce doit être numérotée et légendée pour être citée dans le texte (\emph{cf. photo n°1, annexe 2}).

\courssub{La signature et la fonction du rédacteur}

Le rapport se termine, en bas à droite, par la \cle{signature} du rédacteur, précédée ou suivie de son nom en toutes lettres et de sa \textbf{fonction} : \emph{L'élève chargé de mission, MBAPPE Larissa, classe de Première F3} ; \emph{Le chef d'atelier, NGUEFACK Paul}. La signature engage la responsabilité de celui qui écrit : elle certifie que le contenu est exact et assumé. Un rapport anonyme n'a aucune force.

\courssub{Schéma d'ensemble : la page d'un rapport}

\begin{popfigure}
\begin{center}
\begin{tikzpicture}
% Page
\draw[fill=popcream, draw=popdark, thick] (0,0) rectangle (7.2,10.2);
% En-tête
\node[align=left, font=\footnotesize, anchor=north west] at (0.4,9.9) {RÉPUBLIQUE DU CAMEROUN\\ \emph{Paix -- Travail -- Patrie}\\ Ministère des Enseignements Secondaires\\ Lycée Technique de Bafoussam};
\node[align=left, font=\footnotesize, anchor=north east] at (6.8,9.9) {REPUBLIC OF CAMEROON\\ \emph{Peace -- Work -- Fatherland}};
\draw[popline] (0.4,8.3) -- (6.8,8.3);
% Références
\node[align=left, font=\footnotesize, anchor=north west] at (0.4,8.1) {\textbf{Rapport N°12/LT/2026}\\ À l'attention de Monsieur le Proviseur};
% Date
\node[align=left, font=\footnotesize, anchor=north west] at (4.3,7.3) {Bafoussam, le 15 mars 2026};
% Objet
\node[align=left, font=\footnotesize, anchor=north west] at (0.4,6.6) {\textbf{Objet :} Rapport sur l'état des ateliers de menuiserie.};
% Corps fictif
\foreach \y in {5.9,5.6,5.3,5.0,4.7,4.4,4.1,3.8,3.5,3.2}{\draw[popmuted, line width=1.5pt] (0.4,\y) -- (6.8,\y);}
\node[font=\footnotesize\itshape, text=popmuted] at (3.6,2.9) {(corps du rapport : introduction, développement, conclusion)};
% Pièces jointes
\node[align=left, font=\footnotesize, anchor=north west] at (0.4,2.4) {\textbf{Pièces jointes :} 3 photographies, 1 tableau.};
% Signature
\node[align=center, font=\footnotesize, anchor=north west] at (4.3,2.2) {L'élève chargé de mission,\\[0.5em] \emph{(signature)}\\ MBAPPE Larissa, Première F3};
% Flèches d'annotation
\draw[-{Stealth}, popblue, thick] (8.6,9.2) -- (7.3,9.2) node[pos=0, right, font=\footnotesize, text=popblue, align=left]{1. En-tête :\\ qui écrit};
\draw[-{Stealth}, poporange, thick] (8.6,7.7) -- (4.6,7.7) node[pos=0, right, font=\footnotesize, text=poporange, align=left]{2. Références et\\ destinataire};
\draw[-{Stealth}, popgreen, thick] (8.6,6.9) -- (7.3,6.9) node[pos=0, right, font=\footnotesize, text=popgreen, align=left]{3. Date et lieu};
\draw[-{Stealth}, poppurple, thick] (-1.4,6.3) -- (0.3,6.3) node[pos=0, left, font=\footnotesize, text=poppurple, align=right]{4. Objet :\\ de quoi il s'agit};
\draw[-{Stealth}, poppink, thick] (-1.4,1.9) -- (0.3,1.9) node[pos=0, left, font=\footnotesize, text=poppink, align=right]{5. Pièces\\ jointes};
\draw[-{Stealth}, popgold, thick] (8.6,1.5) -- (6.6,1.5) node[pos=0, right, font=\footnotesize, text=popgold, align=left]{6. Signature et\\ fonction};
\end{tikzpicture}
\end{center}
\legende{Schéma annoté d'une page de rapport : les six éléments de la structure externe. Seule la zone centrale (traits gris) constitue la structure interne, étudiée dans la section suivante.}
\end{popfigure}

\courssub{Exploitation du texte fonctionnel N°3 : un rapport administratif anonymisé}

Observons maintenant un document réel (anonymisé) et traquons-y chaque élément de la structure externe. C'est la démarche du lecteur professionnel : avant de lire le fond, on identifie la forme.

\begin{popfigure}
\begin{center}
\begin{tikzpicture}
\draw[fill=white, draw=popdark, thick] (0,0) rectangle (8,9.6);
\node[align=left, font=\scriptsize, anchor=north west] at (0.3,9.4) {RÉPUBLIQUE DU CAMEROUN\\ \emph{Paix -- Travail -- Patrie}\\ MINESEC -- Délégation Régionale de l'Ouest\\ Lycée Technique de \ldots};
\draw[popline] (0.3,7.9) -- (7.7,7.9);
\node[align=left, font=\scriptsize, anchor=north west] at (0.3,7.7) {\textbf{N°045/LT/DRO/2026}\hspace{2cm} À l'attention de Monsieur le Délégué Régional};
\node[align=left, font=\scriptsize, anchor=north west] at (4.6,6.9) {\ldots, le 22 avril 2026};
\node[align=left, font=\scriptsize, anchor=north west] at (0.3,6.3) {\textbf{Objet :} Rapport sur l'incident survenu à l'atelier d'électrotechnique.};
\foreach \y in {5.5,5.2,4.9,4.6,4.3,4.0,3.7,3.4}{\draw[popmuted, line width=1.5pt] (0.3,\y) -- (7.7,\y);}
\node[align=left, font=\scriptsize, anchor=north west] at (0.3,2.9) {\textbf{P.J. :} 2 photographies, 1 procès-verbal de constat.};
\node[align=center, font=\scriptsize, anchor=north west] at (4.6,2.5) {Le surveillant de permanence,\\[0.4em] \emph{(signature)}\\ T\ldots J\ldots};
% annotations
\draw[popblue, thick, rounded corners] (0.25,8.05) rectangle (5.6,9.5);
\node[font=\scriptsize, text=popblue, anchor=west] at (8.3,8.8) {En-tête complet};
\draw[-{Stealth}, popblue] (8.3,8.7) -- (5.7,8.7);
\draw[poporange, thick, rounded corners] (0.25,7.15) rectangle (7.75,7.85);
\node[font=\scriptsize, text=poporange, anchor=west] at (8.3,7.5) {Référence + destinataire};
\draw[-{Stealth}, poporange] (8.3,7.4) -- (7.85,7.4);
\draw[poppurple, thick, rounded corners] (0.25,5.95) rectangle (7.75,6.5);
\node[font=\scriptsize, text=poppurple, anchor=west] at (8.3,6.2) {Objet : phrase nominale};
\draw[-{Stealth}, poppurple] (8.3,6.1) -- (7.85,6.1);
\draw[popgold, thick, rounded corners] (4.55,1.6) rectangle (7.75,2.6);
\node[font=\scriptsize, text=popgold, anchor=west] at (8.3,2.1) {Signature + fonction};
\draw[-{Stealth}, popgold] (8.3,2.0) -- (7.85,2.0);
\end{tikzpicture}
\end{center}
\legende{Reproduction schématisée et anonymisée d'un rapport administratif camerounais (texte fonctionnel N°3) : repérage visuel des éléments externes.}
\end{popfigure}

\begin{exoresolu}[Identifier la structure externe d'un document]
\textbf{Énoncé.} Voici le début et la fin d'un document : \emph{« Lycée Technique de Garoua. N°08/LTG/2026. À l'attention de Madame la Censeure. Objet : Compte rendu de la sortie pédagogique au barrage de Lagdo. [\ldots] Pièces jointes : 5 photographies. Le professeur accompagnateur, signé : ALIM Georges. »} Relevez les éléments de la structure externe et signalez ce qui manque.
\tcblower
\textbf{Solution détaillée.}
\begin{itemize}
\item \textbf{En-tête} : présent mais incomplet (manquent les mentions républicaines et le ministère de tutelle).
\item \textbf{Référence} : présente (N°08/LTG/2026).
\item \textbf{Destinataire} : présent (Madame la Censeure).
\item \textbf{Objet} : présent, correctement rédigé en phrase nominale.
\item \textbf{Pièces jointes} : annoncées (5 photographies).
\item \textbf{Signature et fonction} : présentes (le professeur accompagnateur, ALIM Georges).
\item \textbf{Éléments manquants} : la \textbf{date} et le \textbf{lieu} de rédaction (\emph{« Fait à Garoua, le\ldots »}). Sans date, le rapport ne peut pas situer les faits ni être valablement archivé : sa valeur administrative est compromise.
\end{itemize}
\end{exoresolu}

\begin{attention}[Les trois oublis qui annulent un rapport]
L'erreur la plus fréquente des candidats comme des rédacteurs pressés est d'oublier la \textbf{date}, la \textbf{signature} ou le \textbf{destinataire}. Or un rapport sans date ne situe aucun fait, un rapport sans signature n'engage personne, et un rapport sans destinataire ne sait pas où il va. Dans les trois cas, le document est \textbf{sans valeur administrative} : il peut être refusé, et à l'examen, ces oublis coûtent des points de présentation. Avant de rendre toute copie ou tout rapport, vérifiez la liste complète : en-tête, référence, destinataire, objet, date et lieu, pièces jointes, signature et fonction.
\end{attention}

\begin{exercice}[À vous de compléter]
Reconstituez la structure externe complète d'un rapport à partir des informations suivantes : vous êtes élève en Première F4 au Lycée Technique de Ngaoundéré ; le 5 mai 2026, vous rendez compte à Monsieur le Proviseur de l'état du parc informatique ; votre rapport est le troisième de l'année ; vous joignez un tableau des machines hors service. Rédigez l'en-tête, la référence, le destinataire, l'objet, la date et la mention des pièces jointes, puis la formule de signature.
\end{exercice}

\begin{corrige}[Exercice : la structure externe complète]
Éléments attendus : en-tête avec \emph{République du Cameroun -- Paix -- Travail -- Patrie}, MINESEC, Lycée Technique de Ngaoundéré ; référence du type \emph{Rapport N°03/LTN/2026} ; destinataire \emph{À l'attention de Monsieur le Proviseur} ; objet en phrase nominale : \emph{Rapport sur l'état du parc informatique} ; date : \emph{Ngaoundéré, le 5 mai 2026} ; pièces jointes : \emph{1 tableau des machines hors service} ; signature : \emph{L'élève chargé de mission}, nom et prénom, \emph{classe de Première F4}.
\end{corrige}

La structure externe est donc le cadre qui donne au rapport son existence officielle. Mais un beau cadre ne suffit pas : il faut maintenant remplir ce cadre. C'est l'objet de la section suivante, consacrée à la structure interne du rapport : l'introduction, le corps et la conclusion.

\coursec{La structure interne : introduction, corps, conclusion}

Dans la section précédente, nous avons étudié la \cle{structure externe} du rapport : l'en-tête, les mentions d'identification (expéditeur, destinataire, objet, date, références), la pagination et la signature. Mais un rapport ne se réduit pas à sa « coquille » : ce qui fait sa valeur, c'est ce qu'il contient et la manière dont ce contenu est organisé. C'est précisément l'objet de cette section : la \cle{structure interne}, c'est-à-dire l'organisation du propos en trois parties obligatoires — l'\cle{introduction}, le \cle{corps} et la \cle{conclusion}.

\courssub{Qu'est-ce que la structure interne ?}

\begin{definition}[Structure interne du rapport]
La \cle{structure interne} d'un rapport est l'organisation du contenu du document en trois parties ordonnées : l'\textbf{introduction} (qui situe le contexte, rappelle la mission et annonce le plan), le \textbf{corps} (qui présente les faits constatés, leur analyse et les recommandations) et la \textbf{conclusion} (qui dresse un bilan synthétique et se termine par une formule de courtoisie administrative).
\end{definition}

L'intuition est simple : un rapport est une \emph{réponse écrite à une mission}. Celui qui le lit (un chef d'établissement, un directeur d'entreprise, un ministre) veut savoir trois choses, dans cet ordre : \emph{pourquoi} ce rapport existe, \emph{ce que} l'auteur a constaté et propose, et \emph{ce qu'il faut retenir} en fin de compte. Ces trois attentes correspondent exactement aux trois parties de la structure interne.

\begin{popfigure}
\begin{center}
\begin{tikzpicture}[
  bloc/.style={draw=popblue, thick, rounded corners=3pt, fill=popblueL, text width=0.88\linewidth, align=left, inner sep=6pt},
  fleche/.style={-{Stealth[length=3mm]}, very thick, poporange}
]
\node[bloc, align=center] (intro) at (0,0){\textcolor{popblue}{\textbf{1. INTRODUCTION}}\\ Contexte et rappel de la mission $\bullet$ période et lieu concernés $\bullet$ méthode d'enquête $\bullet$ annonce du plan};
\node[bloc, fill=poporangeL, draw=poporange, align=center] (corps) at (0,-2.6){\textcolor{poporange}{\textbf{2. CORPS DU RAPPORT}}\\ Faits constatés (ordre chronologique, thématique ou par importance) $\bullet$ analyse (causes, conséquences, comparaisons, chiffres) $\bullet$ recommandations};
\node[bloc, fill=popgreenL, draw=popgreen, align=center] (concl) at (0,-5.2){\textcolor{popgreen}{\textbf{3. CONCLUSION}}\\ Bilan synthétique $\bullet$ formule de courtoisie administrative};
\draw[fleche] (intro.south) -- (corps.north);
\draw[fleche] (corps.south) -- (concl.north);
\end{tikzpicture}
\end{center}
\legende{Schéma de la structure interne du rapport : trois blocs enchaînés, du contexte vers le bilan.}
\end{popfigure}

\courssub{L'introduction : situer la mission et annoncer le plan}

L'introduction d'un rapport n'est pas une introduction de dissertation : elle ne pose pas un problème philosophique, elle \emph{rappelle des faits administratifs}. Elle doit permettre à un lecteur qui découvrirait le document isolément de comprendre immédiatement de quoi il s'agit.

\begin{methode}[Rédiger l'introduction d'un rapport]
Pour rédiger une introduction complète, réponds dans l'ordre aux cinq questions suivantes :
\begin{enumerate}
  \item \textbf{Qui ?} — Qui a reçu la mission et qui l'a confiée ? (« À la suite de la lettre n\textsuperscript{o} 45/DP du 12 mars 2026, par laquelle Monsieur le Directeur m'a chargé de... »)
  \item \textbf{Quoi ?} — Quel est l'objet précis de la mission ? (visite d'inspection, enquête, stage, étude...)
  \item \textbf{Où ?} — Quel est le lieu concerné ? (entreprise, établissement, localité...)
  \item \textbf{Quand ?} — Quelle est la période couverte ? (dates de début et de fin)
  \item \textbf{Comment ?} — Quelle méthode a été utilisée ? (observation, entretiens, questionnaires, consultation de documents...)
\end{enumerate}
Termine ensuite par l'\textbf{annonce du plan} : « Ce rapport présente d'abord..., ensuite..., enfin... »
\end{methode}

\begin{exemplebox}[Introduction d'un rapport de stage à la SODECOTON de Garoua]
« Conformément à la convention de stage signée entre le Lycée technique de Garoua et la SODECOTON, j'ai effectué un stage d'observation et de pratique du \textbf{2 janvier au 27 février 2026} au \textbf{service maintenance mécanique} de l'usine de Garoua. Ce stage avait pour objectifs de découvrir le fonctionnement d'une unité industrielle d'égrenage du coton, d'appliquer les connaissances acquises en atelier et d'observer l'organisation du travail en équipe. Au cours de cette période, j'ai mené mon enquête par l'observation directe des postes de travail, par des entretiens avec le chef de service et par la consultation des registres de maintenance. Le présent rapport présente d'abord la structure d'accueil, ensuite les activités réalisées, enfin les difficultés rencontrées et les suggestions qu'elles inspirent. »
\end{exemplebox}

Observe ce modèle : en huit lignes, le lecteur connaît le cadre (la convention), le qui, le quoi, le où, le quand, le comment, et le plan. Rien de plus n'est nécessaire.

\courssub{Le corps du rapport : faits, analyse, recommandations}

Le corps est la partie la plus longue et la plus importante : il occupe généralement les deux tiers du document. Il comporte trois temps, qu'il ne faut jamais confondre.

\textbf{1. La présentation ordonnée des faits constatés.} Les faits sont des réalités observées, vérifiables, rapportées sans commentaire personnel. Ils doivent être présentés selon un ordre clair, choisi en fonction du sujet :
\begin{itemize}
  \item l'ordre \cle{chronologique} : on suit le déroulement dans le temps (idéal pour un rapport de stage ou de voyage d'étude : semaine par semaine) ;
  \item l'ordre \cle{thématique} : on regroupe les faits par thèmes (idéal pour un rapport d'enquête : état des bâtiments, équipements, effectifs, discipline) ;
  \item l'ordre \cle{d'importance} : on va du fait le plus grave ou le plus important au moins important (idéal pour un rapport d'inspection ou d'incident).
\end{itemize}

\textbf{2. L'analyse des faits.} Après les faits viennent leur interprétation objective : rechercher les \emph{causes} (pourquoi la panne est-elle survenue ?), mesurer les \emph{conséquences} (quel impact sur la production ?), établir des \emph{comparaisons} (avec l'année précédente, avec une autre filière) et s'appuyer sur des \emph{données chiffrées}. C'est ici que les tableaux et diagrammes prennent toute leur valeur : un chiffre bien présenté vaut un long paragraphe.

\begin{popfigure}
\begin{center}
\begin{tikzpicture}
\begin{axis}[
  ybar, bar width=18pt,
  width=0.85\linewidth, height=6cm,
  ymin=0, ymax=100,
  ylabel={Taux de réussite au Probatoire technique (\%)},
  symbolic x coords={F1,F2,F3,F4},
  xtick=data,
  xlabel={Filières techniques},
  nodes near coords,
  nodes near coords style={font=\small, text=popdark},
  every node near coord/.append style={/pgf/number format/.cd, fixed, fixed zerofill, precision=1},
  axis line style={popline},
  tick label style={text=popdark},
  label style={text=popdark}
]
\addplot[fill=popblue, draw=popblue] coordinates {(F1,72.5) (F2,64.0) (F3,58.5) (F4,81.0)};
\end{axis}
\end{tikzpicture}
\end{center}
\legende{Exemple de données chiffrées intégrées au corps d'un rapport scolaire : taux de réussite au Probatoire technique par filière (F1 : mécanique ; F2 : électrotechnique ; F3 : construction ; F4 : chimie). Données fictives à visée pédagogique. Dans le corps du rapport, on écrira par exemple : « La filière F4 obtient le meilleur taux (81 \%), tandis que la filière F3, avec 58,5 \%, accuse un retard de 22,5 points qui s'explique en partie par le manque de matériel de chantier. »}
\end{popfigure}

\textbf{3. Les recommandations ou suggestions.} Le rapport ne s'arrête pas au constat : il propose. Les recommandations sont des \emph{propositions concrètes, réalistes et adressées au destinataire} (c'est lui qui décidera). Elles se formulent souvent à l'infinitif ou avec « il conviendrait de », « nous suggérons de » : « Nous recommandons l'achat de dix postes de soudure avant la rentrée prochaine » ; « Il conviendrait de renforcer l'encadrement des élèves de F3 en travaux pratiques ».

\begin{attention}[Erreurs fréquentes dans le corps du rapport]
\begin{itemize}
  \item \textbf{Placer les recommandations dans l'introduction} : l'introduction annonce la mission, elle ne propose rien. Les suggestions se trouvent à la fin du corps, après les faits et leur analyse.
  \item \textbf{Mélanger faits et opinions} : « Les ateliers sont dans un état lamentable » est un jugement ; « Trois machines sur huit sont hors service depuis janvier » est un fait. Dans un rapport, on écrit le fait d'abord, et l'appréciation, si elle est nécessaire, doit rester mesurée et fondée sur le fait.
  \item \textbf{Accumuler sans ordre} : des faits jetés pêle-mêle fatiguent le lecteur. Choisis un ordre (chronologique, thématique ou par importance) et signale-le.
\end{itemize}
\end{attention}

\courssub{La conclusion : bilan et courtoisie}

La conclusion est la partie la plus courte, mais elle est décisive : c'est souvent elle que le destinataire lit en premier, juste après l'introduction.

\begin{methode}[Rédiger la conclusion d'un rapport]
\begin{enumerate}
  \item \textbf{Synthétiser sans répéter} : en deux ou trois phrases, dégage l'essentiel du constat et des propositions, sans recopier les phrases du corps. Utilise des formulations comme « Au terme de cette mission, il apparaît que... », « En somme... ».
  \item \textbf{Formuler le vœu ou l'attente} : « Nous espérons que ces suggestions retiendront votre attention. »
  \item \textbf{Terminer par une formule de courtoisie administrative} : « Nous restons à votre disposition pour tout complément d'information. », « Veuillez agréer, Monsieur le Directeur, l'expression de notre haute considération. »
\end{enumerate}
\end{methode}

\begin{exemplebox}[Conclusion du rapport de stage à la SODECOTON]
« Au terme de ce stage, il apparaît que la SODECOTON offre un cadre de formation riche et exigeant, où les acquis de l'atelier scolaire trouvent une application directe. Les difficultés signalées — notamment l'insuffisance de matériel de protection individuelle — appellent des mesures simples que nous avons proposées. Nous espérons que ce rapport rendra compte fidèlement de l'intérêt de cette expérience et restons à votre disposition pour tout complément d'information. Veuillez agréer, Monsieur le Proviseur, l'expression de notre respectueuse considération. »
\end{exemplebox}

\courssub{Les connecteurs logiques propres au rapport}

La clarté d'un rapport repose sur des \cle{connecteurs logiques} qui guident le lecteur d'une étape à l'autre. Retiens les principaux :

\begin{center}
\begin{tabularx}{\linewidth}{|l|X|}
\hline
\rowcolor{popblueL} \textbf{Fonction} & \textbf{Connecteurs utiles} \\
\hline
Ordonner les faits & d'abord, ensuite, puis, enfin \\
\hline
Ajouter une information & en outre, par ailleurs, de plus, également \\
\hline
Expliquer (cause) & parce que, en effet, du fait de, grâce à \\
\hline
Montrer une conséquence & c'est pourquoi, par conséquent, ainsi, de sorte que \\
\hline
Comparer & contrairement à, tandis que, par rapport à \\
\hline
Conclure & en conclusion, en somme, au terme de, pour conclure \\
\hline
\end{tabularx}
\end{center}

Ces marqueurs d'organisation sont précieux aussi en \emph{lecture} : dans les textes fonctionnels n\textsuperscript{o} 3 et n\textsuperscript{o} 4 étudiés en réception, c'est en repérant « d'abord / ensuite / enfin », « en outre / par ailleurs » et « en conclusion » que l'on découpe rapidement le document en ses trois parties et que l'on reconstitue son plan.

\courssub{Exploitation des textes fonctionnels n\textsuperscript{o} 3 et n\textsuperscript{o} 4}

\begin{exoresolu}[Repérer les trois parties d'un rapport]
\textbf{Énoncé.} Voici un extrait de rapport (texte fonctionnel n\textsuperscript{o} 3, adapté) : « (a) À la demande de Monsieur le Proviseur, exprimée lors du conseil d'établissement du 5 octobre 2025, nous avons procédé, du 10 au 24 octobre 2025, à l'inspection des ateliers du Lycée technique de Maroua. (b) En outre, le poste de soudure n\textsuperscript{o} 2 est hors service depuis la rentrée. (c) En conclusion, l'état général des ateliers est satisfaisant, à l'exception du parc de soudure, et nous restons à votre disposition pour tout complément d'information. » Indique à quelle partie (introduction, corps, conclusion) appartient chaque segment et justifie.
\tcblower
\textbf{Solution détaillée.}
\begin{itemize}
  \item Le segment (a) appartient à l'\textbf{introduction} : il rappelle qui a confié la mission (Monsieur le Proviseur), l'objet (inspection des ateliers), le lieu (Lycée technique de Maroua) et la période (du 10 au 24 octobre 2025). Ce sont les marques typiques de l'introduction.
  \item Le segment (b) appartient au \textbf{corps} : il présente un fait constaté, vérifiable (un poste hors service), introduit par le connecteur d'addition « en outre », propre à l'enchaînement des constatations.
  \item Le segment (c) appartient à la \textbf{conclusion} : le marqueur « en conclusion » l'annonce ; il contient un bilan synthétique (« l'état général est satisfaisant ») et une formule de courtoisie (« nous restons à votre disposition... »).
\end{itemize}
\end{exoresolu}

\begin{exercice}[À toi de jouer]
À partir du texte fonctionnel n\textsuperscript{o} 4 distribué en classe : 1) souligne en bleu les éléments de l'introduction (mission, lieu, période, méthode) ; 2) souligne en vert deux faits constatés et une donnée chiffrée dans le corps ; 3) encadre les recommandations ; 4) surligne la formule de courtoisie finale ; 5) relève cinq connecteurs logiques et précise leur fonction.
\end{exercice}

\begin{corrige}[Exercice — indications de correction]
L'introduction se reconnaît aux mentions de la mission et des dates ; les faits sont des énoncés vérifiables au présent ou au passé composé, sans jugement ; les recommandations contiennent des verbes à l'infinitif ou des tournures comme « il conviendrait de » ; la formule de courtoisie occupe la dernière phrase ; les connecteurs attendus sont du type d'abord, ensuite, en outre, par ailleurs, enfin, en conclusion.
\end{corrige}

\begin{savaistu}[Un bon rapport se lit « en diagonale »]
Dans l'administration et l'entreprise, un chef de service reçoit parfois des dizaines de rapports par semaine. Il les lit « en diagonale » : l'introduction pour connaître la mission, la conclusion pour connaître l'essentiel du bilan et des propositions. Si ces deux parties sont claires, il plongera dans le corps du document ; si elles sont confuses, le rapport sera classé sans suite. C'est pourquoi introduction et conclusion doivent pouvoir se lire seules, comme un résumé fidèle de l'ensemble.
\end{savaistu}

\begin{aretenir}[L'essentiel]
La structure interne du rapport comporte trois parties : l'\textbf{introduction} (qui, quoi, où, quand, comment + annonce du plan), le \textbf{corps} (faits ordonnés, analyse avec chiffres, recommandations) et la \textbf{conclusion} (bilan synthétique + formule de courtoisie). Les connecteurs logiques (d'abord, ensuite, en outre, par ailleurs, enfin, en conclusion) assurent la lisibilité de l'ensemble. On ne mélange jamais faits et opinions, et l'on ne propose rien dans l'introduction.
\end{aretenir}

\coursec{Exploiter les documents : textes fonctionnels et iconiques}

Dans la section précédente, nous avons appris à construire la structure interne du rapport : une introduction qui présente le contexte et la mission, un corps organisé en parties, une conclusion qui propose des recommandations. Mais un rapport ne se rédige pas à partir de rien : il s'appuie sur des \cle{documents} que le rapporteur a recueillis sur le terrain ou reçus de son administration. Un procès-verbal de réunion, une lettre de mission, une photographie, un tableau de chiffres, un plan de quartier : autant de pièces qu'il faut savoir lire, analyser et transformer en phrases rédigées. C'est précisément l'objet de cette section : apprendre à \emph{exploiter} ces documents pour nourrir le corps du rapport.

\courssub{Lecture méthodique d'un texte fonctionnel}

\begin{definition}[Texte fonctionnel]
Un \cle{texte fonctionnel} est un texte écrit qui remplit une fonction pratique dans la vie professionnelle ou administrative : informer, ordonner, attester, demander, rapporter. Il s'oppose au texte littéraire, qui vise le plaisir esthétique. Exemples : la lettre administrative, la note de service, le procès-verbal, le compte rendu, le rapport lui-même.
\end{definition}

Lire un texte fonctionnel, ce n'est pas le parcourir comme un roman. C'est l'interroger méthodiquement pour en extraire les informations utiles au rapport. Quatre questions guident cette lecture :

\begin{methode}[Lire un texte fonctionnel en quatre questions]
\begin{enumerate}
\item \textbf{Qui parle ?} Identifier l'\cle{émetteur} : nom, fonction, service (exemple : « le Sous-préfet de Douala III »).
\item \textbf{À qui ?} Identifier le \cle{destinataire} : personne ou service visé (exemple : « Monsieur le Délégué du Gouvernement »).
\item \textbf{De quoi s'agit-il ?} Repérer l'\cle{objet}, souvent indiqué dans une ligne « Objet : » ou dans la première phrase.
\item \textbf{Comment est-ce organisé ?} Dégager l'\cle{organisation} : rappel des faits, décision, demande, délai, formule finale.
\end{enumerate}
\end{methode}

\begin{exemplebox}[Lecture d'une lettre de mission]
« \emph{Monsieur ABEGA, technicien communal, est chargé de visiter le marché Sandaga, à la suite des inondations du 12 octobre, et de présenter un rapport au plus tard le 30 octobre.} »
\begin{itemize}
\item Émetteur : l'autorité municipale (le Maire).
\item Destinataire : M. Abega, technicien communal.
\item Objet : visite du marché Sandaga après inondation ; production d'un rapport.
\item Organisation : désignation de l'agent, cause de la mission, résultat attendu, délai.
\end{itemize}
Ces quatre éléments fourniront directement la matière de l'introduction du rapport (contexte, mission, délai).
\end{exemplebox}

\courssub{Lecture d'un texte iconique}

\begin{definition}[Texte iconique]
Un \cle{texte iconique} est un document à dominante visuelle qui transmet une information : \cle{photographie}, \cle{tableau statistique}, \cle{schéma}, \cle{organigramme}, \cle{plan}, graphique, carte. Contrairement au texte rédigé, il ne « dit » rien avec des phrases : c'est au lecteur de transformer l'image ou les chiffres en énoncés.
\end{definition}

Un rapport technique s'appuie très souvent sur de tels documents : photo d'un bâtiment fissuré, tableau de fréquentation d'un service, organigramme d'une entreprise, plan d'un quartier sinistré. Leur lecture obéit à une méthode en quatre étapes.

\begin{methode}[Lire un document iconique en quatre étapes]
\begin{enumerate}
\item \textbf{Identifier} : de quel type de document s'agit-il ? Qui l'a produit, quand, à quel propos ? (lire le titre, la source, la date).
\item \textbf{Décrire} : que voit-on objectivement ? (éléments visibles, colonnes du tableau, unités, légendes) — sans encore interpréter.
\item \textbf{Analyser} : que signifient ces éléments ? (comparaisons, écarts, tendances, relations de cause à effet).
\item \textbf{Exploiter} : quelle information retenir pour le rapport, et sous quelle formulation rédigée ?
\end{enumerate}
\end{methode}

\begin{popfigure}
\begin{center}
\begin{tikzpicture}[node distance=0.55cm, every node/.style={font=\small}]
\node[draw=popblue, fill=popblueL, rounded corners, minimum width=4.6cm, minimum height=1cm, align=center] (e1) {\textbf{1. Identifier}\\ type, source, date, titre};
\node[draw=poporange, fill=poporangeL, rounded corners, minimum width=4.6cm, minimum height=1cm, align=center, below=of e1] (e2) {\textbf{2. Décrire}\\ éléments visibles, unités, légende};
\node[draw=poppurple, fill=poppurpleL, rounded corners, minimum width=4.6cm, minimum height=1cm, align=center, below=of e2] (e3) {\textbf{3. Analyser}\\ écarts, comparaisons, tendances};
\node[draw=popgreen, fill=popgreenL, rounded corners, minimum width=4.6cm, minimum height=1cm, align=center, below=of e3] (e4) {\textbf{4. Exploiter}\\ information retenue, phrase rédigée};
\draw[-{Stealth}, thick, popdark] (e1) -- (e2);
\draw[-{Stealth}, thick, popdark] (e2) -- (e3);
\draw[-{Stealth}, thick, popdark] (e3) -- (e4);
\end{tikzpicture}
\end{center}
\legende{Méthode de lecture d'un document iconique en quatre étapes : identifier, décrire, analyser, exploiter.}
\end{popfigure}

Appliquons cette méthode aux trois grands types de documents iconiques rencontrés dans les textes N°1, N°2 et N°4 de la séquence.

\textbf{La photographie.} On identifie d'abord le lieu, le moment et les personnes ou objets représentés. On décrit ensuite du général vers le particulier : vue d'ensemble, puis détails significatifs (eau stagnante, étals renversés, marchandises abîmées). On analyse enfin : que révèle la photo ? Un problème de drainage, une insalubrité, une baisse d'activité.

\textbf{Le tableau statistique.} On lit le titre, les intitulés de lignes et de colonnes, les unités. On repère les valeurs extrêmes (maximum, minimum), on calcule les écarts, on dégage une tendance (hausse, baisse, stabilité).

\textbf{L'organigramme et le plan.} L'organigramme montre la hiérarchie d'un service : on lit du haut (l'autorité) vers le bas (les exécutants), ce qui permet d'identifier les responsables à citer dans un rapport. Le plan montre la disposition spatiale : on l'oriente (nord), on repère les axes et les zones concernées par le problème étudié.

\courssub{Transformer une donnée iconique en phrase de rapport}

Le cœur du travail du rapporteur est là : convertir ce qu'il voit en ce qu'il écrit. Trois opérations sont à maîtriser.

\textbf{Décrire} : rendre compte fidèlement, avec un vocabulaire précis et neutre. On n'écrit pas « le marché est dans un état épouvantable », mais « les allées du marché sont recouvertes de 20 à 30 cm d'eau stagnante ».

\textbf{Chiffrer} : remplacer les impressions par des données mesurées. « Beaucoup de clients ont déserté » devient « la fréquentation quotidienne est tombée de 2 400 à 900 visiteurs ».

\textbf{Comparer} : mettre deux données en relation pour faire apparaître un écart, souvent exprimé en pourcentage. Le taux de variation se calcule ainsi :
\[ \text{Taux de variation} = \frac{\text{valeur d'arrivée} - \text{valeur de départ}}{\text{valeur de départ}} \times 100 \]

\begin{methode}[Exploiter un tableau chiffré]
\begin{enumerate}
\item Relever les \textbf{valeurs extrêmes} (la plus forte, la plus faible).
\item Calculer les \textbf{écarts} (différence absolue, puis taux de variation en \%).
\item Dégager une \textbf{tendance} (augmentation, diminution, stagnation).
\item Rédiger une \textbf{phrase de synthèse} qui énonce le chiffre, l'écart et l'interprétation.
\end{enumerate}
\end{methode}

\begin{exemplebox}[Exemple chiffré]
Fréquentation du marché Sandaga : 2 400 visiteurs par jour avant les pluies, 900 après. Calcul de l'écart :
\[ \frac{900 - 2400}{2400} \times 100 = \frac{-1500}{2400} \times 100 = -62{,}5\ \% \]
Phrase de synthèse pour le rapport : « La fréquentation du marché est passée de 2 400 à 900 visiteurs par jour, soit une baisse de 62,5 \%. »
\end{exemplebox}

\begin{popfigure}
\begin{center}
\begin{tikzpicture}
\begin{axis}[
  ybar, bar width=14pt,
  width=0.85\linewidth, height=6.2cm,
  ymin=0, ymax=2800,
  ylabel={Visiteurs par jour},
  symbolic x coords={Lundi, Mardi, Mercredi, Jeudi, Vendredi, Samedi},
  xtick=data,
  legend style={at={(0.5,-0.18)}, anchor=north, legend columns=2, font=\small},
  ymajorgrids=true, grid style={popline, dashed},
  every axis/.append style={font=\small}
]
\addplot[fill=popblue, draw=popblue] coordinates {(Lundi,2200) (Mardi,2350) (Mercredi,2400) (Jeudi,2300) (Vendredi,2500) (Samedi,2600)};
\addplot[fill=poporange, draw=poporange] coordinates {(Lundi,800) (Mardi,900) (Mercredi,950) (Jeudi,850) (Vendredi,1000) (Samedi,1100)};
\legend{Avant les pluies, Après les pluies}
\end{axis}
\end{tikzpicture}
\end{center}
\legende{Fréquentation hebdomadaire du marché Sandaga (Douala) avant et après les pluies (données fictives d'entraînement, en visiteurs par jour). La baisse est constante sur toute la semaine, autour de 60 \%.}
\end{popfigure}

\begin{attention}[Erreurs fréquentes]
Deux fautes reviennent sans cesse dans les copies. Premièrement, \textbf{recopier un tableau sans l'analyser} : le rapport n'est pas un album de documents, chaque donnée citée doit être commentée et mise en relation avec le problème étudié. Deuxièmement, \textbf{tirer des conclusions non justifiées par les données} : si le tableau montre une baisse de fréquentation, on peut l'affirmer ; mais on ne peut pas écrire « les commerçants sont ruinés » si aucune donnée ne porte sur leurs revenus. Toute affirmation du rapport doit pouvoir renvoyer à un document ou à une observation.
\end{attention}

\courssub{Intégrer les documents en annexe et y faire référence}

Les documents ne se collent pas n'importe où dans le rapport. La règle professionnelle est la suivante : le corps du texte \emph{cite} et \emph{commente} l'information ; le document complet est placé en \cle{annexe}, à la fin du rapport.

\begin{aretenir}[Règles d'utilisation des annexes]
\begin{itemize}
\item Numéroter chaque annexe : annexe 1, annexe 2, etc., avec un titre (« Annexe 1 : photographie du marché Sandaga, allée centrale, 15 octobre »).
\item Dans le corps du rapport, renvoyer explicitement à l'annexe par une formule du type : « (voir annexe 1) », « comme le montre le tableau de l'annexe 2 ».
\item Ne jamais placer en annexe un document qui n'est pas mentionné dans le corps, et inversement.
\end{itemize}
\end{aretenir}

\courssub{Vérifier la fiabilité d'une information}

Un rapport engage la responsabilité de son auteur : une information fausse ou invérifiée le discrédite entièrement. Avant de consigner une donnée, le rapporteur se pose trois questions :

\begin{itemize}
\item \textbf{La source} : d'où vient l'information ? Un service officiel, un témoin direct, une rumeur ? On privilégie les sources identifiables et autorisées.
\item \textbf{La date} : l'information est-elle récente et encore valable ? Un chiffre de fréquentation datant d'avant les pluies ne décrit pas la situation actuelle.
\item \textbf{La cohérence} : l'information concorde-t-elle avec les autres documents ? Si le tableau indique 900 visiteurs et qu'un témoin affirme « personne ne vient plus », on retient le chiffre et on nuance le témoignage.
\end{itemize}

En cas de doute, le rapporteur le signale honnêtement : « selon les commerçants interrogés (information à confirmer) ».

\begin{savaistu}[Documents iconiques à l'examen]
Au Probatoire et au Baccalauréat technique, l'épreuve de français peut fournir des documents iconiques (photographie, tableau, schéma) à exploiter dans la production d'écrits. Le correcteur attend que le candidat s'appuie réellement sur ces documents : citer un chiffre, décrire un élément visible, justifier une conclusion par une donnée. Un candidat qui ignore les documents fournis perd des points précieux, même si son expression est bonne.
\end{savaistu}

\courssub{Activité guidée : rédiger le paragraphe « constats »}

\begin{exoresolu}[Le marché inondé de Douala]
Énoncé. Vous êtes technicien communal à Douala. Après les fortes pluies du 12 octobre, le Maire vous charge d'inspecter le marché Sandaga. Vous disposez de deux documents : (1) une photographie montrant les allées centrales recouvertes d'eau stagnante, des étals en bois partiellement immergés et quelques commerçants déplaçant leurs marchandises ; (2) le tableau de fréquentation hebdomadaire avant et après les pluies (voir figure ci-dessus : environ 2 400 visiteurs par jour avant, 900 après). Rédigez le paragraphe « constats » du corps de votre rapport.
\tcblower
Solution détaillée. On applique la méthode : identifier, décrire, analyser, exploiter. Le paragraphe rédigé peut prendre la forme suivante :

\emph{« Au cours de ma visite du 15 octobre, j'ai constaté que les allées centrales du marché Sandaga sont recouvertes d'une eau stagnante atteignant par endroits 30 cm de hauteur, ce qui rend plusieurs sections impraticables (voir annexe 1). De nombreux étals en bois sont partiellement immergés et des marchandises ont été endommagées. Cette situation a entraîné une désaffection massive de la clientèle : la fréquentation du marché est passée d'environ 2 400 à 900 visiteurs par jour, soit une baisse de 62,5 \% (voir annexe 2). Cette baisse, constante sur l'ensemble de la semaine, confirme que l'inondation des allées constitue la cause principale du ralentissement de l'activité commerciale. »}

On observe que chaque phrase repose sur un document : la photographie pour la description, le tableau pour les chiffres, avec renvoi aux annexes et calcul du taux de variation. Aucune affirmation n'est gratuite.
\end{exoresolu}

\begin{exercice}[À vous de jouer]
À partir du tableau de fréquentation de la figure ci-dessus, rédigez deux phrases de rapport : l'une comparant les fréquentations du samedi avant et après les pluies (avec le taux de variation), l'autre dégageant la tendance générale de la semaine. Vous terminerez par un renvoi à une annexe imaginaire.
\end{exercice}

\begin{corrige}[Exercice]
Samedi avant les pluies : 2 600 visiteurs ; samedi après les pluies : 1 100 visiteurs. Taux de variation :
\[ \frac{1100 - 2600}{2600} \times 100 = \frac{-1500}{2600} \times 100 \approx -57{,}7\ \% \]
Phrases attendues : « Le samedi, jour de plus forte affluence, la fréquentation est tombée de 2 600 à 1 100 visiteurs, soit une baisse d'environ 57,7 \% (voir annexe 2). » Puis : « Sur l'ensemble de la semaine, la baisse se maintient autour de 60 \% chaque jour, ce qui traduit une désaffection durable de la clientèle et non un phénomène passager. »
\end{corrige}

\begin{aretenir}[L'essentiel de la section]
Un rapport se nourrit de documents. Le texte fonctionnel se lit en quatre questions (émetteur, destinataire, objet, organisation) ; le texte iconique se lit en quatre étapes (identifier, décrire, analyser, exploiter). Toute donnée iconique doit être transformée en phrase rédigée, chiffrée et comparée ; tout document est placé en annexe numérotée et cité dans le corps par « voir annexe n » ; toute information est vérifiée (source, date, cohérence) avant d'être consignée. Ces compétences seront mises en œuvre dans la section suivante, consacrée aux tonalités satirique et tragique, puis à la rédaction complète du rapport.
\end{aretenir}

\coursec{Les tonalités satirique et tragique}

Dans la section précédente, nous avons appris à exploiter des documents variés : textes fonctionnels (articles, lettres, comptes rendus) et textes iconiques (photographies, tableaux, schémas). Or, en lisant ces documents, vous avez sans doute remarqué qu'un même fait peut être raconté de façons très différentes : un article peut \emph{se moquer} d'un comportement, un récit peut nous plonger dans \emph{l'angoisse} face à un malheur inévitable. Cette « couleur » particulière que l'auteur donne à son texte s'appelle la \cle{tonalité} (ou \cle{registre}). Dans cette section, nous allons définir la tonalité, étudier en profondeur deux tonalités au programme — la tonalité \cle{satirique} et la tonalité \cle{tragique} —, apprendre à repérer leurs marques linguistiques, et comprendre pourquoi elles doivent être \emph{évitées} dans un rapport, texte qui exige neutralité et objectivité.

\courssub{La tonalité : définition et rappel}

\begin{definition}[La tonalité (ou registre)]
La \cle{tonalité} (ou \cle{registre}) d'un texte est l'ensemble des procédés linguistiques et stylistiques par lesquels l'auteur exprime une attitude, une émotion dominante, et cherche à produire un effet particulier sur le lecteur ou l'auditeur (rire, pitié, terreur, indignation, admiration...).
\end{definition}

\begin{aretenir}[Les tonalités déjà connues]
Au cours des années précédentes, vous avez rencontré plusieurs tonalités :
\begin{itemize}
\item la tonalité \cle{comique} : elle fait rire (comique de mots, de gestes, de situation, de caractère) ;
\item la tonalité \cle{lyrique} : elle exprime les sentiments personnels (amour, tristesse, joie), souvent à la première personne ;
\item la tonalité \cle{épique} : elle magnifie les héros et les grandes actions ;
\item la tonalité \cle{pathétique} : elle suscite la pitié, l'émotion douloureuse ;
\item la tonalité \cle{didactique} : elle cherche à instruire, à expliquer.
\end{itemize}
Nous ajoutons aujourd'hui deux tonalités au programme de Première technique : la tonalité \cle{satirique} et la tonalité \cle{tragique}.
\end{aretenir}

\courssub{La tonalité satirique : dénoncer en riant}

\textbf{Intuition.} Imaginez un élève qui arrive chaque matin avec une heure de retard. Au lieu de le gronder, le professeur déclare solennellement : « Félicitons Monsieur pour sa ponctualité légendaire : il n'est en retard que de cinquante-neuf minutes aujourd'hui, c'est un record ! » Toute la classe rit, mais le message est clair : le retard est dénoncé. C'est exactement le mécanisme de la satire.

\begin{definition}[Tonalité satirique]
La tonalité \cle{satirique} est le registre qui \emph{critique en riant} : elle dénonce un défaut, un vice, un travers ou un abus (la corruption, la paresse, l'hypocrisie, l'avarice...) par la moquerie, afin de corriger les mœurs ou de provoquer une prise de conscience. On la résume par la formule latine \emph{castigat ridendo mores} : « elle corrige les mœurs en riant ».
\end{definition}

\textbf{Les procédés de la satire.} Pour atteindre son but, la satire utilise des outils précis :

\begin{itemize}
\item l'\cle{ironie} : dire le contraire de ce que l'on pense, en laissant entendre sa véritable pensée (« Quel honnête homme ! Il n'a détourné que trois millions cette année. ») ;
\item l'\cle{antiphrase} : cas particulier de l'ironie, elle consiste à affirmer le contraire de ce qu'on veut faire comprendre (« Bravo pour ta discrétion ! » dit à quelqu'un qui crie) ;
\item la \cle{caricature} : portrait exagéré qui grossit les défauts physiques ou moraux d'une personne pour la ridiculiser ;
\item l'\cle{exagération} (ou hyperbole satirique) : amplifier un travers jusqu'à l'absurde pour le rendre visible (« Il est tellement avare qu'il économise même sa salive. ») ;
\item la \cle{récriture} parodique : imiter un style noble ou administratif pour parler de choses basses, et inversement.
\end{itemize}

\begin{exemplebox}[La satire de la corruption dans la presse camerounaise]
Un éditorialiste camerounais écrit : « Les caisses de l'État ont maigri pendant que certains ventres s'arrondissaient. »
\begin{itemize}
\item \textbf{Procédés} : antithèse (« maigri » / « s'arrondissaient »), métaphore corporelle, euphémisme (« certains » désigne sans nommer).
\item \textbf{Visée} : dénoncer le détournement des deniers publics sans procès en diffamation.
\item \textbf{Effet} : le lecteur sourit, mais il comprend et s'indigne : la satire informe \emph{et} juge.
\end{itemize}
\end{exemplebox}

\begin{savaistu}[La satire dans la tradition orale camerounaise]
La satire n'est pas une invention occidentale ! Dans la tradition orale camerounaise — contes, devinettes, proverbes, chansons des griots —, on corrige les comportements sociaux en riant. Le conte met en scène des personnages-types (l'avare, le paresseux, le glouton, souvent incarnés par des animaux comme la hyène ou l'araignée) dont les mésaventures ridiculisent le défaut. L'auditoire rit, et chacun reconnaît le travers sans que personne ne soit nommé : c'est une satire collective et pédagogique, ancêtre de nos caricatures de presse.
\end{savaistu}

\courssub{La tonalité tragique : le destin inéluctable}

\textbf{Intuition.} Dans certaines histoires, le personnage a beau lutter, tout le pousse vers sa perte : un avertissement ignoré, une faute irréparable, un enchaînement fatal d'événements. Le lecteur sent que la catastrophe \emph{doit} arriver, et il éprouve à la fois pitié pour le personnage et terreur devant la force du destin. C'est le tragique.

\begin{definition}[Tonalité tragique]
La tonalité \cle{tragique} est le registre qui représente un conflit \emph{insoluble} — entre le devoir et la passion, entre l'homme et le destin, entre deux obligations contradictoires — et qui suscite chez le lecteur la \emph{pitié} et la \emph{terreur} face à un destin inéluctable, aboutissant le plus souvent à la mort ou à la ruine du personnage.
\end{definition}

\textbf{Les procédés du tragique.}

\begin{itemize}
\item la \cle{fatalité} : le destin, les dieux, une prédiction ou une malédiction pèsent sur le personnage ; le récit annonce souvent la catastrophe à l'avance (prolepse) ;
\item la \cle{gradation} : accumulation de termes de plus en plus forts qui précipitent le lecteur vers le dénouement (« Il hésita, trembla, recula, puis s'effondra. ») ;
\item l'\cle{hyperbole} : exagération qui amplifie la douleur ou le malheur (« Tout le ciel pleurait avec elle. ») ;
\item le \cle{pathétique} : mise en scène de la souffrance pour émouvoir (plaintes, cris, scènes d'adieux) ;
\item le \cle{conflit de devoirs} : le personnage est déchiré entre deux choix également condamnables (aimer ou obéir, vivre ou rester fidèle).
\end{itemize}

\begin{exemplebox}[Le tragique dans le récit]
« Depuis la nuit où le vieux devin avait parlé, Kanga savait. Chaque matin, il regardait le fleuve monter ; chaque soir, il entendait les tambours se rapprocher. Il pria, il supplia, il offrit ses chèvres, son champ, sa maison : rien n'y fit. Le jour venu, il marcha seul vers la rive, et le village entier retenait son souffle. »
\begin{itemize}
\item \textbf{Marques} : prédiction initiale (fatalité), gradation (« il pria, il supplia, il offrit »), anaphore (« chaque matin... chaque soir »), solennité du dénouement.
\item \textbf{Effet} : pitié pour Kanga, terreur devant l'inévitable.
\end{itemize}
\end{exemplebox}

\courssub{Repérer les marques linguistiques : tableau comparatif}

\begin{methode}[Identifier une tonalité en quatre étapes]
\begin{enumerate}
\item \textbf{Relever les procédés} : cherche l'ironie, l'antiphrase, la caricature (satire) ; la gradation, l'hyperbole, les annonces du destin (tragique).
\item \textbf{Dégager le champ lexical dominant} : champ du rire, du ridicule, du vice (satire) ; champ de la mort, du destin, de la souffrance (tragique).
\item \textbf{Identifier l'attitude de l'auteur} : se moque-t-il pour corriger ? plaint-il un personnage écrasé par le sort ?
\item \textbf{Nommer l'effet produit} sur le lecteur : rire critique et indignation (satire) ; pitié et terreur (tragique).
\end{enumerate}
Conclus toujours par une phrase du type : « Ce texte est de tonalité satirique car... », en citant au moins deux procédés relevés.
\end{methode}

\begin{popfigure}
\begin{center}
\begin{tikzpicture}[font=\small]
% Cadres du tableau
\fill[popblueL] (0,3.4) rectangle (12.6,4.3);
\fill[poporangeL] (0,1.9) rectangle (12.6,3.4);
\fill[poppurpleL] (0,0.4) rectangle (12.6,1.9);
\draw[popline, thick] (0,0.4) rectangle (12.6,4.3);
\draw[popline] (0,3.4)--(12.6,3.4);
\draw[popline] (0,1.9)--(12.6,1.9);
\draw[popline] (2.6,0.4)--(2.6,4.3);
\draw[popline] (7.6,0.4)--(7.6,4.3);
% En-têtes
\node[font=\small\bfseries, text=popdark] at (1.3,3.85) {Critère};
\node[font=\small\bfseries, text=popdark] at (5.1,3.85) {Tonalité satirique};
\node[font=\small\bfseries, text=popdark] at (10.1,3.85) {Tonalité tragique};
% Ligne visée
\node[font=\small\bfseries, text=popdark, align=center] at (1.3,2.65) {Visée};
\node[align=left, text width=4.6cm] at (5.1,2.65) {Dénoncer un vice, un abus par la moquerie ; corriger les mœurs};
\node[align=left, text width=4.6cm] at (10.1,2.65) {Susciter pitié et terreur devant un destin inéluctable};
% Ligne procédés
\node[font=\small\bfseries, text=popdark, align=center] at (1.3,1.15) {Procédés et marques};
\node[align=left, text width=4.6cm] at (5.1,1.15) {Ironie, antiphrase, caricature, exagération ; champ lexical du rire et du ridicule};
\node[align=left, text width=4.6cm] at (10.1,1.15) {Gradation, hyperbole, pathétique, fatalité ; champ lexical de la mort et du destin};
\end{tikzpicture}
\end{center}
\legende{Tableau comparatif des tonalités satirique et tragique : visée, procédés et marques linguistiques.}
\end{popfigure}

\begin{popfigure}
\begin{center}
\begin{tikzpicture}[font=\small]
% Encadré du texte
\draw[popblue, thick, rounded corners, fill=popcream] (0,0) rectangle (9.2,3.2);
\node[align=left, text width=8.4cm] at (4.6,1.6)
{\guillemotleft~Notre honorable directeur est d'une générosité sans bornes : il a offert à l'école un tableau neuf... payé avec l'argent des fournitures des élèves. Quant aux bancs promis, ils voyagent encore, sans doute en première classe.~\guillemotright};
% Annotations avec flèches
\node[align=center, text=poporange, font=\small\bfseries] (ironie) at (11.9,2.7) {ironie /\\ antiphrase};
\draw[-{Stealth}, poporange, thick] (ironie.west) to[bend right=10] (7.2,2.55);
\node[align=center, text=poppurple, font=\small\bfseries] (caric) at (11.9,1.3) {caricature du\\ faux généreux};
\draw[-{Stealth}, poppurple, thick] (caric.west) to[bend left=5] (6.6,1.75);
\node[align=center, text=popgreen, font=\small\bfseries] (exag) at (11.9,-0.1) {exagération\\ comique};
\draw[-{Stealth}, popgreen, thick] (exag.west) to[bend left=10] (5.9,0.7);
\end{tikzpicture}
\end{center}
\legende{Extrait satirique annoté (chronique scolaire camerounaise) : les flèches désignent les procédés de la satire.}
\end{popfigure}

\begin{attention}[Erreur fréquente : ironie ou humour ?]
Ne confonds pas l'\cle{ironie} et l'\cle{humour}. L'humour fait rire \emph{sans visée critique} : un jeu de mots, une situation cocasse suffisent. L'ironie, elle, dit \emph{le contraire} de ce qu'on pense \emph{pour dénoncer}. Toute ironie peut faire sourire, mais tout rire n'est pas satirique. Pour qu'il y ait satire, il faut : le rire \textbf{+} une cible (un vice, un abus) \textbf{+} une intention de correction. De même, ne confonds pas le tragique (destin inéluctable, conflit insoluble) avec le simple triste ou le pathétique (émotion douloureuse sans fatalité).
\end{attention}

\courssub{Emploi raisonné : où utiliser (et ne pas utiliser) ces tonalités}

\textbf{Où les trouve-t-on ?} La tonalité satirique est à sa place dans un \emph{article de presse d'opinion}, un \emph{éditorial}, une \emph{chronique}, une \emph{caricature}, un \emph{discours} polémique ou une pièce de théâtre de critique sociale : elle informe tout en jugeant. La tonalité tragique appartient au \emph{récit}, au \emph{théâtre}, à la \emph{poésie élégiaque} : elle émeut et fait réfléchir sur la condition humaine.

\textbf{Pourquoi les éviter dans un rapport ?} Le rapport est un \cle{texte fonctionnel} : sa fonction est d'\emph{informer objectivement} un destinataire (chef de service, commission, autorité) pour l'aider à décider. Or :
\begin{itemize}
\item la satire \emph{juge et ridiculise} : elle introduit la subjectivité, expose au reproche de partialité, voire à des poursuites ;
\item le tragique \emph{dramatise} : il amplifie les faits au lieu de les exposer, et nuit à la crédibilité du document.
\end{itemize}
Dans un rapport, on écrit donc dans un registre \cle{neutre} et \cle{objectif} : faits vérifiables, chiffres, vocabulaire précis, phrases déclaratives, absence de moquerie et de dramatisation.

\begin{exoresolu}[Transformer un passage satirique en paragraphe de rapport]
\textbf{Énoncé.} Voici un extrait de chronique scolaire : « Notre honorable cantine est un modèle du genre : les élèves y font un régime minceur gratuit, puisque sur trois cents inscrits, deux cents repartent l'assiette à moitié pleine... de riz sans sauce. Bravo aux cuisiniers, ces grands économistes ! » Réécris ce passage en paragraphe de rapport neutre et objectif.
\tcblower
\textbf{Solution détaillée.}
Étape 1 — Je repère les procédés satiriques : ironie (« modèle du genre », « Bravo »), antiphrase (« grands économistes »), exagération (« régime minceur gratuit »).

Étape 2 — J'extrais les \emph{faits} derrière la moquerie : la nourriture servie est insuffisante en quantité et en qualité ; environ deux tiers des élèves (200 sur 300, soit environ 66,7 \%) ne mangent pas à leur faim ; le riz est servi sans sauce.

Étape 3 — Je reformule dans un registre neutre, avec chiffres et constats vérifiables :

« Lors de notre visite du 12 mars, nous avons constaté que la cantine scolaire ne répond pas aux besoins des élèves. Sur trois cents inscrits, environ deux cents élèves, soit les deux tiers, quittent le réfectoire sans avoir terminé leur repas. Le menu observé se limitait à du riz servi sans sauce. Nous recommandons une révision des quantités préparées et un contrôle de la qualité des repas. »

\textbf{Bilan} : les faits sont conservés, les chiffres précisés, la moquerie supprimée, une recommandation ajoutée — c'est l'écriture du rapport.
\end{exoresolu}

\begin{exercice}[Repérage et transformation]
1. Relis l'extrait annoté de la figure 2 (« Notre honorable directeur... »). Releve dans le texte deux procédés satiriques et cite les mots exacts qui les réalisent.

2. Indique la tonalité dominante de ce passage et justifie : « Maudit soit le jour où j'ai vendu la terre de mes pères ! Depuis, le malheur marche avec moi : j'ai perdu ma femme, puis mon fils, puis ma vue. Demain, je le sais, c'est moi que la terre reprendra. »

3. Transforme la phrase satirique suivante en phrase de rapport : « Le service des permis bat des records de rapidité : trois mois pour tamponner une feuille ! »
\end{exercice}

\begin{corrige}[Exercice : Repérage et transformation]
1. Procédés attendus (au choix) : l'ironie/antiphrase (« honorable », « générosité sans bornes ») ; la caricature (le directeur faux généreux qui se paye sur l'argent des élèves) ; l'exagération comique (« voyagent encore, sans doute en première classe »).

2. Tonalité tragique. Justification : champ lexical de la mort et du malheur (« maudit », « malheur », « perdu », « la terre reprendra ») ; fatalité annoncée (« Demain, je le sais ») ; gradation des pertes (« ma femme, puis mon fils, puis ma vue ») ; effet de pitié et de terreur.

3. Phrase de rapport possible : « Le service des permis accuse des délais de traitement excessifs : l'instruction d'un dossier simple nécessite environ trois mois. Nous recommandons une réorganisation du circuit de validation. »
\end{corrige}

\begin{aretenir}[L'essentiel]
\begin{itemize}
\item La \cle{tonalité satirique} critique en riant : ironie, antiphrase, caricature, exagération ; visée : corriger les mœurs.
\item La \cle{tonalité tragique} met en scène un destin inéluctable : fatalité, gradation, hyperbole, pathétique ; effet : pitié et terreur.
\item Pour identifier une tonalité : procédés + champ lexical + attitude de l'auteur + effet sur le lecteur.
\item Ces deux tonalités sont légitimes dans la presse d'opinion et les récits, mais \emph{interdites} dans un rapport, qui exige un registre neutre et objectif.
\end{itemize}
\end{aretenir}

\coursec{De la situation au plan : préparer et rédiger le rapport}

Dans la section précédente, nous avons étudié les tonalités satirique et tragique, qui concernent surtout les textes littéraires : on y joue avec les émotions du lecteur, on y cherche à faire rire ou à émouvoir. Le rapport, lui, appartient à une tout autre famille de textes : le texte \cle{fonctionnel}. Ici, point d'ironie ni de pathos : le rapporteur doit être \cle{objectif}, clair et ordonné. Mais cette objectivité ne s'improvise pas. Un bon rapport est le fruit d'un \cle{travail préparatoire} méthodique : analyser la situation, collecter les informations, bâtir un plan, rédiger, puis relire et améliorer. C'est tout ce cheminement, de la situation de départ jusqu'à la version finale, que nous allons maintenant suivre pas à pas.

\courssub{Analyser la situation de communication}

\begin{definition}[La situation de communication du rapport]
La \cle{situation de communication} d'un rapport est l'ensemble des circonstances dans lesquelles le rapport est commandé et rédigé. Elle se résume en quatre questions : \textbf{qui} commande le rapport ? \textbf{sur quoi} porte-t-il ? \textbf{pour qui} est-il rédigé ? \textbf{dans quel délai} doit-il être rendu ?
\end{definition}

Imaginez qu'un chef d'atelier vous demande de « faire un petit compte de ce qui s'est passé hier ». Avant même d'écrire la moindre phrase, vous devez clarifier la commande. C'est exactement ce que fait tout professionnel : un rapport flou naît presque toujours d'une situation mal analysée.

\begin{methode}[Analyser la situation en quatre questions]
\begin{enumerate}
\item \textbf{Qui commande ?} Identifier l'autorité qui demande le rapport (le proviseur, le chef d'atelier, le maire, le chef d'entreprise). Cela détermine le \cle{destinataire} et le niveau de langue : on s'adresse à un supérieur avec déférence et précision.
\item \textbf{Sur quoi ?} Délimiter précisément l'objet du rapport : un incident, une visite, un stage, une panne, une activité. Tout ce qui sort de cet objet doit être écarté.
\item \textbf{Pour qui ?} Déterminer le lecteur : est-ce la personne qui commande, ou un tiers (un ministère, une entreprise partenaire) ? Le lecteur conditionne les informations à donner et celles qu'on peut sous-entendre.
\item \textbf{Dans quel délai ?} Noter la date de remise. Un rapport rendu en retard, même excellent, perd sa valeur : il sert à prendre une décision, souvent urgente.
\end{enumerate}
\end{methode}

\begin{exemplebox}[Analyse d'une situation concrète]
Situation : « Le proviseur du lycée technique de Bafoussam vous demande un rapport sur la panne d'électricité survenue mardi dans les ateliers, à remettre vendredi, pour être transmis à la délégation régionale. »
\begin{itemize}
\item Qui commande ? Le proviseur.
\item Sur quoi ? La panne d'électricité de mardi dans les ateliers (pas les pannes antérieures, pas les problèmes d'eau).
\item Pour qui ? Le proviseur d'abord, puis la délégation régionale : il faut donc des faits vérifiables et des chiffres.
\item Délai ? Vendredi : trois jours pour observer, interroger et rédiger.
\end{itemize}
\end{exemplebox}

\courssub{Collecter et trier les informations}

Une fois la situation comprise, le rapporteur devient enquêteur. Il rassemble des \cle{informations} de natures différentes, qu'il faudra ensuite trier.

\begin{definition}[Les sources d'information du rapport]
Le rapporteur mobilise quatre grandes sources : les \cle{observations} directes (ce qu'il a vu et constaté lui-même), les \cle{témoignages} (ce que les personnes présentes déclarent), les \cle{données chiffrées} (mesures, dates, durées, quantités) et les \cle{documents} (textes fonctionnels : lettres, factures, registres ; documents iconiques : photographies, schémas, plans).
\end{definition}

\begin{experience}[Mener l'enquête sur le terrain]
Protocole de collecte pour le rapport sur la panne :
\begin{enumerate}
\item \textbf{Observer} : se rendre dans les ateliers, noter l'état des installations (tableau électrique, câbles, machines à l'arrêt).
\item \textbf{Interroger} : recueillir les témoignages du surveillant d'atelier, de deux élèves présents, du technicien appelé en réparation. Poser des questions factuelles : quand ? où ? qu'avez-vous vu ?
\item \textbf{Relever les chiffres} : heure du début de la panne (9 h 15), heure du rétablissement (14 h 30), nombre de machines touchées (6), nombre d'élèves privés de cours pratiques (environ 120).
\item \textbf{Photographier ou schématiser} : un document iconique (photo du tableau électrique endommagé, croquis de l'atelier) peut être annexé au rapport.
\end{enumerate}
Observation : on obtient une masse d'informations désordonnée, mêlant faits certains et rumeurs.
Interprétation : le tri est indispensable. On distingue les \cle{faits} vérifiables (« le disjoncteur général a sauté à 9 h 15 ») des \cle{opinions} (« c'est la faute de l'entretien ») : seuls les faits ont leur place dans le corps du rapport ; les hypothèses doivent être signalées comme telles.
\end{experience}

\begin{attention}[Piège : mélanger faits et opinions]
Une erreur fréquente consiste à écrire « l'installation électrique du lycée est lamentable » au lieu de « trois câbles du tableau général présentent une usure visible ». La première phrase est un jugement subjectif ; la seconde est un constat objectif, vérifiable. Dans un rapport, on \textbf{montre} les faits et on laisse le lecteur juger.
\end{attention}

\courssub{Élaborer le plan : choisir un ordre}

Les informations triées doivent maintenant être \cle{classées}. C'est l'étape décisive : le plan est le squelette du rapport.

\begin{definition}[Les trois ordres possibles]
Le plan d'un rapport suit généralement l'un de ces trois ordres :
\begin{itemize}
\item l'ordre \cle{chronologique} : on raconte les faits dans l'ordre où ils se sont produits (avant, pendant, après) ;
\item l'ordre \cle{thématique} : on regroupe les informations par thèmes (état des lieux, moyens humains, moyens matériels) ;
\item l'ordre \cle{problème-solutions} : on expose les constats, puis les causes, puis les conséquences, puis les recommandations. C'est l'ordre le plus fréquent dans les rapports d'incident.
\end{itemize}
\end{definition}

\begin{methode}[Construire le plan en trois gestes]
\begin{enumerate}
\item \textbf{Classer} les informations collectées en 2 à 4 grandes parties, selon l'ordre choisi. Chaque partie doit répondre à une question simple : que s'est-il passé ? pourquoi ? avec quelles conséquences ? que faut-il faire ?
\item \textbf{Subdiviser} chaque partie en paragraphes : un paragraphe = une seule idée, annoncée par une phrase-clé.
\item \textbf{Vérifier} que le plan couvre tout l'objet du rapport, sans répétition et sans hors-sujet : chaque information collectée doit trouver sa place, et chaque partie doit servir la commande.
\end{enumerate}
\end{methode}

\begin{exemplebox}[Plan d'un rapport sur la panne d'électricité au lycée technique de Bafoussam]
Ordre choisi : problème-solutions, en quatre parties.
\begin{enumerate}
\item \textbf{Constats} : panne survenue mardi à 9 h 15 dans les ateliers ; six machines hors service ; tableau général endommagé.
\item \textbf{Causes} : usure de trois câbles du tableau général ; surcharge due au branchement simultané des machines de l'atelier de mécanique.
\item \textbf{Conséquences} : interruption des cours pratiques de 9 h 15 à 14 h 30, soit 5 h 15 d'arrêt ; environ 120 élèves privés de travaux pratiques ; retard dans le programme de deux classes.
\item \textbf{Recommandations} : remplacer les câbles usés ; faire contrôler l'installation par un technicien agréé ; répartir les horaires d'utilisation des machines.
\end{enumerate}
\end{exemplebox}

\begin{attention}[Erreur fréquente : rédiger sans plan]
Beaucoup d'élèves se jettent directement dans la rédaction. Résultat : un texte \textbf{désordonné} (les causes apparaissent après les recommandations), \textbf{redondant} (la même information revient deux ou trois fois) et \textbf{incomplet} (une conséquence importante est oubliée). Cinq minutes consacrées au plan font gagner trente minutes de rédaction et évitent une mauvaise note.
\end{attention}

\begin{popfigure}
\begin{center}
\begin{tikzpicture}[
node distance=0.55cm,
etape/.style={rectangle, rounded corners=3pt, draw=popblue, fill=popblueL, text width=3.1cm, align=center, minimum height=0.9cm, font=\small},
fleche/.style={-{Stealth[length=2.5mm]}, thick, popdark}]
\node[etape, align=center] (s){\textbf{1. Situation}\\ qui ? sur quoi ? pour qui ? délai ?};
\node[etape, below=of s, align=center] (c){\textbf{2. Collecte}\\ observations, témoignages, chiffres, documents};
\node[etape, below=of c, align=center] (p){\textbf{3. Plan}\\ classer en 2 à 4 parties ordonnées};
\node[etape, below=of p, align=center] (r){\textbf{4. Rédaction}\\ introduction, corps, conclusion};
\node[etape, below=of r, align=center] (re){\textbf{5. Relecture}\\ grille de correction, lecture croisée};
\node[etape, below=of re, fill=popgreenL, draw=popgreen, align=center] (f){\textbf{6. Version finale}\\ rapport amélioré et remis};
\draw[fleche] (s) -- (c);
\draw[fleche] (c) -- (p);
\draw[fleche] (p) -- (r);
\draw[fleche] (r) -- (re);
\draw[fleche] (re) -- (f);
\draw[fleche, poporange, dashed] (re.east) to[out=0,in=0] node[right, font=\scriptsize, align=left, popdark]{retour en\\ arrière si\\ nécessaire} (p.east);
\end{tikzpicture}
\end{center}
\legende{Organigramme des six étapes de production du rapport. La flèche en pointillés rappelle que la relecture peut renvoyer à une étape antérieure : on corrige le plan avant de corriger les phrases.}
\end{popfigure}

\courssub{Rédiger l'introduction et la conclusion à partir du plan}

Le plan étant fixé, la rédaction peut commencer. On rédige d'abord l'introduction et la conclusion, car elles encadrent tout le propos.

\begin{definition}[Contenu de l'introduction et de la conclusion]
L'\cle{introduction} du rapport présente les circonstances de la commande et l'objet du rapport : qui a demandé le rapport, quand, à quel sujet, et comment l'enquête a été menée. La \cle{conclusion} résume les constats essentiels et formule les recommandations ou les souhaits ; elle ne contient \textbf{aucune information nouvelle}.
\end{definition}

\begin{exemplebox}[Introduction et conclusion rédigées]
\textbf{Introduction} : « Par note de service en date du 14 mars, Monsieur le Proviseur m'a chargé d'établir un rapport sur la panne d'électricité survenue le mardi 13 mars dans les ateliers du lycée technique de Bafoussam. Pour mener à bien cette mission, j'ai procédé à des observations sur place et recueilli les témoignages du surveillant d'atelier et du technicien de maintenance. »

\textbf{Conclusion} : « En définitive, la panne du 13 mars, due à l'usure de l'installation et à une surcharge, a privé environ 120 élèves de travaux pratiques pendant plus de cinq heures. Il apparaît urgent de remplacer les câbles défectueux et de faire contrôler l'ensemble de l'installation. Je reste à la disposition de Monsieur le Proviseur pour tout complément d'information. »
\end{exemplebox}

\begin{attention}[Pièges de l'introduction et de la conclusion]
Dans l'introduction, ne racontez pas déjà toute l'affaire : annoncez seulement la commande, l'objet et la méthode. Dans la conclusion, n'introduisez jamais un fait nouveau (« j'ai d'ailleurs appris que... ») : une information nouvelle en conclusion est le signe d'un plan mal construit ; elle doit être replacée dans le corps.
\end{attention}

\courssub{Rédiger le corps du rapport}

Le corps développe le plan, partie après partie, paragraphe après paragraphe.

\begin{propriete}[Règles de rédaction du corps]
\begin{itemize}
\item \textbf{Un paragraphe par idée} : chaque paragraphe commence par une phrase-clé qui annonce l'idée, suivie de phrases qui la précisent par des faits et des chiffres.
\item \textbf{Des phrases courtes} : une phrase = une information. Les longues phrases embrouillent le lecteur.
\item \textbf{Un vocabulaire précis} : « le disjoncteur général » et non « le truc » ; « à 9 h 15 » et non « dans la matinée ».
\item \textbf{Les temps du passé} : le \cle{passé composé} pour les faits ponctuels et datés (« la panne a commencé à 9 h 15 ») ; l'\cle{imparfait} pour le décor, les descriptions et les actions en cours (« les élèves travaillaient sur les tours lorsque les machines se sont arrêtées »).
\end{itemize}
\end{propriete}

\begin{exoresolu}[Rédiger un paragraphe du corps]
Énoncé. À partir des informations suivantes, rédigez le paragraphe « Causes » du rapport sur la panne : câbles usés (3), surcharge, branchement simultané des machines de mécanique, témoignage du technicien.
\tcblower
Solution détaillée. On commence par une phrase-clé, puis on enchaîne les faits au passé composé et à l'imparfait :

« L'enquête a permis d'identifier deux causes à cette panne. D'une part, le technicien de maintenance a constaté que trois câbles du tableau général présentaient une usure avancée : leur gaine était fissurée et les fils étaient visibles. D'autre part, au moment de la panne, toutes les machines de l'atelier de mécanique fonctionnaient simultanément, ce qui a provoqué une surcharge de l'installation. Le disjoncteur général a alors sauté, privant d'électricité l'ensemble des ateliers. »

On remarque : une seule idée par phrase, des chiffres précis (trois câbles), le passé composé pour les événements (« a constaté », « a sauté ») et l'imparfait pour l'état et le décor (« était fissurée », « fonctionnaient »).
\end{exoresolu}

\courssub{Relire, confronter et améliorer la production}

Un premier jet n'est jamais un rapport fini. La dernière étape, prévue par le programme officiel, est la \cle{confrontation} des productions et leur amélioration.

\begin{methode}[Améliorer sa production en quatre vérifications]
\begin{enumerate}
\item \textbf{Structure externe} : l'en-tête (lieu, date, objet), les titres des parties, la signature sont-ils présents et correctement placés ?
\item \textbf{Enchaînement logique} : les parties suivent-elles l'ordre du plan ? Les connecteurs (« d'abord », « ensuite », « enfin », « en définitive ») guident-ils le lecteur ?
\item \textbf{Objectivité} : ai-je supprimé tous les jugements personnels ? Chaque affirmation est-elle un fait vérifiable ?
\item \textbf{Correction langagière} : orthographe, accords, ponctuation, emploi correct du passé composé et de l'imparfait.
\end{enumerate}
\end{methode}

\begin{popfigure}
\begin{center}
\begin{tabular}{|l|c|c|c|}
\hline
\rowcolor{popblueL} \textbf{Critère} & \textbf{Oui} & \textbf{Partiellement} & \textbf{Non} \\
\hline
Structure externe respectée & & & \\
\hline
Introduction : commande et objet annoncés & & & \\
\hline
Plan logique (ordre clair, pas de hors-sujet) & & & \\
\hline
Un paragraphe = une idée & & & \\
\hline
Faits objectifs, chiffres précis & & & \\
\hline
Temps du passé correctement employés & & & \\
\hline
Conclusion sans information nouvelle & & & \\
\hline
Orthographe et ponctuation soignées & & & \\
\hline
\end{tabular}
\end{center}
\legende{Grille d'évaluation pour la correction croisée des rapports. Chaque binôme coche les cases en lisant le texte de l'autre, puis annote en marge les passages à retravailler.}
\end{popfigure}

\begin{methode}[La lecture croisée en binômes]
\begin{enumerate}
\item \textbf{Échanger} les rapports au sein du binôme : chacun lit le texte de l'autre, en silence, une première fois sans écrire.
\item \textbf{Annoter} à la seconde lecture : souligner les passages confus, cocher la grille d'évaluation, noter en marge une suggestion précise (« chiffre manquant », « opinion à supprimer », « accord à revoir »).
\item \textbf{S'expliquer} : chaque élève justifie ses annotations à son camarade, en citant le passage concerné.
\item \textbf{Réécrire} : chacun reprend son propre rapport et produit une \cle{version améliorée}, en tenant compte des remarques acceptées.
\end{enumerate}
\end{methode}

\begin{savaistu}[Critiquer le texte d'autrui améliore le sien]
La confrontation des productions entre pairs n'est pas une simple activité de classe : elle est explicitement prévue par le programme officiel de Français en Première technique. Les chercheurs en didactique de l'écriture ont montré qu'un élève qui apprend à repérer les défauts du texte d'un camarade devient capable de les repérer dans le sien. C'est le même principe que dans les métiers techniques : le contrôle qualité d'une pièce fabriquée par un collègue rend plus exigeant sur sa propre fabrication. Au Baccalauréat technique, cette habitude de relecture fait souvent la différence entre une copie moyenne et une bonne copie.
\end{savaistu}

\begin{exercice}[De la situation à la version finale]
Situation : le directeur de l'entreprise où tu effectues ton stage te demande un rapport sur un accident de machine survenu dans l'atelier, à remettre dans deux jours.
\begin{enumerate}
\item Analyse la situation de communication en répondant aux quatre questions.
\item Dresse la liste des informations à collecter et des personnes à interroger.
\item Élabore un plan en trois ou quatre parties (ordre problème-solutions).
\item Rédige l'introduction et la conclusion du rapport.
\item Rédige le paragraphe « Constats » en respectant les règles du corps (une idée par paragraphe, phrases courtes, temps du passé).
\item Échange ta production avec un camarade, remplis la grille d'évaluation, puis produis ta version améliorée.
\end{enumerate}
\end{exercice}

\begin{corrige}[Exercice : De la situation à la version finale]
\begin{enumerate}
\item Qui commande ? Le directeur de l'entreprise. Sur quoi ? L'accident de machine dans l'atelier. Pour qui ? Le directeur (et éventuellement l'assurance ou l'inspection du travail). Délai ? Deux jours.
\item Informations à collecter : date et heure de l'accident, machine concernée, nature des blessures, état de la machine, circonstances. Personnes à interroger : l'ouvrier blessé, le chef d'atelier, les témoins.
\item Plan : 1. Constats (faits, blessures, dégâts) ; 2. Causes (défaut de sécurité, erreur de manipulation) ; 3. Conséquences (arrêt de production, hospitalisation) ; 4. Recommandations (réparer le carter de protection, former le personnel).
\item Introduction : « À la suite de l'accident survenu le [date] sur la fraiseuse de l'atelier, Monsieur le Directeur m'a chargé d'établir un rapport sur cet événement. J'ai procédé à des constatations sur place et recueilli les témoignages du chef d'atelier et de deux ouvriers présents. » Conclusion : « Cet accident, causé par l'absence de carter de protection, a entraîné l'arrêt de la fraiseuse et l'hospitalisation d'un ouvrier. Il est recommandé de remettre en état les dispositifs de sécurité et de rappeler les consignes au personnel. »
\item Paragraphe « Constats » : « L'accident s'est produit le mardi à 10 h 40 sur la fraiseuse numéro 2. L'ouvrier qui l'utilisait a été blessé à la main gauche. La machine était en marche et son carter de protection était retiré au moment des faits. »
\item La lecture croisée suit la grille : structure, logique, objectivité, langue. La version améliorée corrige les points signalés (par exemple : remplacer « la machine dangereuse » par « la machine dont le carter était retiré »).
\end{enumerate}
\end{corrige}

\begin{aretenir}[L'essentiel de la section]
\begin{itemize}
\item Avant d'écrire, analyser la situation : \textbf{qui ? sur quoi ? pour qui ? dans quel délai ?}
\item Collecter et trier les informations : faits vérifiables d'un côté, opinions de l'autre.
\item Construire un plan en 2 à 4 parties ordonnées (chronologique, thématique ou problème-solutions) : jamais de rédaction sans plan.
\item Rédiger : introduction (commande, objet, méthode), corps (un paragraphe par idée, phrases courtes, passé composé et imparfait), conclusion (bilan et recommandations, rien de nouveau).
\item Relire avec la grille, confronter en binômes, annoter, puis produire une version améliorée.
\end{itemize}
\end{aretenir}

\coursec{Compte rendu, remédiation et évaluation}

Après avoir conduit les élèves de l'analyse de la situation à la rédaction effective de l'introduction, du corps et de la conclusion (\emph{cf.} section précédente « De la situation au plan : préparer et rédiger le rapport »), il est indispensable de ménager un temps de \textbf{regard rétrospectif} et de \textbf{consolidation}. Cette séance finale vise trois objectifs indissociables : faire le bilan collectif des acquis, corriger les erreurs persistantes par des activités de remédiation ciblées, et évaluer les compétences dans une situation authentique proche des exigences du Probatoire.

\courssub{Bilan collectif des acquis : structure et rédaction du rapport}

Avant toute remédiation, l'enseignant guide la classe dans une \textbf{réactivation synthétique} des savoirs et savoir-faire travaillés au fil de la séquence. Ce n'est pas une simple énumération, mais une mise en cohérence : chaque élément de structure répond à une fonction communicationnelle précise.

\begin{aretenir}[Synthèse des acquis attendus en fin de séquence]
\textbf{Structure externe} : Page de garde (identification complète), Sommaire, Liste des abréviations/sigles (si nécessaire), Corps du rapport, Annexes.
\textbf{Structure interne} :
\begin{itemize}
    \item \textbf{Introduction} : Accroche contextuelle $\rightarrow$ Objet (verbe d'action + nom) $\rightarrow$ Intérêt/Portée $\rightarrow$ Démarche/Plan annoncé.
    \item \textbf{Corps} : Parties équilibrées (2 à 3), titres parlants, arguments étayés par des faits/chiffres (documents), connecteurs logiques, objectivité (ton neutre, impersonnel).
    \item \textbf{Conclusion} : Bilan synthétique (réponse à l'objet) $\rightarrow$ Recommandations concrètes, hiérarchisées, réalistes $\rightarrow$ Ouverture/Perspective.
\end{itemize}
\textbf{Langue} : Concordance des temps (passé composé/imparfait pour les faits, présent pour les constats, futur/conditionnel pour les recommandations), vocabulaire technique précis, phrases nominalisées pour la concision.
\end{aretenir}

L'enseignant peut projeter ou distribuer un \textbf{rapport modèle} (anonymisé) et demander aux élèves, par groupes, d'identifier les composantes listées ci-dessus. Cette activité de \emph{lecture à visée analytique} renforce la méta-cognition.

\courssub{Repérage des difficultés récurrentes et remédiation ciblée}

L'analyse des productions d'élèves (devoirs, brouillons, évaluations formatives) fait émerger des \textbf{erreurs systémiques}. Les traiter collectivement par la remédiation est plus efficace que la correction individuelle isolée.

\begin{popfigure}
\centering
\begin{tikzpicture}[
    cell/.style={rectangle, draw=popline, minimum height=1.1cm, align=center, text width=5.2cm, font=\small},
    header/.style={fill=popblue, text=white, font=\small\bfseries, minimum height=0.9cm},
    row1/.style={fill=popblueL!50},
    row2/.style={fill=popcream},
]
% Header
\node[header, text width=3.5cm] (h1) at (0,0) {Erreur fréquente};
\node[header, text width=5.2cm] (h2) at (3.5,0) {Cause probable};
\node[header, text width=5.2cm] (h3) at (8.7,0) {Correction attendue / Stratégie de remédiation};

% Row 1
\node[cell, row1, anchor=north west, align=center] (r1c1) at (0,-0.9){Objet mal formulé :\\ « Faire le rapport des activités »};
\node[cell, row1, align=center] (r1c2) at (3.5,-0.9){Confusion objet/sujet ;\\ verbe vague (faire, faire le point) ;\\ absence de finalité.};
\node[cell, row1, align=center] (r1c3) at (8.7,-0.9){\textbf{Modèle} : Verbe d'action précis + Nom de la mission + Finalité.\\ Ex : « \emph{Présenter le bilan des activités de la coopérative 2023-2024 pour orienter le budget 2025} »\.\\ \textbf{Exercice} : Transformer 5 sujets en objets.};

% Row 2
\node[cell, row2, anchor=north west, align=center] (r2c1) at (0,-2.0){Recommandations absentes,\\ vagues ou irréalistes};
\node[cell, row2, align=center] (r2c2) at (3.5,-2.0){Oubli de la fonction prescriptive ;\\ confusion conclusion/bilan ;\\ manque de pragmatisme.};
\node[cell, row2, align=center] (r2c3) at (8.7,-2.0){\textbf{Règle} : Une recommandation = Verbe d'action (Renforcer, Créer, Réviser) + Objet + Délai/Acteur + Indicateur.\\ \textbf{Activité} : À partir d'un bilan fictif, rédiger 3 recommandations SMART.};

% Row 3
\node[cell, row1, anchor=north west, align=center] (r3c1) at (0,-3.1){Subjectivité / Ton inapproprié :\\ « Je pense que c'est bien »,\\ « C'est nul »};
\node[cell, row1, align=center] (r3c2) at (3.5,-3.1){Non-maîtrise de l'impersonnalité ;\\ vocabulaire affectif ;\\ absence de distance critique.};
\node[cell, row1, align=center] (r3c3) at (8.7,-3.1){\textbf{Technique} : Passif / « Il convient de » / « L'analyse révèle que »\.\\ Remplacer « Je trouve » par « Il ressort que »\.\\ \textbf{Exercice} : Réécrire 5 phrases subjectives en style objectif.};

% Row 4
\node[cell, row2, anchor=north west, align=center] (r4c1) at (0,-4.2){Fautes de concordance\\ des temps (Cdt)};
\node[cell, row2, align=center] (r4c2) at (3.5,-4.2){Mélange récit/rétrospectif ;\\ indicateur temporel mal géré ;\\ conditionnel/futur confus.};
\node[cell, row2, align=center] (r4c3) at (8.7,-4.2){\textbf{Rappel} :\\ - Faits passés $\rightarrow$ P. Composé / Imparfait.\\ - Constats actuels $\rightarrow$ Présent.\\ - Recommandations $\rightarrow$ Futur / Conditionnel (si hypothèse).\\ \textbf{Dictée ciblée} sur un paragraphe mélangeant les 3 temps.};
\end{tikzpicture}
\legende{Tableau de remédiation : erreurs types, causes et stratégies de correction (à afficher ou distribuer).}
\end{popfigure}

\courssub{Activités de remédiation guidées}

Trois ateliers courts (10-15 min chacun) permettent de traiter ces points en classe.

\begin{exoresolu}[Réécriture d'une introduction défaillante]
\textbf{Consigne} : Voici une introduction d'élève. Identifiez les 4 manquements par rapport au canevas standard, puis réécrivez-la correctement.

\textbf{Production d'élève} :
\emph{« Bonjour Monsieur le Proviseur. Je suis le président de la coopérative. Je vais vous faire le rapport de l'année. On a vendu des cahiers et on a fait de l'argent. Mais il y a eu des vols. Je vais vous dire ce qu'on a fait. »}

\textbf{Analyse des manquements} :
\begin{enumerate}
    \item Absence d'accroche contextuelle (cadre officiel, date).
    \item Objet mal formulé (verbe « dire », pas de finalité).
    \item Ton familier (\"On a\", \"faire de l'argent\", \"Je vais vous dire\").
    \item Plan non annoncé.
\end{enumerate}

\textbf{Proposition de réécriture} :
\emph{« Monsieur le Proviseur, L'année scolaire 2023-2024 s'achève pour la Coopérative des Élèves du Lycée Technique de Bafoussam, structure essentielle d'appui à la vie scolaire. \textbf{Objet} : Présenter le bilan moral et financier des activités de la coopérative afin d'évaluer la gestion et d'orienter les décisions budgétaires de l'exercice prochain. \textbf{Intérêt} : Ce rapport permet de garantir la transparence financière attendue par l'administration et les adhérents. \textbf{Démarche} : Après avoir retracé les principales activités menées, nous analyserons les résultats financiers avant d'identifier les dysfonctionnements constatés. Nous formulerons enfin des recommandations pour une gestion plus rigoureuse. »}

\tcblower
\textbf{Apport pédagogique} : Cet exercice montre concrètement le passage du « langage de conversation » au « langage administratif ». On insiste sur la \textbf{nominalisation} (\emph{la présentation du bilan} vs \emph{je vais présenter}) et l'usage des \textbf{connecteurs de plan} (\emph{après avoir... nous analyserons... avant d'identifier... enfin}).
\end{exoresolu}

\begin{exercice}[Remise en ordre du corps de rapport]
\textbf{Consigne} : Les paragraphes suivants constituent le corps d'un rapport sur la gestion des déchets au lycée, mais ils sont mélangés. Numérotez-les de 1 à 6 pour reconstruire une progression logique (Description $\rightarrow$ Analyse $\rightarrow$ Causes $\rightarrow$ Conséquences $\rightarrow$ Solutions existantes $\rightarrow$ Perspectives).

\begin{enumerate}
    \item[A] Les poubelles de tri sélectif installées en 2022 sont inutilisées : les bacs « plastique » et « papier » contiennent des déchets mélangés.
    \item[B] Cette situation résulte de l'absence de campagne de sensibilisation auprès des nouveaux élèves et du manque de formation des agents d'entretien.
    \item[C] Le lycée produit en moyenne 150 kg de déchets par semaine, dont 60\% de matières recyclables (papiers, bouteilles) et 40\% de déchets organiques (cantine).
    \item[D] En conséquence, la totalité des déchets est enfouie sur le site municipal, privant l'établissement de revenus potentiels (vente de recyclables) et alourdissant la facture d'évacuation.
    \item[E] Une convention avec l'ONG « Vert-Cameroun » pour la collecte hebdomadaire du tri est à l'étude pour la rentrée prochaine.
    \item[F] Actuellement, seul le ramassage global est assuré par la mairie, deux fois par semaine, sans valorisation.
\end{enumerate}
\end{exercice}

\begin{corrige}[Exercice : Remise en ordre du corps de rapport]
\textbf{Ordre logique : C $\rightarrow$ A $\rightarrow$ F $\rightarrow$ B $\rightarrow$ D $\rightarrow$ E}
\begin{itemize}
    \item \textbf{C} : Description quantitative (Constat initial / Données brutes).
    \item \textbf{A} : Analyse qualitative du dispositif existant (Dysfonctionnement du tri).
    \item \textbf{F} : Description de la procédure actuelle (Gestion en place).
    \item \textbf{B} : Analyse causale (Pourquoi ça ne marche pas ?).
    \item \textbf{D} : Analyse des conséquences (Impact financier/écologique).
    \item \textbf{E} : Perspective / Solution envisagée (Ouverture vers recommandations).
\end{itemize}
\textbf{Structure des parties suggérée} :
\textbf{Partie I : État des lieux de la production et de la gestion des déchets} (C, A, F).
\textbf{Partie II : Analyse des dysfonctionnements et perspectives} (B, D, E).
\end{corrige}

\begin{exercice}[Correction de phrases : concordance des temps et objectivité]
\textbf{Consigne} : Corrigez les phrases suivantes (temps des verbes, style impersonnel, vocabulaire technique).
\begin{enumerate}
    \item Le directeur \textbf{a dit} qu'il \textbf{va} acheter de nouveaux ordinateurs l'an prochain.
    \item \textbf{Je trouve que} les élèves \textbf{n'ont pas respecté} le règlement intérieur.
    \item Nous \textbf{allons} proposer qu'on \textbf{fait} une formation pour les profs.
    \item L'année dernière, la recette \textbf{est} de 500 000 FCFA et la dépense \textbf{est} de 600 000 FCFA.
\end{enumerate}
\end{exercice}

\begin{corrige}[Exercice : Correction de phrases]
\begin{enumerate}
    \item Le directeur \textbf{a annoncé} qu'il \textbf{envisageait} (ou \textbf{prévoit} s'il s'agit d'un projet actuel ferme) l'acquisition de nouveaux ordinateurs pour l'exercice suivant.
    \item \textbf{Il ressort que} les élèves \textbf{n'ont pas respecté} le règlement intérieur. (Ou : \textbf{Le non-respect du règlement intérieur par les élèves est constaté}.)
    \item \textbf{Il est recommandé d'organiser} une session de formation à l'intention des enseignants. (Passif / Infinitif impersonnel).
    \item L'année dernière, les recettes \textbf{se sont élevées} à 500 000 FCFA alors que les dépenses \textbf{ont atteint} 600 000 FCFA. (Passé composé pour faits clos ; vocabulaire précis : \emph{s'élever à, atteindre}).
\end{enumerate}
\end{corrige}

\courssub{Réinvestissement des tonalités satirique et tragique}

Le programme de Langue française (Morphosyntaxe) impose la maîtrise des \textbf{tonalités} dans la production écrite. Bien que le rapport exige une tonalité \textbf{neutre/objective}, la capacité à basculer vers le \textbf{satirique} (dénonciation par l'ironie, l'hyperbole, le second degré) ou le \textbf{tragique} (gravité, fatalité, pathos, champ lexical de la souffrance/irréversible) est un marqueur de maturité langagière exigible au Probatoire (épreuve d'expression écrite, sujet d'invention).

\begin{definition}[Rappel : Tonalités satirique et tragique]
\textbf{Satirique} : Vise à corriger les mœurs par le rire. Procédés : Ironie (dire le contraire de sa pensée), Hyperbole (exagération), Antiphrase, Métonymie dévalorisante, Accumulation grotesque. \emph{Marqueurs} : Guillemets ironiques, adjectifs péjoratifs décalés, questions rhétoriques.
\textbf{Tragique} : Vise à émouvoir par la peinture d'une situation sans issue. Procédés : Champ lexical de la mort/fin/douleur, Rythme lent (phrases longues, points-virgules), Figures d'insistance (anaphore, gradation), Impuissance du sujet. \emph{Marqueurs} : Verbes d'état, modalisateurs de certitude négative (\emph{hélas, fatalement}), absence de perspective.
\end{definition}

\begin{exemplebox}[Exercice de style : Même fait, trois tonalités]
\textbf{Fait divers} : « La cantine du lycée a servi un repas avarié ; 30 élèves hospitalisés. »

\textbf{1. Ton Neutre (Rapport)} :
\emph{« Le 12 mars 2024, une intoxication alimentaire collective a touché 30 élèves après le déjeuner à la cantine. L'enquête administrative a révélé une rupture de la chaîne du froid pour le poisson servi. Les victimes ont été prises en charge au centre de santé du district. Des mesures conservatoires (fermeture cuisine, audit fournisseurs) ont été immédiates. »}

\textbf{2. Ton Satirique (Chronique/Édito élève)} :
\emph{« Ah, la fameuse « Spécialité du Chef » du mardi ! Ce poisson, vétéran des congélateurs depuis l'ère glaciaire, a enfin décidé de se venger. Trente estomacs vaillants ont rendu l'âme (ou presque) au centre de santé. On applaudit la performance : transformer la cantine en service de réanimation, il fallait oser. Prochain menu : « Roulette russe à la sauce tomate » ? »}

\textbf{3. Ton Tragique (Récit/Témoignage)} :
\emph{« Le soleil de midi ne réchauffait plus les plateaux, il éclairait l'irréparable. L'odeur, d'abord appétissante, portait en elle le silence de la mort. Trente destins basculèrent dans la nausée, la fièvre, l'angoisse des mères accourues. La cantine, hier refuge, devint le théâtre d'une trahison invisible. Et le poisson, ce morceau de glace empoisonné, gît toujours au fond des poubelles, témoin muet d'une négligence qui tue. »}
\end{exemplebox}

\begin{methode}[Identifier et produire une tonalité]
\begin{enumerate}
    \item \textbf{Repérer le lexique} : Mots connotés (péjoratifs/mélancoliques) vs termes techniques (neutres).
    \item \textbf{Analyser la syntaxe} : Phrases exclamatives/interrogatives (satirique) ; phrases amples, subordonnées, anaphores (tragique) ; phrases nominalisées, passives (neutre).
    \item \textbf{Identifier l'intention} : Faire rire pour corriger (satirique) ? Faire pleurer/réfléchir sur la condition humaine (tragique) ? Informer/décider (neutre).
    \item \textbf{S'exercer} : Réécrire un même paragraphe factuel (ex: panne d'électricité au lycée) dans les trois tonalités.
\end{enumerate}
\end{methode}

\courssub{Évaluation sommative : rédaction complète d'un rapport}

L'évaluation sommative doit placer l'élève en situation de \textbf{production autonome} intégrant l'ensemble de la séquence : analyse de la situation, exploitation de documents (textes fonctionnels, iconiques), planification, rédaction, relecture.

\begin{exemplebox}[Sujet type — Évaluation Sommative (Durée : 50 min)]
\textbf{Situation} : Vous êtes le Président de la Coopérative des Élèves du Lycée Technique de Bafoussam. L'Assemblée Générale de fin d'année approche. Le Proviseur vous demande de lui remettre le \textbf{rapport annuel d'activités 2023-2024} pour validation avant diffusion aux adhérents.

\textbf{Documents fournis} :
\begin{itemize}
    \item \textbf{Doc 1 (Texte fonctionnel)} : Extrait du règlement intérieur de la coopérative (missions du bureau, règles financières).
    \item \textbf{Doc 2 (Tableau chiffré)} : Bilan financier simplifié (Recettes : cotisations, ventes, subventions / Dépenses : achats stock, frais bancaires, dons, matériel). Solde : +150 000 FCFA.
    \item \textbf{Doc 3 (Texte fonctionnel)} : Courrier du fournisseur « Papeterie du Centre » signalant 3 retards de paiement sur l'année.
    \item \textbf{Doc 4 (Iconique)} : Photo de la nouvelle salle d'étude aménagée grâce aux bénéfices 2022-2023 + Photo de poubelles débordantes près de la coopérative (problème hygiène).
\end{itemize}

\textbf{Consignes} :
\begin{enumerate}
    \item Rédigez la page de garde complète.
    \item Rédigez l'introduction (accroche, objet, intérêt, plan).
    \item Rédigez le corps du rapport en 2 parties équilibrées (titres parlants).
    \item Rédigez la conclusion (bilan, 3 recommandations SMART, ouverture).
    \item Soignez la présentation (paragraphes, connecteurs, impersonnalité, concordance des temps).
\end{enumerate}
\end{exemplebox}

\begin{methode}[Gérer son temps pendant l'évaluation (50 min)]
\begin{enumerate}
    \item \textbf{Analyse de la situation \& documents (5 min)} : Identifier le destinateur, le destinataire, l'objet, la finalité. Surligner les chiffres clés et infos critiques dans les docs. Noter les dysfonctionnements (retards paiement, hygiène).
    \item \textbf{Élaboration du plan détaillé (10 min)} : Rédiger l'introduction au brouillon. Structurer le corps :
    \begin{itemize}
        \item \textbf{Partie I : Bilan d'activités et résultats financiers} (Activités menées, Analyse recettes/dépenses Doc 2, Réalisation salle d'étude Doc 4).
        \item \textbf{Partie II : Difficultés rencontrées et axes d'amélioration} (Retards paiement Doc 3, Problème hygiène Doc 4, Respect règlement Doc 1).
    \end{itemize}
    Pré-rédiger les 3 recommandations. Écrire la conclusion au brouillon.
    \item \textbf{Rédaction au net (30 min)} : Transférer proprement. Veiller à l'enchaînement logique (connecteurs). Vérifier l'accord des participes passés, la concordance des temps.
    \item \textbf{Relecture active (5 min)} : \textbf{Check-list} : Page de garde ? Objet clair ? 2 parties ? Recommandations SMART ? « Je » éliminé ? Temps verbaux cohérents ? Orthographe lexicale (vocabulaire technique) ?
\end{enumerate}
\end{methode}

\begin{attention}[Erreur fatale : Négliger la relecture]
L'expérience des jurys du Probatoire montre que la \textbf{relecture finale} est le moment où l'on récupère 2 à 3 points sur 20 (orthographe d'usage, accords, ponctuation, cohérence des titres). Beaucoup d'élèves rendent leur copie dès la fin de la rédaction. \textbf{Interdiction de rendre la copie avant la sonnerie.} Les 5 dernières minutes sont sacrées. Relisez \emph{à l'envers} (dernière phrase vers la première) pour repérer les fautes d'inattention.
\end{attention}

\begin{popfigure}
\centering
\begin{tikzpicture}
\begin{axis}[
    ybar,
    bar width=0.6cm,
    width=\linewidth,
    height=8cm,
    ymin=0, ymax=22,
    ylabel={Points sur 20},
    symbolic x coords={Structure externe, Introduction, Corps (2 parties), Conclusion, Langue \& Style},
    xtick=data,
    x tick label style={rotate=30, anchor=east, font=\small},
    ymajorgrids=true,
    grid style=dashed,
    nodes near coords,
    nodes near coords align={vertical},
    every node near coord/.append style={font=\small\bfseries, text=popdark},
    title style={font=\bfseries, text=popblue},
    title={Barème indicatif de l'évaluation sommative (Total 20 pts)},
    fill=popblue!70,
]
\addplot coordinates {
    (Structure externe, 2)
    (Introduction, 4)
    (Corps (2 parties), 8)
    (Conclusion, 4)
    (Langue \& Style, 2)
};
\end{axis}
\end{tikzpicture}
\legende{Barème de pondération : Le corps du rapport (contenu, structure, exploitation docs) porte le poids majeur. La langue sanctionne la maîtrise du code.}
\end{popfigure}

\courssub{Auto-évaluation et grille de critères}

L'auto-évaluation développe l'autonomie et la métacognition. Elle se fait \textbf{après} la remise de la copie corrigée (ou en peer-review avant remise), à l'aide d'une grille critériée connue à l'avance.

\begin{aretenir}[Grille d'auto-évaluation — À compléter par l'élève (Oui / Partiellement / Non)]
\begin{center}
\begin{tabularx}{\linewidth}{|l|X|X|X|}\hline
\textbf{Critère} & \textbf{Oui (Maîtrisé)} & \textbf{Partiellement (À revoir)} & \textbf{Non (Non acquis)} \\ \hline
\textbf{Structure externe} : Page de garde complète (titre, auteur, destinataire, date, structure) & & & \\ \hline
\textbf{Introduction} : Accroche contextuelle + Objet précis (V+N) + Intérêt + Plan annoncé & & & \\ \hline
\textbf{Corps} : 2 parties équilibrées, titres parlants, arguments étayés par docs, connecteurs logiques & & & \\ \hline
\textbf{Conclusion} : Bilan synthétique + 3 Recommandations SMART (Verbe + Objet + Délai/Acteur) + Ouverture & & & \\ \hline
\textbf{Langue} : Impersonnalité (pas de « je »), Concordance des temps (PC/Imparfait/Présent/Futur), Vocabulaire technique & & & \\ \hline
\textbf{Présentation} : Aération, paragraphes, lisibilité, respect consignes (longueur) & & & \\ \hline
\end{tabularx}
\end{center}
\end{aretenir}

\begin{methode}[Utiliser la grille pour progresser]
\begin{enumerate}
    \item Cochez honnêtement chaque ligne.
    \item Pour chaque « Partiellement » ou « Non », identifiez \textbf{un exemple précis} dans votre copie (ex: « Mon objet p.2 est vague »).
    \item Fixez-vous \textbf{un objectif unique} pour le prochain écrit utilitaire (ex: « Systématiser les recommandations SMART »).
\end{enumerate}
\end{methode}

\begin{savaistu}[Le rapport : une compétence professionnelle transversale]
Au Cameroun, la rédaction de rapports est une \textbf{compétence cœur de métier} pour les techniciens supérieurs (BTS, HND, DUT) et les fonctionnaires.
\begin{itemize}
    \item \textbf{Stage / PFE} : Le rapport de stage est l'examen de fin de cycle (BTS, Licence Pro). Sa structure (intro, corps, conclusion, recommandations) est identique.
    \item \textbf{Entreprise} : Rapports d'activité mensuels, rapports d'incident, rapports d'audit, comptes rendus de réunion. L'objectivité, la clarté et l'orientation « décision » (recommandations) sont exigées.
    \item \textbf{Concours administratifs (ENAM, ENSET, Concours professionnels)} : L'épreuve de « Note de synthèse » ou « Rapport » est classique. Elle teste la capacité à traiter un dossier documentaire sous contrainte de temps.
    \item \textbf{Vie associative / Coopératives / GIC} : Tout responsable doit rendre des comptes (rapport moral, rapport financier) à l'assemblée générale. C'est une obligation légale (Loi 92/006 sur les coopératives, OHADA Acte uniforme société coopérative).
\end{itemize}
Maîtriser le rapport en Première Technique, c'est acquérir un \textbf{capital employabilité} durable.
\end{savaistu}

\newpage

\coursec{Activité d'intégration}

\coursec{Leçon 11 : Production d'écrits utilitaires — Rédiger un rapport (Introduction, Corps, Conclusion)}

\courssub{Objectifs de l'activité}
\begin{itemize}
    \item Mobiliser les connaissances sur la \cle{structure externe} (page de garde, sommaire, annexes) et la \cle{structure interne} (introduction, corps, conclusion) du rapport technique.
    \item Exploiter des documents hétérogènes (texte iconique, tableau chiffré, texte fonctionnel administratif, extrait satirique, modèle de rapport) pour extraire, sélectionner et organiser l'information pertinente.
    \item Appliquer la \cle{neutralité} du langage technique en transformant un extrait à tonalité \cle{satirique} et un extrait à tonalité \cle{tragique} en énoncés objectifs.
    \item Rédiger un rapport complet, structuré et argumenté, respectant les conventions de la communication professionnelle (objectivité, précision, impersonnalité).
    \item S'auto-évaluer et améliorer sa production grâce à une grille de correction critériée (confrontation des productions).
\end{itemize}

\courssub{Situation}
\begin{exemplebox}[Contexte professionnel simulé]
\textbf{Lycée Technique Industriel de Douala-Bassa (LTIDB)}, Octobre 2026.\\
À la suite d'un épisode pluviométrique exceptionnel (187 mm en 48 h, selon la station météo de Douala), les ateliers du pôle « Génie Mécanique et Électrique » ont subi d'importantes infiltrations et une inondation partielle. Le Proviseur, M. \textsc{Ebongué}, mandate une commission d'élèves de \textbf{Première Technique} (spécialités F1, F2, F3) pour réaliser une \textbf{inspection technique} et produire un \textbf{rapport d'expertise} destiné au Conseil d'Administration et à la Tutelle (MINESEC / MINESUP).
\end{exemplebox}

\courssub{Dossier documentaire fourni aux groupes (4 à 5 élèves)}

\textbf{Document 1 (Texte iconique) :} Photographie de l'atelier d'électrotechnique (prise le 12/10/2026 à 10h15). On y observe : niveau d'eau résiduel (~15 cm), armoires électriques entrouvertes, câbles trempés, outillage au sol, traces de boue sur les machines-outils (tour parallèle, fraiseuse).

\begin{popfigure}
\begin{tikzpicture}[scale=0.7]
% Simplified workshop sketch
\draw[popblue, thick, fill=popblue!5] (0,0) rectangle (10,6); % Room
\draw[popblue!50, fill=popblue!20] (0,0) rectangle (10,1.5); % Water
\draw[poporange, thick] (1,2) rectangle (3,4) node[midway, align=center] {Armoire\\Électrique};
\draw[poporange, thick] (4,2) rectangle (6,3.5) node[midway, align=center] {Tour\\Parallèle};
\draw[poporange, thick] (7,2) rectangle (9,3.5) node[midway, align=center] {Fraiseuse};
\draw[popmuted, decorate, decoration={snake, amplitude=2pt, segment length=10pt}] (0.5,1.5) -- (9.5,1.5); % Mud line
\foreach \x in {1.5, 2.5, 4.5, 5.5, 7.5, 8.5} {
    \fill[popmuted] (\x,1.2) circle (0.15); % Debris
}
\draw[->, popdark, thick] (10.5,3) -- (11.5,3) node[right, align=left] {Niveau eau\\~15 cm};
\draw[->, popdark, thick] (10.5,1) -- (11.5,1) node[right] {Boue / Débris};
\end{tikzpicture}
\legende{Document 1 : Croquis schématique de l'atelier d'électrotechnique après décrue (vue de coupe).}
\end{popfigure}

\textbf{Document 2 (Tableau chiffré) :} État des lieux quantitatif des équipements endommagés (estimation provisoire par le Chef de Travaux, M. \textsc{Ndongo}).
\begin{center}
\begin{tabularx}{\linewidth}{|l|X|c|c|c|}\hline
\rowcolor{popblueL}
\textbf{Atelier} & \textbf{Équipement / Matériel} & \textbf{Qté} & \textbf{Coût unitaire (FCFA)} & \textbf{Coût total (FCFA)} \\ \hline
Électrotechnique & Armoire de commande programmable (API Siemens S7-1200) & 2 & 1\,250\,000 & 2\,500\,000 \\ \hline
Électrotechnique & Multimètres numériques Fluke 117 & 5 & 45\,000 & 225\,000 \\ \hline
Électrotechnique & Câblage réseau (brassage baie, 120 ml) & 1 lot & 380\,000 & 380\,000 \\ \hline
Mécanique & Tour parallèle à commande numérique (Fanuc) - Carte mère & 1 & 1\,800\,000 & 1\,800\,000 \\ \hline
Mécanique & Jeu de mèches HSS / Outillage de coupe & 1 lot & 120\,000 & 120\,000 \\ \hline
Menuiserie & Scie circulaire à format (moteur triphasé) & 1 & 650\,000 & 650\,000 \\ \hline
Menuiserie & Stock de bois (Okoumé, Iroko - 15 m³) & 15 & 45\,000 & 675\,000 \\ \hline
Transversal & Système d'aspiration / ventilation (moteurs brûlés) & 3 & 210\,000 & 630\,000 \\ \hline
\rowcolor{popcream}
\multicolumn{4}{|r|}{\textbf{TOTAL ESTIMATIF DES DOMMAGES MATÉRIELS}} & \textbf{6\,980\,000} \\ \hline
\end{tabularx}
\end{center}

\textbf{Document 3 (Texte fonctionnel - Décision administrative) :} Extrait de la Décision N° 011/LTIDB/PROV/2026 du 13/10/2026.
\begin{exemplebox}[Décision de mission]
\textbf{RÉPUBLIQUE DU CAMEROUN} \hfill \textbf{Paix – Travail – Patrie} \\
\textbf{MINESEC / MINESUP} \\
\textbf{LYCÉE TECHNIQUE INDUSTRIEL DE DOUALA-BASSA} \\
\textbf{DÉCISION N° 011/LTIDB/PROV/2026 DU 13/10/2026} \\
\textbf{Portant mission d'inspection des ateliers post-inondation.}
Le Proviseur du Lycée Technique Industriel de Douala-Bassa,
Vu le règlement intérieur ; Vu le procès-verbal de constat d'huissier N° 45/2026 ;
Décide :
\textbf{Article 1 :} Est constituée une commission d'inspection technique composée d'élèves de Première Technique (F1, F2, F3) sous la supervision de M. \textsc{Ndongo}, Chef de Travaux.
\textbf{Article 2 :} La commission est chargée de : constater l'étendue des dégâts ; dresser l'inventaire des équipements endommagés ; estimer le coût de la remise en état ; formuler des recommandations techniques.
\textbf{Article 3 :} Le rapport final est attendu pour le 20 octobre 2026 au plus tard sur la table du Proviseur.
\textbf{Article 4 :} La présente décision prend effet à compter de sa signature.
\vspace{0.5em}
\hfill \textbf{Le Proviseur,} \hspace{2cm} M. \textsc{Ebongué}
\end{exemplebox}

\textbf{Document 4 (Extrait de presse - Tonalité satirique) :} « \emph{Au LTIDB, la piscine olympique est enfin ouverte ! Plus besoin d'aller à la piscine municipale, les ateliers de mécanique et d'électrotechnique offrent un bassin naturel gratuit, chauffé par les moteurs des tours qui ont pris l'eau. Nos chers élèves pourront désormais pratiquer la natation synchronisée entre deux cours de câblage. Le Proviseur, grand architecte de ce centre aquatique improvisé, a oublié de prévoir les maillots de bain dans la liste des fournitures scolaires. Quant aux armoires électriques, elles font désormais office d'aquariums géants pour poissons-chats. Une rentrée 2026... mouillée ! }» (Extrait du \emph{Canard Déchaîné}, 14/10/2026).

\textbf{Document 5 (Extrait associatif - Tonalité tragique) :} « \emph{Ô désastre ! Ô douleur indicible ! Les temples sacrés du savoir-faire, nos ateliers de mécanique et d'électrotechnique, gisent désormais sous les eaux souillées, tels des géants blessés à mort. Le rugissement des tours parallèles s'est tu, remplacé par le clapotis funeste de la boue. L'avenir de nos jeunes techniciens, forgé dans l'acier et le cuivre, est-il donc voué à pourrir dans l'indifférence générale ? Les armoires, cœurs battants de l'automatisme, sont devenues des tombeaux pour l'intelligence artificielle. Puisse le ciel entendre la plainte de la jeunesse camerounaise privée de son outil de travail ! }» (Extrait de la \emph{Gazette de l'Association des Anciens Élèves}, 15/10/2026).

\textbf{Document 6 (Modèle de rapport - Structure type MINESEC) :} Fournissant les rubriques obligatoires : Page de garde, Sommaire, Introduction (Contexte, Objet, Méthodologie, Plan), Corps (Constats, Analyse, Causes, Recommandations), Conclusion (Synthèse, Perspectives), Annexes (Photos, Devis, PV de réception).

\courssub{Consignes}
\begin{enumerate}
    \item \textbf{Analyse de la situation (15 min) :} Identifier la \cle{nature de la commande} (qui demande ? qui rédige ? pour qui ? quel délai ?), le \cle{type de rapport} (d'expertise / d'inspection / de sinistre) et l'\cle{objectif de communication} (informer, alerter, proposer des solutions chiffrées). Exploiter le Document 3 (texte fonctionnel) pour valider le cadre officiel.
    \item \textbf{Traitement des documents (30 min) :}
    \begin{enumerate}
        \item Décrire le Document 1 en langage technique (vocabulaire précis : \emph{infiltration, corrosion, court-circuit, immobilisation de l'outil de production}).
        \item Exploiter le Document 2 : Calculer le pourcentage de dommages par atelier et le coût global. Identifier les postes les plus critiques (coût \textgreater{} 1\,000\,000 FCFA).
        \item \textbf{Transposition de registre (Satirique) :} Réécrire le Document 4 en \cle{langage neutre et objectif} (style administratif/technique), en extrayant les 3 faits réels sous-jacents à la satire.
        \item \textbf{Transposition de registre (Tragique) :} Réécrire le Document 5 en \cle{langage neutre et objectif}, en extrayant les 3 faits réels sous-jacents à la tragédie.
    \end{enumerate}
    \item \textbf{Élaboration du plan détaillé (15 min) :} Construire le sommaire (structure interne) du rapport en 3 parties principales (Constats, Analyse/Chiffrage, Recommandations) avec au moins 2 sous-parties chacune. Justifier l'ordre logique choisi (chronologique ? thématique ? gravité ?).
    \item \textbf{Rédaction collaborative (45 min) :} Rédiger l'\textbf{Introduction} (contexte, objet, méthode, plan annoncé), le \textbf{Corps} (intégrer description iconique, tableau analysé, causes techniques probables, 3 recommandations hiérarchisées : urgence, court terme, moyen terme) et la \textbf{Conclusion} (synthèse chiffrée, impact pédagogique, appel à décision).
    \item \textbf{Confrontation et amélioration (20 min) :} Échanger les productions entre groupes. Utiliser la \textbf{Grille d'évaluation} (voir Ressources) pour noter : Structure (5 pts), Contenu/Exploitation docs (5 pts), Langage/Style (5 pts), Présentation (5 pts). Rédiger une \textbf{version finale améliorée} intégrant les corrections.
\end{enumerate}

\courssub{Ressources à mobiliser}
\begin{definition}[Structure externe du rapport (MINESEC)]
\begin{itemize}
    \item \textbf{Page de garde :} Identité de l'établissement, titre du rapport, auteur(s), destinataire, date, référence.
    \item \textbf{Sommaire :} Liste des parties/sous-parties avec pagination.
    \item \textbf{Annexes :} Documents bruts (photos, tableaux, devis, PV) référencés dans le corps (ex: \emph{voir Annexe 1}).
\end{itemize}
\end{definition}

\begin{definition}[Structure interne type (Introduction - Corps - Conclusion)]
\begin{itemize}
    \item \textbf{Introduction (I) :} \textbf{Contexte} (circonstances), \textbf{Objet} (mission), \textbf{Méthodologie} (visite, entretiens, docs consultés), \textbf{Plan annoncé} (transition).
    \item \textbf{Corps (II, III, IV...) :} Parties logiques. \textbf{II. Constats descriptifs} (ce qu'on a vu/mesuré). \textbf{III. Analyse technique et financière} (interprétation, chiffrage, causes). \textbf{IV. Recommandations} (actions correctives/préventives, priorisées).
    \item \textbf{Conclusion :} \textbf{Bilan synthétique} (réponse à l'objet), \textbf{Impact/Enjeux} (pédagogique, sécurité, budgétaire), \textbf{Perspective/Suite à donner} (validation, marché, suivi).
\end{itemize}
\end{definition}

\begin{propriete}[Tonalités : Satirique, Tragique vs Neutre/Technique]
\begin{itemize}
    \item \textbf{Satirique :} Ironie, hyperbole, subjectivité, vocabulaire connoté (\emph{piscine, aquarium, gratuit}), but divertissant/critique.
    \item \textbf{Tragique :} Pathos, emphase, vocabulaire lyrique (\emph{temples, géants blessés, tombeaux, plainte}), figures de style (métaphore, personnification), but émouvoir/accuser.
    \item \textbf{Neutre/Technique (exigible au Probatoire) :} Objectivité, impersonnalité (passif, infinitif, « il convient de »), vocabulaire précis (\emph{inondation, immersion, dysfonctionnement, coût de remplacement}), phrases déclaratives, absence de jugement de valeur.
\end{itemize}
\end{propriete}

\begin{methode}[De la situation au plan de rapport]
\begin{enumerate}
    \item Lire la commande et les documents ; surligner les \cle{mots-clés} (chiffres, termes techniques, verbes d'action).
    \item Classer les informations par \cle{thèmes} (Bâtiment, Équipements, Pédagogie, Sécurité, Budget).
    \item Hiérarchiser : \textbf{Constats} (faits bruts) $\rightarrow$ \textbf{Analyse} (liens causes/effets, coûts) $\rightarrow$ \textbf{Action} (recommandations SMART : Spécifiques, Mesurables, Atteignables, Réalistes, Temporelles).
    \item Formuler les titres de parties/sous-parties (noms infinitifs ou nominaux, parallélisme syntaxique).
\end{enumerate}
\end{methode}

\begin{aretenir}[Grille d'évaluation simplifiée (20 pts)]
\begin{center}
\begin{tabularx}{\linewidth}{|X|c|X|}\hline
\rowcolor{popblueL}
\textbf{Critère} & \textbf{Pts} & \textbf{Indicateurs de réussite} \\ \hline
Structure externe/interne & 5 & Page de garde complète, sommaire cohérent, I-C-C respecté, annexes citées. \\ \hline
Exploitation documents & 5 & Doc 1 décrit techniquement, Doc 2 chiffré/analysé (\%), Doc 3 (décision) cité, Doc 4 \& 5 transposés (3 faits chacun), Doc 6 suivi. \\ \hline
Langage / Style & 5 & Neutre, impersonnel, vocabulaire technique précis, connecteurs logiques, pas de familier/satirique/tragique. \\ \hline
Qualité rédactionnelle & 5 & Orthographe, syntaxe, ponctuation, lisibilité, mise en page (paragraphes, alinéas). \\ \hline
\end{tabularx}
\end{center}
\end{aretenir}

\courssub{Démarche et corrigé commenté}
\begin{exoresolu}[Rapport d'inspection des ateliers post-inondation - Version modèle corrigée]
\textbf{Énoncé :} Produire le rapport final attendu du groupe « excellence » après confrontation et remédiation.
\tcblower
\textbf{RAPPORT D'INSPECTION TECHNIQUE DES ATELIERS SUITE AUX INONDATIONS D'OCTOBRE 2026}

\textbf{Page de garde}\\
\begin{center}
\begin{tabular}{|c|}\hline
\rowcolor{popblueL} \textbf{RÉPUBLIQUE DU CAMEROUN} \\ \hline
\rowcolor{popblueL} \textbf{Paix – Travail – Patrie} \\ \hline
\rowcolor{popblueL} \textbf{MINESEC / MINESUP} \\ \hline
\rowcolor{popblueL} \textbf{LYCÉE TECHNIQUE INDUSTRIEL DE DOUALA-BASSA} \\ \hline
\vspace{1cm}
\textbf{RAPPORT D'INSPECTION TECHNIQUE} \\ \vspace{0.2cm}
\textbf{ÉTAT DES LIEUX DES ATELIERS APRÈS LES INONDATIONS} \\ \vspace{0.2cm}
\textbf{DU 10 AU 12 OCTOBRE 2026} \\ \vspace{1.5cm}
\textbf{Objet :} Constat des dommages, estimation financière et propositions de réhabilitation. \\ \vspace{1.5cm}
\textbf{Rédigé par :} Commission d'élèves de Première Technique (Spéc. F1, F2, F3) \\ \vspace{0.5cm}
\textbf{Destinataire :} M. le Proviseur / Président du Conseil d'Administration \\ \vspace{0.5cm}
\textbf{Date de remise :} 20 Octobre 2026 \\ \vspace{0.5cm}
\textbf{Référence :} N° 014/LTIDB/PROV/2026 \\ \hline
\end{tabular}
\end{center}

\textbf{Sommaire}
\begin{enumerate}
    \item \textbf{INTRODUCTION} .................................................................................. 3
    \item \textbf{CONSTATS DES DOMMAGES MATÉRIELS ET ENVIRONNEMENTAUX} .............. 4
    \item \hspace{0.5cm} 1.1. État des infrastructures et des locaux .............................................. 4
    \item \hspace{0.5cm} 1.2. État des équipements et outillages par atelier ................................. 4
    \item \textbf{ANALYSE TECHNIQUE ET FINANCIÈRE DES PRÉJUDICES} ...................... 5
    \item \hspace{0.5cm} 2.1. Répartition analytique des coûts de remise en état ......................... 5
    \item \hspace{0.5cm} 2.2. Identification des causes techniques aggravantes .......................... 6
    \item \textbf{RECOMMANDATIONS ET PLAN D'ACTION PRIORISÉ} .............................. 6
    \item \hspace{0.5cm} 3.1. Mesures d'urgence (J+1 à J+7) : Sécurisation et sauvetage .............. 6
    \item \hspace{0.5cm} 3.2. Mesures à court terme (Mois 1-2) : Réparation et remplacement ....... 7
    \item \hspace{0.5cm} 3.3. Mesures à moyen terme (Année 2027) : Prévention et résilience ....... 7
    \item \textbf{CONCLUSION} .................................................................................... 8
    \item \textbf{ANNEXES} ......................................................................................... 9
\end{enumerate}

\courssub{Introduction}
\textbf{Contexte.} Du 10 au 12 octobre 2026, la ville de Douala a enregistré des précipitations cumulées de 187 mm en 48 heures (source : Météo Cameroun), provoquant la saturation du réseau d'assainissement du quartier Bassa. Le Lycée Technique Industriel de Douala-Bassa (LTIDB) a subi une montée des eaux pluviales ayant envahi les ateliers du pôle « Génie Mécanique et Électrique » (Électrotechnique, Mécanique, Menuiserie) ainsi que les locaux techniques attenants.

\textbf{Objet.} Par décision N° 011/LTIDB/PROV/2026 en date du 13/10/2026 (cf. \textbf{Annexe 2}), M. le Proviseur a mandaté une commission d'élèves de Première Technique pour procéder à l'inspection visuelle et technique desdits ateliers, évaluer l'étendue des dégâts matériels, en estimer le coût financier et formuler des recommandations techniques en vue de la reprise rapide des activités pédagogiques.

\textbf{Méthodologie.} La commission a procédé, les 14 et 15 octobre 2026, à : (1) une visite in situ des trois ateliers avec relevés photographiques (cf. \textbf{Annexe 1 / Doc 1}) ; (2) un inventaire contradictoire des équipements endommagés en présence des Chefs de Travaux (cf. \textbf{Annexe 3 / Doc 2}) ; (3) l'exploitation du registre de maintenance et des factures d'acquisition pour l'estimation des coûts de remplacement ; (4) la consultation du modèle de rapport administratif en vigueur (cf. \textbf{Annexe 5 / Doc 6}) ; (5) l'analyse des réactions externes (presse, association) pour mesurer l'impact médiatique (cf. \textbf{Annexe 4 / Doc 4 \& 5}).

\textbf{Plan annoncé.} Ce rapport présente successivement les constats descriptifs (I), l'analyse technique et financière des préjudices (II), puis les recommandations structurées par horizon temporel (III), avant de conclure sur les enjeux de la reprise pédagogique.

\courssub{I. Constats des dommages matériels et environnementaux}
\courssub{État des infrastructures et des locaux}
L'inspection visuelle (cf. \textbf{Annexe 1 : Photographies}) révèle une hauteur d'eau résiduelle moyenne de 15 cm au sol de l'atelier d'électrotechnique au moment de la visite (J+4). Les sols en béton lissé présentent des traces de boue et d'hydrocarbures (fuites moteurs). Les murs bas (h < 1 m) montrent des traces d'humidité capillaire. Le système d'évacuation des eaux pluviales (caniveaux périphériques) est obstrué par des déchets végétaux et des sédiments, empêchant l'écoulement gravitaire. L'installation électrique générale (tableau principal, gaines techniques) a été immergée sur une hauteur de 30 cm, imposant une vérification diélectrique complète avant toute ré-energisation.

\courssub{État des équipements et outillages par atelier}
\begin{itemize}
    \item \textbf{Atelier Électrotechnique :} Deux armoires de commande API (Siemens S7-1200) présentent des traces d'oxydation sur les borniers d'entrée/sortie et les cartes processeurs. Cinq multimètres Fluke 117 affichent des dysfonctionnements d'affichage (condensation interne). Le brassage de la baie réseau (120 ml de câble RJ45 Cat 6) est inutilisable (gainage gorgé d'eau).
    \item \textbf{Atelier Mécanique :} Le tour parallèle à commande numérique (Fanuc) ne démarre pas ; diagnostic préliminaire : carte mère HS (court-circuit par humidité). L'outillage de coupe (mèches HSS, plaquettes carbure) stocké en tiroirs bas est oxydé (rouille de surface).
    \item \textbf{Atelier Menuiserie :} La scie circulaire à format (moteur triphasé 4 kW) a subi une immersion du bobinage. Le stock de bois (Okoumé, Iroko - 15 m³) a absorbé l'eau (taux d'humidité > 30\%), risquant déformation et développement fongique (moisissures).
    \item \textbf{Transversal :} Les trois moteurs du système d'aspiration centralisée des copeaux/poussières sont hors service (roulements grippés, isolation défaillante).
\end{itemize}

\courssub{II. Analyse technique et financière des préjudices}
\courssub{Répartition analytique des coûts de remise en état}
L'exploitation du Document 2 (inventaire chiffré, \textbf{Annexe 3}) permet d'établir la ventilation budgétaire suivante. Le coût total estimatif des dommages matériels s'élève à \textbf{6\,980\,000 FCFA}.

\begin{center}
\begin{tabular}{|l|c|c|c|}\hline
\rowcolor{popblueL}
\textbf{Atelier / Poste} & \textbf{Coût total (FCFA)} & \textbf{\% du total} & \textbf{Niveau de criticité} \\ \hline
Électrotechnique (API, Multimètres, Câblage) & 3\,105\,000 & 44,5 \% & \textbf{Critique} (Arrêt TP) \\ \hline
Mécanique (Tour CN, Outillage) & 1\,920\,000 & 27,5 \% & \textbf{Critique} (Arrêt TP) \\ \hline
Menuiserie (Scie, Bois) & 1\,325\,000 & 19,0 \% & \textbf{Élevé} (Stock perdu) \\ \hline
Transversal (Aspiration) & 630\,000 & 9,0 \% & \textbf{Moyen} (Sécurité/Hygiène) \\ \hline
\textbf{TOTAL} & \textbf{6\,980\,000} & \textbf{100 \%} &  \\ \hline
\end{tabular}
\end{center}

\textbf{Analyse chiffrée :} Les ateliers Électrotechnique et Mécanique concentrent \textbf{72 \%} du coût global. La perte de l'API et du Tour CN (coût unitaire > 1,2 M FCFA) constitue un blocage majeur pour les épreuves pratiques du Probatoire. Le stock de bois, bien que moins coûteux unitairement, représente une perte de matière première non récupérable à 100\%.

\courssub{Identification des causes techniques aggravantes}
Au-delà de l'aléa climatique (pluie exceptionnelle), trois facteurs techniques ont amplifié les dégâts :
\begin{enumerate}
    \item \textbf{Implantation basse des équipements sensibles :} Armoires API et variateurs de vitesse posés au niveau du sol (zéro surélévation), contrairement aux normes NF C 15-100 (hauteur minimale 30 cm au-dessus du niveau fini de sol).
    \item \textbf{Défaillance du système d'évacuation :} Caniveaux non curés depuis > 2 ans (absence de plan de maintenance préventive du bâti).
    \item \textbf{Absence de procédure d'urgence :} Aucune consigne de mise en sécurité rapide (coupure générale, surélévation outillage) n'était affichée ni connue des élèves/encadrants.
\end{enumerate}

\courssub{III. Recommandations et plan d'action priorisé}
\courssub{Mesures d'urgence (J+1 à J+7) : Sécurisation et sauvetage}
\begin{enumerate}
    \item \textbf{Coupure et consignation électrique générale} des ateliers concernés (vérification absence de tension) avant toute intervention. Affichage « Interdit - Travaux ».
    \item \textbf{Séchage et décontamination} des armoires API et du Tour CN par air chaud sec (T < 50°C) et alcool isopropylique pour les circuits imprimés, réalisé par le Chef de Travaux Électrotechnique assisté d'un prestataire agréé Siemens/Fanuc.
    \item \textbf{Évacuation et séchage du stock de bois} sous abri ventilé (bâche + ventilation forcée) pour limiter le pourrissement ; tri des pièces récupérables (objectif : sauver 30\% du volume).
    \item \textbf{Nettoyage complet des locaux} (décapage boue, désinfection sols) par entreprise de nettoyage industriel.
\end{enumerate}

\courssub{Mesures à court terme (Mois 1-2) : Réparation et remplacement}
\begin{enumerate}
    \item \textbf{Lancement de consultations (3 devis minimum)} pour le remplacement des équipements « Critiques » : 2 API S7-1200, 1 Carte mère Tour Fanuc, 1 Moteur triphasé Scie, 3 Moteurs aspiration. Budget prévisionnel : 5\,500\,000 FCFA (hors main d'œuvre).
    \item \textbf{Réparation/étalonnage} des multimètres et outillage de coupe (affûtage, traitement antirouille) en interne (équipe maintenance).
    \item \textbf{Remise en conformité électrique} : Surélévation des armoires sur plots béton (h = 40 cm), installation de presses-étanches IP67, pose de seuils de porte étanches (h = 15 cm) aux entrées d'ateliers.
\end{enumerate}

\courssub{Mesures à moyen terme (Année 2027) : Prévention et résilience}
\begin{enumerate}
    \item \textbf{Étude hydraulique du site} par un bureau d'études (BET) pour redimensionner le réseau d'assainissement (caniveaux, regards, pompe de relevage si nécessaire). Budget estimé : 3\,000\,000 FCFA.
    \item \textbf{Élaboration d'un Plan de Continuité d'Activité (PCA) ateliers} : Fiches de poste « Crue », exercices de simulation annuels, kit d'urgence (bâches, pompes vide-cave, sangles) stocké en salle des professeurs.
    \item \textbf{Intégration au budget d'investissement 2027} (ligne « Rénovation ateliers ») pour pérenniser les travaux de mise hors d'eau.
\end{enumerate}

\courssub{Conclusion}
La crue d'octobre 2026 a exposé la vulnérabilité structurelle des ateliers du LTIDB face aux aléas climatiques de plus en plus fréquents. Le bilan matériel s'élève à \textbf{6\,980\,000 FCFA}, dont 72 \% concentrés sur les équipements numériques de pointe (API, Tour CN) indispensables aux épreuves certificatives. L'analyse technique met en évidence des non-conformités préexistantes (implantation au sol, drainage défaillant) qui ont transformé un aléa météorologique en sinistre pédagogique majeur.

La mise en œuvre immédiate des mesures d'urgence (séchage, sécurisation) et le lancement rapide des marchés de remplacement (court terme) sont la condition \emph{sine qua non} de la tenue des TP du second trimestre. Parallèlement, l'investissement dans la résilience infrastructurelle (drainage, surélévation, PCA) s'impose comme une priorité stratégique pour l'établissement.

La commission recommande au Conseil d'Administration de valider le plan d'action ci-dessus et de solliciter auprès de la Tutelle (MINESEC) une dotation exceptionnelle « Fonds de secours climatique » pour couvrir le gap budgétaire non absorbable par le budget de fonctionnement courant.

\courssub{Annexes}
\begin{itemize}
    \item \textbf{Annexe 1 :} Photographies des ateliers (Électrotechnique, Mécanique, Menuiserie) - 6 clichés datés/horodatés (Document 1).
    \item \textbf{Annexe 2 :} Décision de mission N° 011/LTIDB/PROV/2026 (Document 3).
    \item \textbf{Annexe 3 :} Inventaire détaillé des dommages (Document 2 complet + signatures Chefs de Travaux).
    \item \textbf{Annexe 4 :} Extraits de presse et associatifs (Documents 4 \& 5) et leurs transpositions en langage neutre (voir ci-dessous).
    \item \textbf{Annexe 5 :} Modèle de rapport utilisé (Document 6).
    \item \textbf{Annexe 6 :} Relevés météorologiques Douala (10-12/10/2026).
\end{itemize}

\courssub{Annexe 4 : Transposition des Documents 4 et 5 (Satirique / Tragique $\rightarrow$ Neutre)}
\begin{center}
\begin{tabularx}{\linewidth}{|X|X|}\hline
\rowcolor{popblueL}
\textbf{Extrait satirique (Doc 4)} & \textbf{Réécriture neutre / technique (Rapport)} \\ \hline
« La piscine olympique est enfin ouverte ! » & \textbf{Fait 1 :} Inondation des ateliers (Électrotechnique, Mécanique) par eaux pluviales, hauteur ~15 cm résiduelle. \\ \hline
« Armoires électriques font office d'aquariums » & \textbf{Fait 2 :} Immersion partielle des armoires de commande (API, variateurs) entraînant risque de court-circuit et corrosion. \\ \hline
« Rentrée 2026... mouillée ! » & \textbf{Fait 3 :} Perturbation majeure des activités pratiques (TP) ; impossibilité d'utiliser les machines-outils numériques et conventionnelles. \\ \hline
\end{tabularx}
\end{center}
\vspace{0.5em}
\begin{center}
\begin{tabularx}{\linewidth}{|X|X|}\hline
\rowcolor{poppurpleL}
\textbf{Extrait tragique (Doc 5)} & \textbf{Réécriture neutre / technique (Rapport)} \\ \hline
« Temples sacrés du savoir-faire... gisent sous les eaux » & \textbf{Fait 1 :} Inondation des locaux techniques (ateliers Mécanique, Électrotechnique, Menuiserie) par eaux pluviales. \\ \hline
« Rugissement des tours s'est tu... cœurs battants devenus tombeaux » & \textbf{Fait 2 :} Arrêt total des machines-outils numériques (Tour CN Fanuc, API Siemens) et conventionnelles (Scie, Tour parallèle) suite à immersion. \\ \hline
« Avenir de nos jeunes techniciens... privé de son outil de travail » & \textbf{Fait 3 :} Compromission de la formation pratique (TP, épreuves certificatives Probatoire) par indisponibilité de l'outil de production. \\ \hline
\end{tabularx}
\end{center}
\vspace{0.5em}
\textit{Commentaire :} Les deux transpositions ont éliminé les marques subjectives (ironie, hyperbole, pathos, métaphores) pour ne conserver que l'information factuelle, mesurable et exploitable techniquement.

\textbf{Fin du rapport.}
\end{exoresolu}

\courssub{Bilan de l'activité}
\begin{aretenir}[Compétences validées par l'activité d'intégration]
Cette activité a permis de mettre en évidence, dans une situation professionnelle authentique, la maîtrise des savoir-faire suivants :
\begin{enumerate}
    \item \textbf{Analyse de la commande :} Identification correcte du destinataire, de l'expéditeur, de l'objet et du type de rapport (expertise/sinistre) via l'exploitation du texte fonctionnel (Décision administrative).
    \item \textbf{Traitement documentaire croisé :} Capacité à passer de l'image (iconique) au texte descriptif technique, du tableau brut (chiffres) à l'analyse relative (pourcentages, criticité), et de la satire/tragédie (presse, associatif) à l'objectivité administrative (double transposition de registre).
    \item \textbf{Structuration logique :} Respect de l'architecture standardisée (Externe : page de garde, sommaire, annexes ; Interne : Introduction \emph{Contextualisée/Objectivée/Planifiée}, Corps \emph{Descriptif/Analytique/Prospectif}, Conclusion \emph{Synthétique/Engageante}).
    \item \textbf{Qualité langagière :} Usage constant de l'impersonnalité (passif, infinitif, « il convient de »), du vocabulaire spécialisé (\emph{immersion, oxydation, diélectrique, surélévation, PCA, SMART}), des connecteurs logiques (\emph{par ailleurs, en conséquence, par conséquent, eu égard à}).
    \item \textbf{Démarche qualité :} La phase de confrontation (grille 20 pts) a révélé les écarts types (oubli annexes, plan non annoncé, langage familier persistant, chiffrage non exploité, non-exploitation du texte fonctionnel) et permis une remédiation immédiate par réécriture collaborative.
\end{enumerate}
L'élève de Première Technique est désormais apte à produire un rapport technique de niveau probatoire, document de décision à part entière, transférable en stage en entreprise ou en milieu professionnel.
\end{aretenir}

\newpage

\coursec{Méthodes et exercices résolus}

\coursec{Production d'écrits utilitaires : rédiger un rapport}

\courssub{Analyser la situation de communication et identifier le type de rapport}

\begin{methode}[Analyser la situation d'énonciation avant d'écrire]
    \textbf{Objectif :} Déterminer les contraintes du rapport (type, destinataire, objectif, ton) pour adapter la structure et le lexique.
    \begin{enumerate}
        \item \textbf{Identifier la commande :} Repérer le verbe d'action (rédiger, présenter, analyser, proposer) et le sujet imposé.
        \item \textbf{Définir le type de rapport :}
            \begin{itemize}
                \item \emph{Rapport de stage / d'activité :} relato factuel, chronologique, bilan compétences.
                \item \emph{Rapport d'incident / d'accident :} circonstantiel, causes, conséquences, mesures correctives.
                \item \emph{Rapport de mission / de visite :} objectifs, déroulement, observations, recommandations.
                \item \emph{Rapport d'enquête / d'étude :} problématique, méthodologie, résultats, analyse, conclusions.
            \end{itemize}
        \item \textbf{Préciser le destinataire :} Hiérarchie (formel, respectueux), client (pédagogique, persuasif), jury (académique, rigoureux). Adaptez le niveau de langue (soutenu technique).
        \item \textbf{Fixer l'objectif :} Informer ? Convaincre ? Décider ? Alerter ? Cela oriente la conclusion (recommandations vs simple bilan).
        \item \textbf{Lister les sources :} Notes de terrain, documents fournis (textes fonctionnels : procès-verbaux, statistiques, courriers ; textes iconiques : photos, schémas, organigrammes).
    \end{enumerate}
    \textbf{Réflexe Bac :} Avant de plancher, remplissez une \emph{fiche d'analyse} : Type = ... | Destinataire = ... | Objectif = ... | Ton = Neutre/Objectif (sauf parties analytiques).
\end{methode}

\begin{exoresolu}[Analyse de commande : Rapport de stage en entreprise]
    \textbf{Énoncé :} Vous avez effectué un stage de trois semaines dans l'entreprise \emph{« Cameroun Agro-Industrie »} (CAI), service Maintenance. Le Directeur des Ressources Humaines vous demande un rapport de fin de stage pour évaluer votre insertion et la pertinence des tâches confiées par rapport à votre formation en Maintenance Industrielle (Première Technique). Rédigez l'introduction de ce rapport.
    \tcblower
    \textbf{Solution détaillée :}
    \begin{enumerate}
        \item \textbf{Analyse de la situation (Fiche d'analyse) :}
            \begin{itemize}
                \item \textbf{Type :} Rapport de stage / d'activité (formation technique).
                \item \textbf{Rédacteur :} Élève de 1ère Tech, spécialité Maintenance Industrielle (statut : stagiaire).
                \item \textbf{Destinataire :} DRH de CAI (interne à l'entreprise) + Tuteur de stage + Jury scolaire (double destinataire $\rightarrow$ ton formel, technique mais accessible).
                \item \textbf{Objet :} Évaluation de l'insertion et adéquation stage/formation.
                \item \textbf{Durée :} 3 semaines (chronologie possible).
                \item \textbf{Sources :} Fiche de poste, cahier de stage, entretiens avec tuteur, organigramme service maintenance (texte iconique), planning des interventions (texte fonctionnel).
            \end{itemize}
        \item \textbf{Rédaction de l'introduction (Modèle standard) :}
            \begin{quote}
                \textbf{Objet :} Rapport de stage de fin de formation – Maintenance Industrielle – Du 03/06/2025 au 21/06/2025.
                
                \textbf{Contexte (Accroche) :} Dans le cadre de la validation du Probatoire en Maintenance Industrielle, j'ai effectué un stage d'immersion professionnelle au sein de l'entreprise Cameroun Agro-Industrie (CAI), du 3 au 21 juin 2025.
                
                \textbf{Sujet :} Ce rapport rend compte de mon affectation au service Maintenance, sous la tutelle de M. \textsc{Ngono}, Chef de service. Il porte sur l'analyse des interventions préventives et curatives réalisées sur la ligne de conditionnement.
                
                \textbf{Problématique / Objectif :} Dans quelle mesure les activités de maintenance réalisées en entreprise permettent-elles de consolider les acquis théoriques de la classe de Première Technique et de développer les compétences professionnelles attendues ?
                
                \textbf{Annonce du plan :} Nous présenterons d'abord la structure de l'entreprise et le service d'affectation (I), puis nous détaillerons les activités techniques menées en les confrontant au référentiel de formation (II), pour finir par un bilan critique de l'apport du stage et des recommandations (III).
            \end{quote}
        \item \textbf{Justification des choix :} Le ton est \emph{objectif} et \emph{professionnel}. L'introduction intègre les 4 composantes obligatoires : Contexte, Sujet, Problématique/Objectif, Plan. Le vocabulaire technique (préventif, curatif, ligne de conditionnement, référentiel) est présent.
    \end{enumerate}
    \textbf{Conclusion :} L'analyse préalable garantit une introduction ciblée, évitant le hors-sujet (ex: raconter sa vie au lieu de l'adéquation stage/formation).
\end{exoresolu}

\courssub{Construire la structure externe et le plan détaillé (structure interne)}

\begin{methode}[Maquetter le rapport : Structure externe et Plan interne]
    \textbf{Objectif :} Produire un document lisible, normé et navigable. La structure externe est la « coquille », le plan interne est le « squelette logique ».
    \begin{enumerate}
        \item \textbf{Structure externe (Normes administratives/techniques) :}
            \begin{enumerate}
                \item \textbf{Page de garde :} Logo établissement/entreprise, Titre complet (Objet), Nom/N° candidat, Spécialité, Année scolaire, Lieu/Date.
                \item \textbf{Page de remerciements :} Ordre protocolaire (Dieu, Parents, Administration, Tuteur, Personnel).
                \item \textbf{Sommaire / Table des matières :} Numérotation décimale (1, 1.1, 1.1.1), alignement points de suite, numéros de page.
                \item \textbf{Liste des abréviations / sigles / tableaux / figures :} Si $>$ 5 occurrences.
                \item \textbf{Corps du rapport :} Introduction, Développement (Parties/Chapitres), Conclusion, Recommandations.
                \item \textbf{Annexes :} Documents bruts (PV, photos, schémas, fiches techniques) référencés dans le corps (ex: \emph{voir Annexe 3}).
            \end{enumerate}
        \item \textbf{Structure interne (Plan dialectique ou thématique) :}
            \begin{itemize}
                \item \textbf{Plan chronologique :} Idéal pour rapport de stage/journal de bord (Semaine 1, Semaine 2...). \emph{Risque :} effet « catalogue ».
                \item \textbf{Plan thématique / analytique (Recommandé Bac) :} Regroupe par \emph{problèmes} ou \emph{compétences}.
                    \begin{itemize}
                        \item I. Contexte \& Cadre (Entreprise, Service, Moyens).
                        \item II. Activités \& Analyse technique (Tâches, Problèmes résolus, Outils/Logiciels).
                        \item III. Bilan \& Perspectives (Compétences acquises, Difficultés, Recommandations).
                    \end{itemize}
                \item \textbf{Plan « Problème-Solution » :} Pour rapport d'incident/enquête (Constat $\rightarrow$ Causes $\rightarrow$ Solutions).
            \end{itemize}
        \item \textbf{Règle d'or :} Équilibre des parties ($\approx$ même volume), parallélisme des titres (même structure syntaxique : ex: tous les titres en groupes nominaux ou tous à l'infinitif).
    \end{enumerate}
    \textbf{Formule à retenir :} \cle{Titre Partie = Thème + Angle d'attaque}. Ex: \emph{« La maintenance préventive : organisation et exécution »} (Thème = Maintenance préventive ; Angle = Orga/Exécution).
\end{methode}

\begin{popfigure}
\begin{tikzpicture}[
    node distance=1.2cm and 1.5cm,
    box/.style={rectangle, draw=popblue, thick, fill=popblueL, rounded corners, minimum width=2.5cm, minimum height=1cm, align=center, font=\small},
    arrow/.style={-Stealth, thick, popdark}
]
\node[box] (garde) {Page de garde};
\node[box, below=of garde] (remerc) {Remerciements};
\node[box, below=of remerc] (som) {Sommaire / Listes};
\node[box, below=of som] (corps) {Corps du rapport};
\node[box, below=of corps] (ann) {Annexes};

\node[box, right=of corps, xshift=2cm] (intro) {Introduction};
\node[box, below=of intro] (dev) {Développement (I, II, III)};
\node[box, below=of dev] (concl) {Conclusion};
\node[box, below=of concl] (reco) {Recommandations};

\draw[arrow] (garde) -- (remerc);
\draw[arrow] (remerc) -- (som);
\draw[arrow] (som) -- (corps);
\draw[arrow] (corps) -- (ann);

\draw[arrow] (corps.east) -- ++(0.5,0) |- (intro.west);
\draw[arrow] (intro) -- (dev);
\draw[arrow] (dev) -- (concl);
\draw[arrow] (concl) -- (reco);
\end{tikzpicture}
\legende{Structure externe type d'un rapport technique (norme administrative).}
\end{popfigure}

\begin{exoresolu}[Élaboration d'un plan thématique équilibré]
    \textbf{Énoncé :} À partir de la situation du stage chez CAI (Maintenance Industrielle), l'élève a listé ses activités : 1) Visite usine + consignes sécurité. 2) Lecture plans hydrauliques/pneumatiques. 3) Maintenance préventive hebdo (graissage, vérif pressions). 4) Panne presse hydraulique (recherche fuite, changement joint). 5) Maintenance curative convoyeur (changement roulements). 6) Réunion hebdo + rédaction fiches intervention. 7) Formation logiciel GMAO (Gestion Maintenance Assistée par Ordinateur). 8) Bilan avec tuteur.
    \textbf{Travail demandé :} Proposer un plan thématique détaillé (3 parties, 2-3 sous-parties chacune) avec titres parallèles. Justifier le rejet du plan chronologique.
    \tcblower
    \textbf{Solution détaillée :}
    \begin{enumerate}
        \item \textbf{Critique du plan chronologique :} « Semaine 1 : Visite, Sécurité, Plans. Semaine 2 : Préventive, Panne presse. Semaine 3 : Curative convoyeur, GMAO, Bilan. »
        \textbf{Défaut :} Mélange les genres (formation + production), disperse l'analyse technique (préventif/curatif séparés), rend le bilan invisible. \textbf{Rejeté.}
        \item \textbf{Construction du plan thématique (Analytique) :}
        \textbf{Regroupement par compétences/objets techniques :}
        \begin{itemize}
            \item \textbf{Groupe A (Cadre \& Outils) :} Activités 1, 2, 6 (Réunions), 7 (GMAO). $\rightarrow$ \emph{Intégration et environnement technique.}
            \item \textbf{Groupe B (Maintenance Préventive) :} Activité 3. $\rightarrow$ \emph{Planification et exécution.}
            \item \textbf{Groupe C (Maintenance Curative / Dépannage) :} Activités 4, 5. $\rightarrow$ \emph{Diagnostic et remise en état.}
            \item \textbf{Groupe D (Bilan) :} Activité 8. $\rightarrow$ \emph{Évaluation \& Recommandations.}
        \end{itemize}
        \item \textbf{Plan détaillé retenu (3 Parties équilibrées) :}
        \begin{center}
            \begin{tabularx}{\linewidth}{|c|X|}\hline
            \rowcolor{popblueL} \textbf{Partie} & \textbf{Titres (Parallélisme : GN + Double infinitif/GN)} \\\hline
            \textbf{I} & \textbf{Intégration au service maintenance : découverte, sécurité et outils numériques} \\
            \hline
            I.1 & Découverte de l'environnement industriel et des consignes de sécurité \\\hline
            I.2 & Lecture de plans et appropriation de la GMAO \\\hline
            \textbf{II} & \textbf{Réalisation des opérations de maintenance : du préventif au curatif} \\
            \hline
            II.1 & Exécution de la maintenance préventive planifiée : rigueur et traçabilité \\\hline
            II.2 & Intervention curative sur pannes hydrauliques et mécaniques : diagnostic et réparation \\\hline
            \textbf{III} & \textbf{Bilan de la période de stage : acquis, limites et projections professionnelles} \\
            \hline
            III.1 & Évaluation des compétences techniques et transversales acquises \\\hline
            III.2 & Difficultés rencontrées et recommandations pour l'entreprise et la formation \\\hline
            \end{tabularx}
        \end{center}
        \item \textbf{Justification :} 3 parties équilibrées. Titres parallèles (GN + Complément). Progression logique : Contexte $\rightarrow$ Cœur de métier (Préventif/Curatif ensemble) $\rightarrow$ Réflexivité. Intègre les textes fonctionnels (plans, GMAO, fiches intervention) et iconiques (schémas hydrauliques).
    \end{enumerate}
    \textbf{Conclusion :} Le plan thématique valorise la compétence technique (savoir-faire) plutôt que la simple chronologie (savoir-raconter).
\end{exoresolu}

\courssub{Rédiger l'introduction : les 4 composantes obligatoires}

\begin{methode}[Rédiger une introduction « entonnoir » conforme aux attendus]
    \textbf{Objectif :} Situer, annoncer, engager. L'introduction se rédige \textbf{en dernier} (quand le corps est stable) mais se place \textbf{en premier}.
    \begin{enumerate}
        \item \textbf{La phrase d'accroche / Contexte général (1 phrase) :} Cadre global (formation, examen, obligation réglementaire). \emph{Ex: « Dans le cadre de la préparation du Probatoire de Maintenance Industrielle... »}
        \item \textbf{La présentation du sujet / Circonstances (2-3 phrases) :} Qui ? Quoi ? Où ? Quand ? Pour qui ? (Entreprise, service, dates, tuteur, fonction du stagiaire).
        \item \textbf{La problématique / L'objectif assigné (1 phrase interrogative ou infinitive) :} La question à laquelle le rapport répond. \emph{Ex: « Comment la pratique en entreprise consolide-t-elle les acquis théoriques ? » ou « Analyser l'efficacité du plan de maintenance préventive. »}
        \item \textbf{L'annonce du plan (1 phrase) :} « Ce rapport s'articule en trois parties. La première examine... La seconde analyse... La troisième évalue... »
    \end{enumerate}
    \textbf{Interdits :} « Je vais vous raconter mon stage », « Ce rapport est très intéressant », détails techniques (chiffres, noms de pièces) réservés au corps, familiarité (tutoiement, style oral).
    \textbf{Ton :} Impersonnel (passif, « nous », infinitif) ou « Je » professionnel (rapport de stage). \textbf{Préférence Bac :} « Nous » (inclusif, objectif) ou Passif.
\end{methode}

\begin{exoresolu}[Rédaction d'une introduction complète (Rapport d'incident)]
    \textbf{Énoncé :} Le 15 octobre 2024, vers 10h30, une fuite d'ammoniac (NH$_3$) s'est produite sur le compresseur C-04 de l'unité de réfrigération de l'usine \emph{SOCAPALM} (Dibombari). Vous êtes Technicien Maintenance (1ère Tech). Le Chef d'atelier vous demande un rapport d'incident pour la direction. Rédigez l'introduction.
    \tcblower
    \textbf{Solution détaillée :}
    \begin{quote}
        \textbf{Objet :} Rapport d'incident – Fuite ammoniac Compresseur C-04 – Unité Réfrigération – 15/10/2024.
        
        \textbf{1. Contexte :} Dans le cadre de la surveillance continue des installations frigorifiques de l'usine SOCAPALM de Dibombari, la sécurité des procédés impose la traçabilité de tout événement majeur.
        
        \textbf{2. Circonstances (Sujet) :} Le mardi 15 octobre 2024, à 10h30, une fuite significative d'ammoniac (NH$_3$) a été détectée au niveau du joint d'étanchéité de l'arbre du compresseur C-04, entraînant l'arrêt d'urgence de l'unité et l'évacuation partielle du secteur.
        
        \textbf{3. Objectif / Problématique :} Ce rapport a pour but d'analyser les causes techniques de cette défaillance, d'évaluer les conséquences sur la production et la sécurité, et de proposer des mesures correctives et préventives pour éviter la récidive.
        
        \textbf{4. Annonce du plan :} Nous exposerons dans un premier temps les circonstances exactes et la chronologie de l'intervention (I). Dans un second temps, nous analyserons les causes racines techniques et organisationnelles (II). Enfin, nous formulerons des recommandations techniques et managériales (III).
    \end{quote}
    \textbf{Analyse :} Respect strict du canevas. Vocabulaire précis (joint d'étanchéité, arrêt d'urgence, causes racines). Ton impersonnel/objectif. Pas d'émotion (pas de « tragique » ici, le tragique est une tonalité littéraire à identifier dans des textes supports, pas à écrire dans un rapport technique standard).
\end{exoresolu}

\courssub{Rédiger le corps du rapport : exploitation des sources et argumentation technique}

\begin{methode}[Construire des paragraphes argumentés et exploiter les documents]
    \textbf{Objectif :} Transformer les notes brutes et documents (textes fonctionnels, iconiques) en démonstration structurée.
    \begin{enumerate}
        \item \textbf{Un paragraphe = Une idée force (Thèse) + Preuves (Données/Observations) + Commentaire technique (Analyse).}
        \item \textbf{Exploitation des textes fonctionnels (PV, tableaux, fiches, courriers) :}
            \begin{itemize}
                \item Ne pas recopier : \emph{Synthétiser} (chiffres clés, tendances, écarts).
                \item Intégrer par référence : « Le tableau 2 (Annexe 3) révèle une hausse de 15\% des pannes hydrauliques... »
                \item Citer la source : « Selon le PV de réunion du 12/06 (Annexe 1), le chef service a décidé... »
            \end{itemize}
        \item \textbf{Exploitation des textes iconiques (Photos, schémas, organigrammes, courbes) :}
            \begin{itemize}
                \item \textbf{Décrire} (ce qu'on voit) $\rightarrow$ \textbf{Analyser} (ce que ça veut dire techniquement) $\rightarrow$ \textbf{Interpréter} (lien avec le sujet).
                \item Légender impérativement : \emph{Figure 3 : Vue éclatée du joint défectueux C-04 (source : notice constructeur).}
            \end{itemize}
        \item \textbf{Connecteurs logiques (Cohérence) :}
            \begin{center}
                \begin{tabular}{|c|l|}\hline
                \textbf{Relation} & \textbf{Connecteurs (Variés)} \\\hline
                Ajout & De plus, En outre, Par ailleurs, Également \\\hline
                Cause/Conséquence & Par conséquent, Dès lors, C'est pourquoi, Entraînant \\\hline
                Opposition/Concession & Cependant, Néanmoins, Alors que, Bien que \\\hline
                Illustration & Par exemple, Notamment, C'est le cas de \\\hline
                Conclusion partielle & En définitive, Au total, Force est de constater que \\\hline
            \end{tabular}
            \end{center}
        \item \textbf{Tonalités (Programme Morphosyntaxe) :} Le rapport technique exige le \emph{ton neutre/objectif}. La \emph{tonalité satirique} (ironie, moquerie) ou \emph{tragique} (pathos, fatalité) sont à \textbf{repérer dans les textes supports littéraires} (ex: extrait de roman, article d'opinion), \textbf{pas à produire} dans le rapport technique. Attention à ne pas confondre « ton tragique » (littéraire) et « gravité d'un incident » (factuel).
    \end{enumerate}
\end{methode}

\begin{exoresolu}[Rédaction d'un paragraphe du corps intégrant sources fonctionnelles et iconiques]
    \textbf{Énoncé :} Dans la Partie II de votre rapport de stage CAI (Maintenance préventive/curative), vous devez traiter de l'intervention sur la presse hydraulique (Panne : perte de pression). Documents à disposition :
    \begin{itemize}
        \item \textbf{Texte fonctionnel (Extrait Fiche intervention n°45) :} « Date: 12/06. Equip: Presse 500T. Défaut: Pression chute 180$\rightarrow$120 bar en 10min. Action: Arrêt. Contrôle visuel: fuite raccord flexible HP. Remplacement flexible ref X-22. Remise en route: Pression stable 180 bar. Durée: 45 min. Signé: Tuteur. »
        \item \textbf{Texte iconique (Schéma hydraulique simplifié) :} Pompe $\rightarrow$ Distributeur $\rightarrow$ Vérin. Flexible HP entre Distributeur et Vérin. Capteur pression sur vérin.
    \end{itemize}
    \textbf{Travail :} Rédiger le sous-paragraphe II.2.1 « Intervention curative sur panne hydraulique : diagnostic et réparation » ($\approx$ 80 mots).
    \tcblower
    \textbf{Solution détaillée :}
    \begin{quote}
        \textbf{II.2.1 Intervention curative sur panne hydraulique : diagnostic et réparation}
        
        Le 12 juin, une perte de pression brutale sur la presse 500T (chute de 180 à 120 bar en 10 minutes, \emph{cf. Fiche intervention n°45, Annexe 2}) a imposé l'arrêt immédiat de l'équipement. L'analyse du schéma hydraulique (\emph{Figure 4, Annexe 5}) a permis d'isoler le circuit haute pression entre le distributeur et le vérin. L'inspection visuelle a confirmé une rupture du flexible HP (réf. X-22) au niveau du raccord tournant, point de contrainte mécanique identifié sur le plan. Le remplacement de l'élément défectueux, réalisé en 45 minutes, a rétabli la pression nominale de 180 bar sans oscillation. Cette intervention illustre l'importance de la lecture de plans pour un diagnostic ciblé et la nécessité de prévoir un stock de flexibles critiques pour réduire le temps d'arrêt (TRS).
    \end{quote}
    \textbf{Analyse technique de la rédaction :}
    \begin{itemize}
        \item \textbf{Thèse :} Diagnostic rapide grâce aux documents.
        \item \textbf{Preuves intégrées :} Chiffres précis (180/120 bar, 10 min, 45 min, ref X-22), références croisées (Fiche, Figure, Annexe).
        \item \textbf{Analyse/Commentaire :} « point de contrainte mécanique identifié sur le plan », « illustre l'importance de la lecture de plans », « nécessité de prévoir un stock » (recommandation technique anticipée).
        \item \textbf{Connecteurs :} « a imposé » (conséquence), « a permis » (moyen), « a confirmé » (preuve), « illustre » (généralisation).
        \item \textbf{Ton :} Neutre, technique, passé composé (faits), présent (vérités générales).
    \end{itemize}
\end{exoresolu}

\courssub{Rédiger la conclusion et les recommandations : synthèse et ouverture}

\begin{methode}[Clore le rapport : Bilan, Réponse à la problématique, Recommandations]
    \textbf{Objectif :} La conclusion n'est pas un résumé plat. Elle doit prouver que l'objectif est atteint et ouvrir sur l'action.
    \begin{enumerate}
        \item \textbf{Rappel de la problématique / Objectif :} « Ce stage visait à confronter les acquis théoriques à la réalité industrielle... »
        \item \textbf{Synthèse des résultats majeurs (2-3 points forts) :} « L'analyse a montré que la maintenance préventive réduit de 30\% les arrêts imprévus (Partie II) et que la GMAO optimise la traçabilité (Partie I). » \emph{Chiffres à l'appui si possible.}
        \item \textbf{Réponse directe à la problématique :} « Ainsi, l'adéquation formation/entreprise est globalement satisfaisante, hormis sur la gestion des pièces de rechange critiques. »
        \item \textbf{Recommandations (Numérotées, actionnables, réalistes) :} Distinguer :
            \begin{itemize}
                \item \emph{Pour l'entreprise :} « 1. Constituer un stock de sécurité flexibles HP. 2. Former les opérateurs à la détection précoce. »
                \item \emph{Pour la formation / Personnel :} « 3. Intégrer un module GMAO avancé en Terminale. 4. Simuler des pannes hydrauliques complexes en TP. »
            \end{itemize}
        \item \textbf{Ouverture / Perspective (Optionnelle mais valorisée) :} « L'évolution vers la maintenance prédictive (capteurs IoT) constitue le prochain axe de progression pour le service. »
    \end{enumerate}
    \textbf{Erreur fatale :} Introduire une \textbf{nouvelle information} ou un \textbf{nouvel argument} absent du corps. La conclusion \textbf{ne cite pas de source} (pas de « cf. Annexe »).
    \textbf{Longueur :} $\approx$ 10-15\% du corps (ex: 1/2 page pour un rapport de 5 pages).
\end{methode}

\begin{exoresolu}[Rédaction de la conclusion avec recommandations distinctes]
    \textbf{Énoncé :} Rédigez la conclusion du rapport de stage CAI (Maintenance Industrielle) dont le plan est : I. Intégration \& Outils, II. Opérations Maintenance (Préventif/Curatif), III. Bilan. Problématique : « Adéquation stage/formation ? ».
    \tcblower
    \textbf{Solution détaillée :}
    \begin{quote}
        \textbf{Conclusion}
        
        Au terme de ce stage de trois semaines au service Maintenance de CAI, l'objectif d'évaluer l'adéquation entre la formation de Première Technique en Maintenance Industrielle et la réalité du terrain est atteint.
        
        \textbf{Bilan synthétique :} L'immersion a permis de valider l'acquisition des gestes professionnels de base (lecture de plans ISO, utilisation clé dynamométrique, respect consignes sécurité) et de manipuler l'outil GMAO (Mobility Work) pour la traçabilité des 12 interventions préventives et 3 curatives réalisées. L'analyse de la panne presse (Partie II) a mis en évidence la pertinence de la démarche diagnostique enseignée (Schéma $\rightarrow$ Mesure $\rightarrow$ Action).
        
        \textbf{Toutefois, deux écarts notables persistent :} (1) La gestion des pièces de rechange (rupture stock flexible X-22) relève de la logistique, peu abordée en cours. (2) La maintenance prédictive (analyse vibratoire, thermographie) est absente du référentiel actuel alors qu'elle est déployée sur le compresseur principal.
        
        \textbf{Recommandations :}
        \begin{enumerate}
            \item \textbf{À l'entreprise :} Mettre en place un stock tampon pour les 10 références critiques (flexibles, joints, roulements) et planifier une formation interne « Analyse vibratoire niveau 1 » pour l'équipe.
            \item \textbf{À l'établissement (MINESEC) :} Intégrer un module « Gestion des stocks et Achats techniques » en Terminale et initier à la maintenance 4.0 (capteurs IoT) via un partenariat industriel.
        \end{enumerate}
        
        \textbf{Perspective :} Ce stage confirme mon orientation vers le BTS Maintenance des Systèmes, option Systèmes de Production, avec un intérêt marqué pour la fiabilisation des installations.
    \end{quote}
    \textbf{Validation :} Problématique reprise. Chiffres du corps réutilisés (12 prev, 3 cur, 10 refs). Recommandations séparées Entreprise/École. Ouverture sur orientation. Aucune nouvelle info. Ton professionnel.
\end{exoresolu}

\courssub{Relire, corriger et améliorer : cohérence, syntaxe, présentation}

\begin{methode}[Check-list de révision finale (Auto-évaluation / Pair-évaluation)]
    \textbf{Objectif :} Garantir la qualité formelle et linguistique exigée au Probatoire (notation sur 20 : Fond 10 / Forme 10).
    \begin{enumerate}
        \item \textbf{Structure \& Contenu (Fond - 10 pts) :}
            \begin{itemize}
                \item [ ] Page de garde complète ? Sommaire à jour (numéros pages) ?
                \item [ ] Introduction : 4 composantes ? Problématique claire ? Plan annoncé respecté ?
                \item [ ] Corps : 3 parties équilibrées ? Titres parallèles ? Argumentation technique (thèse/preuves/analyse) ? Exploitation docs (fonctionnels/iconiques) citée (Annexes/Figures) ?
                \item [ ] Conclusion : Reprend problématique ? Synthèse chiffrée ? Recommandations concrètes (Entreprise/École) ? Pas de nouveauté ?
                \item [ ] Annexes : Présentes, numérotées, référencées dans le texte ?
            \end{itemize}
        \item \textbf{Langue \& Style (Forme - 10 pts) :}
            \begin{itemize}
                \item [ ] \textbf{Syntaxe :} Phrases complètes, pas de phrases nominales seules (sauf titres). Subordination (relatives, conjonctives) pour lier idées.
                \item [ ] \textbf{Vocabulaire :} Précision technique (pas de « chose », « truc », « machin »). Termes exacts : \emph{flexible, raccord, pression nominale, TRS, GMAO, étanchéité}.
                \item [ ] \textbf{Connecteurs :} Variés, corrects (pas de « et » en début de phrase, pas de « du coup » oral).
                \item [ ] \textbf{Orthographe/Grammaire :} Accords (participes passés, adjectifs techniques), homophones (à/a, ou/où, ce/se, son/sont). \textbf{Attention :} Accord du participe passé avec « avoir » (COD avant) : « Les pannes \textbf{que} j'ai \textbf{réparées} ».
                \item [ ] \textbf{Ponctuation :} Virgules pour isoler les incises, points-virgules pour énumérations complexes, deux-points pour annoncer explication/liste.
                \item [ ] \textbf{Présentation :} Police unique (Times New Roman 12), interligne 1.5, marges normales, pagination, titres hiérarchisés (gras/souligné), figures/tableaux légendés et centrés.
            \end{itemize}
        \item \textbf{Tonalités (Contrôle) :} Absence de familier, de subjectif (« je pense que », « c'est super »), de satirique/tragique. Ton \textbf{neutre, objectif, impersonnel}.
    \end{enumerate}
    \textbf{Astuce :} Lire \textbf{à voix haute} (détecte les phrases trop longues, lourdes, mal ponctuées). Faire relire par un pair (grille d'évaluation partagée).
\end{methode}

\begin{exoresolu}[Correction d'un extrait défectueux (Révision Forme/Fond)]
    \textbf{Énoncé :} Corrigez et justifiez les erreurs dans cet extrait de corps de rapport (Partie II) :
    \begin{quote}
        « \emph{Dans la deuxieme partie on a vu la maintenance preventive. C'est tres important pour eviter les pannes. Le tableau 1 montre que y'a moins de pannes quand on fait la preventive. Du coup, le chef il a dit qu'il faut continuer. Moi j'ai graissé les roulements et j'ai vérifié les niveaux d'huile. C'etait facile. Par contre la panne du compresseur c'etait tragique, y'avait de l'huile partout, on a eu peur. Heureusement qu'on a reussi a reparer. } »
    \end{quote}
    \tcblower
    \textbf{Solution détaillée : Correction et Justifications}
    \begin{center}
        \begin{tabularx}{\linewidth}{|l|X|X|}\hline
        \rowcolor{popblueL} \textbf{Extrait original (Erreur)} & \textbf{Correction (Style Rapport Technique)} & \textbf{Justification (Code erreur)} \\\hline
        « Dans la deuxieme partie on a vu... » & « La deuxième partie examine la maintenance préventive. » & \textbf{S1 :} « On a vu » (oral, subjectif) $\rightarrow$ Impersonnel/Présent de description. « Deuxième » (ordinal) prend 'e'. \\\hline
        « C'est tres important pour eviter les pannes. » & « Cette stratégie est essentielle pour prévenir les défaillances. » & \textbf{V1 :} « Important » (vague) $\rightarrow$ « Essentielle/Stratégique ». « Éviter » $\rightarrow$ « Prévenir » (terme technique). Accents manquants (très, éviter). \\\hline
        « Le tableau 1 montre que y'a moins de pannes... » & « Le tableau 1 (Annexe 3) indique une réduction de 27\% des arrêts non planifiés grâce à la préventive. » & \textbf{S2 :} « Y'a » (oral) $\rightarrow$ « Il y a » / restructuration. « Montre » $\rightarrow$ « Indique / Révèle ». Chiffre précis obligatoire (exploit doc). \\\hline
        « Du coup, le chef il a dit qu'il faut continuer. » & « Par conséquent, le chef de service a préconisé le maintien du plan préventif. » & \textbf{S3 :} « Du coup » (oral) $\rightarrow$ « Par conséquent / Dès lors ». « Le chef il a dit » (dislocation, familier) $\rightarrow$ « Le chef de service a préconisé / recommandé ». \\\hline
        « Moi j'ai graissé... C'etait facile. » & « J'ai effectué le graissage des roulements et le contrôle des niveaux d'huile, opérations standardisées. » & \textbf{S4 :} « Moi je » (insistance inutile) $\rightarrow$ Supprimé ou « J'ai effectué ». « C'était facile » (subjectif, non pro) $\rightarrow$ « Opérations standardisées / maîtrisées ». Accent « c'était ». \\\hline
        « Par contre la panne du compresseur c'etait tragique... » & « En revanche, l'incident sur le compresseur a engendré une fuite d'huile massive, générant un risque majeur pour la sécurité. » & \textbf{S5/Tonalité :} « Par contre » (oral) $\rightarrow$ « En revanche / Cependant ». \textbf{ERREUR MAJEURE :} « C'était tragique » $\rightarrow$ \textbf{Confusion tonalité littéraire / fait technique}. Le « tragique » est une tonalité littéraire (destin funeste), pas un adjectif technique. Remplacer par factuel : « fuite massive », « risque majeur ». « On a eu peur » (subjectif/émotion) $\rightarrow$ SUPPRIMÉ (hors rapport). « Heureusement » (subjectif) $\rightarrow$ SUPPRIMÉ. « Réparer » $\rightarrow$ « Remettre en état / Résoudre l'incident ». \\\hline
        \end{tabularx}
    \end{center}
    \textbf{Version finale corrigée :}
    \begin{quote}
        « La deuxième partie examine la maintenance préventive. Cette stratégie est essentielle pour prévenir les défaillances : le tableau 1 (Annexe 3) indique une réduction de 27\% des arrêts non planifiés lors de son application stricte. Par conséquent, le chef de service a préconisé le maintien du plan préventif. J'ai effectué le graissage des roulements et le contrôle des niveaux d'huile, opérations standardisées. En revanche, l'incident sur le compresseur a engendré une fuite d'huile massive, générant un risque majeur pour la sécurité. La réparation a été réalisée dans les délais grâce à la disponibilité des joints en stock. »
    \end{quote}
\end{exoresolu}

\courssub{Exploiter textes fonctionnels et iconiques comme sources probantes}

\begin{methode}[Traiter les documents fournis (Textes fonctionnels N° 2, 3, 4 ; Textes iconiques N° 1, 2, 4) pour le rapport]
    \textbf{Objectif :} Le rapport technique s'appuie sur des \emph{preuves documentées}. Le programme impose l'exploitation de 4 textes fonctionnels et 4 textes iconiques au fil de la séquence.
    \begin{enumerate}
        \item \textbf{Typologie des textes fonctionnels (Programme) :}
            \begin{itemize}
                \item \emph{Texte N° 2 :} Procès-verbal (PV) de réunion, Compte-rendu d'intervention, Fiche de poste, Cahier des charges.
                \item \emph{Texte N° 3 :} Tableau statistique (production, pannes, coûts), Planning/Gantt, Budget, Fiche technique constructeur.
                \item \emph{Texte N° 4 :} Courrier administratif (demande, réponse, mise en demeure), Note de service, Email professionnel, Règlement intérieur extrait.
            \end{itemize}
        \item \textbf{Typologie des textes iconiques (Programme) :}
            \begin{itemize}
                \item \emph{Texte N° 1 :} Organigramme (structure entreprise/service), Schéma de principe (hydraulique, électrique, pneumatique), Plan d'implantation.
                \item \emph{Texte N° 2 :} Photographie (état des lieux, pièce défectueuse, poste de travail), Vue éclatée, Coupe technique.
                \item \emph{Texte N° 4 :} Courbe d'évolution (température, pression, vibration, production), Diagramme circulaire (répartition pannes), Histogramme (coûts maintenance), Carte mentale (arbre des causes / 5M).
            \end{itemize}
        \item \textbf{Méthode d'exploitation (3 étapes) :}
            \begin{enumerate}
                \item \textbf{Identifier :} Nature, Source, Date, Auteur, Destinataire, Objectif du document.
                \item \textbf{Extraire :} Informations pertinentes (Chiffres, Décisions, Noms, Fonctions, Cotes, Tendances). \emph{Ne pas tout recopier.}
                \item \textbf{Intégrer :} Dans le corps du rapport par \emph{référence croisée} (Appel de note / Parenthèse).
                    \begin{itemize}
                        \item \emph{Texte fonctionnel :} « Selon le PV de la réunion du 12/06 (Doc 2, Annexe 1), la direction a validé l'achat du stock sécurité. »
                        \item \emph{Tableau/Stat (Doc 3) :} « L'histogramme (Fig 3, Doc 3) révèle que 60\% des pannes sont d'origine hydraulique. »
                        \item \emph{Iconique Schéma/Photo (Doc 1/2) :} « Le schéma hydraulique (Fig 1, Doc 1) localise la fuite au niveau du vérin V3. » / « La photo (Fig 2, Doc 2) montre l'usure anormale du joint torique. »
                        \item \emph{Courbe/Diagramme (Doc 4) :} « La courbe de vibration (Fig 4, Doc 4) franchit le seuil d'alerte le 15/10 à 10h25. »
                    \end{itemize}
            \end{enumerate}
        \item \textbf{Annexes :} Chaque document exploité doit être reproduit (ou résumé s'il est long) en Annexe, numéroté dans l'ordre d'apparition dans le rapport.
    \end{enumerate}
\end{methode}

\begin{exoresolu}[Intégration multi-sources dans une argumentation technique]
    \textbf{Énoncé :} Pour le rapport d'incident SOCAPALM (Fuite Ammoniac C-04), vous disposez de 4 documents. Rédigez le paragraphe d'analyse des causes (Partie II) en les exploitant tous.
    \begin{itemize}
        \item \textbf{Doc 2 (Fonctionnel - PV Réunion sécurité 10/10) :} « Décision : Report de la révision quinquennale compresseurs C-01 à C-04 au budget 2025 (manque trésorerie). Validé par Dir. Usine. »
        \item \textbf{Doc 3 (Fonctionnel - Tableau Suivi Maintenance) :} « C-04 : Dernière révision majeure : 2018. Heures de fonctionnement depuis révision : 42 500 h (Préconisé constructeur : 40 000 h). 3 alertes vibration mineures non traitées (Jan, Mar, Juin 2024). »
        \item \textbf{Doc 1 (Iconique - Organigramme Maintenance) :} Chef Atelier $\rightarrow$ 2 Chefs Équipe $\rightarrow$ 8 Techniciens. \textbf{Poste "Analyste Vibrations" : VACANT depuis 6 mois.}
        \item \textbf{Doc 4 (Iconique - Courbe Vibrations C-04) :} Courbe RMS (mm/s) : Stable 2,5 jusqu'en Jan 2024 $\rightarrow$ Croissance exponentielle $\rightarrow$ 7,8 mm/s le 15/10 (Seuil Alarme 4,5 / Seuil Arrêt 7,0).
    \end{itemize}
    \tcblower
    \textbf{Solution détaillée :}
    \begin{quote}
        \textbf{II. Analyse des causes racines de la fuite ammoniac sur C-04}
        
        L'investigation technique, croisant les données de suivi et l'organisation de la maintenance, met en évidence une convergence de facteurs organisationnels et techniques.
        
        \textbf{1. Cause organisationnelle majeure : Report de la révision majeure.} Le procès-verbal de la réunion de sécurité du 10 octobre (Doc 2, Annexe 1) acte le report de la révision quinquennale du compresseur C-04, préconisée par le constructeur à 40 000 heures, pour raisons budgétaires. Le tableau de suivi maintenance (Doc 3, Annexe 2) confirme que C-04 totalisait 42 500 heures de fonctionnement au 15 octobre, soit \textbf{2 500 heures au-delà de l'échéance constructeur}.
        
        \textbf{2. Défaillance de la surveillance prédictive : Poste vacant.} L'organigramme du service maintenance (Doc 1, Annexe 3) révèle la vacance du poste d'Analyste Vibrations depuis six mois. Cette absence explique l'absence de traitement des trois alertes vibratoires mineures signalées en janvier, mars et juin 2024 (Doc 3).
        
        \textbf{3. Confirmation par l'analyse vibratoire (Cause technique immédiate).} La courbe d'évolution des vibrations RMS (Doc 4, Figure 1) illustre une dégradation exponentielle de l'état mécanique : franchissement du seuil d'alarme (4,5 mm/s) dès février 2024 et atteinte du seuil d'arrêt d'urgence (7,0 mm/s) à 7,8 mm/s le jour de l'incident (15/10, 10h25). Cette signature vibratoire est caractéristique d'un défaut de roulement ou de désalignement progressif, non détecté faute d'analyste.
        
        \textbf{Synthèse causale (Arbre des causes / 5M) :} La fuite du joint d'étanchéité (conséquence) résulte de la vibration excessive (cause technique) due à l'usure roulement non surveillée (cause détection), elle-même due à la vacance du poste analyste (cause organisationnelle) et au report de la révision préventive majeure (cause racine gestion).
    \end{quote}
    \textbf{Analyse de l'intégration :} Chaque document est cité \textbf{une fois} à sa première utilisation (Doc X, Annexe Y / Figure Z). Les données brutes (heures, seuils, dates) sont transformées en \textbf{argumentation causale}. Le ton est neutre (« acte », « confirme », « révèle », « illustre »). L'arbre des causes (5M : Main d'œuvre, Milieu, Machine, Méthode, Management) structure la conclusion du paragraphe.
\end{exoresolu}

\begin{popfigure}
\begin{tikzpicture}[
    mindmap,
    concept color=popblue,
    level 1 concept/.append style={level distance=3.5cm, sibling angle=72, concept color=popblue!60},
    level 2 concept/.append style={level distance=2.5cm, concept color=poporange!60},
    text=white,
    font=\small
]
\node[concept, align=center]{Fuite Joint\\Étanchéité (Conséquence)}
    child[grow=90] { node[concept, align=center]{Vibration Excessive\\ (Cause Technique)}
        child[grow=60] { node[concept, align=center]{Usure Roulement\\ Non Surveillée} }
        child[grow=120] { node[concept, align=center]{Désalignement\\ Progressif} }
    }
    child[grow=18] { node[concept, align=center]{Défaillance Détection\\ (Cause Orga) }
        child[grow=-18] { node[concept, align=center]{Poste Analyste\\ Vacant (Main d'œuvre)} }
    }
    child[grow=-54] { node[concept, align=center]{Gestion Maintenance\\ (Cause Racine) }
        child[grow=-90] { node[concept, align=center]{Report Révision\\ Quinquennale (Management)} }
        child[grow=-36] { node[concept, align=center]{Budget Restreint\\ (Milieu)} }
    }
    child[grow=-126] { node[concept, align=center]{Méthode\\ (Méthode) }
        child[grow=-162] { node[concept, align=center]{Absence Surveillance\\ Prédictive} }
    }
    child[grow=-198] { node[concept, align=center]{Machine\\ (Machine) }
        child[grow=-234] { node[concept, align=center]{Compresseur C-04\\ > 40 000 h} }
    };
\end{tikzpicture}
\legende{Arbre des causes (Méthode 5M) pour l'incident SOCAPALM C-04.}
\end{popfigure}

\begin{attention}[Les 5 erreurs qui coûtent des points au Bac (Probatoire Technique - Français)]
    \begin{enumerate}
        \item \textbf{Confondre « Tonalité littéraire » et « Ton du rapport » :} Utiliser un vocabulaire \emph{satirique} (ironie, moquerie : « c'est n'importe quoi ») ou \emph{tragique} (pathos, fatalité : « c'était horrible, on a eu peur ») dans la rédaction du rapport technique. \textbf{Règle :} Le rapport exige un \textbf{ton neutre, objectif, impersonnel, technique}. Les tonalités satirique/tragique sont à \textbf{analyser} dans les textes supports littéraires fournis (Réception), pas à \textbf{produire} (Production).
        \item \textbf{Introduction « Coquille vide » :} Omettre l'une des 4 composantes (souvent la \textbf{Problématique} ou l'\textbf{Annonce du plan} précise). Rédiger une introduction générique (« Ce rapport parle de mon stage... ») sans citer l'entreprise, les dates, le tuteur, ni la question technique précise.
        \item \textbf{Corps « Catalogue / Journal de bord » :} Suivre un plan chronologique (Lundi..., Mardi...) au lieu d'un plan \textbf{thématique/analytique} (Préventif / Curatif / Gestion). Absence d'\textbf{analyse} (que du récit : « j'ai fait ci, j'ai fait ça ») sans exploitation des documents fournis (textes fonctionnels/iconiques non cités, annexes fantômes).
        \item \textbf{Conclusion « Nouveau départ » :} Introduire de nouveaux arguments, de nouveaux chiffres, ou de nouveaux documents absents du corps. Oublier les \textbf{Recommandations} (ou les rendre vagues : « il faut faire attention » au lieu de « installer un capteur vibration sur C-04 avant le 30/11 »). Ne pas répondre explicitement à la problématique.
        \item \textbf{Négligence linguistique « Coût Forme » (10 pts) :} Familier oralisé (« du coup », « y'a », « moi je », « c'est facile »), syntaxe hachée (phrases nominales), vocabulaire imprécis (« truc », « chose », « machiner »), accords fautifs (participes passés, termes techniques), absence de connecteurs logiques, présentation soignée (sommaire faux, annexes non référencées, figures sans légende).
    \end{enumerate}
\end{attention}

\newpage

\coursec{Exercices}

\courssub{Je vérifie mes connaissances}

\begin{exercice}[QCM : Structure et vocabulaire du rapport]
Chaque question ne comporte qu'une seule réponse exacte. Cochez la bonne proposition.
\begin{enumerate}
    \item La structure \textbf{externe} d'un rapport comprend obligatoirement :
    \begin{enumerate}
        \item L'introduction, le développement, la conclusion.
        \item La page de garde, la table des matières, le corps du texte, les annexes.
        \item L'objet, les causes, les conséquences, les solutions.
        \item Le titre, l'auteur, la date, le destinataire, l'objet, la signature.
    \end{enumerate}
    \item Dans l'\textbf{introduction} d'un rapport, on ne trouve \textbf{pas} :
    \begin{enumerate}
        \item Le contexte et le déclencheur de la mission.
        \item L'objet précis de la mission.
        \item L'analyse détaillée des résultats et des constatations.
        \item La méthode de travail et l'annonce du plan.
    \end{enumerate}
    \item Le \textbf{corps} du rapport doit présenter les faits :
    \begin{enumerate}
        \item De manière subjective, en donnant son avis personnel.
        \item Selon un ordre chronologique ou logique, avec objectivité.
        \item Uniquement sous forme de tableaux statistiques bruts.
        \item En utilisant la tonalité satirique pour critiquer les dysfonctionnements.
    \end{enumerate}
    \item La \textbf{conclusion} d'un rapport technique se termine généralement par :
    \begin{enumerate}
        \item Une formule de politesse adressée au destinataire.
        \item Des recommandations concrètes et opérationnelles.
        \item Un résumé littéraire de l'intrigue.
        \item La signature de l'auteur (qui se place avant la conclusion).
    \end{enumerate}
\end{enumerate}
\end{exercice}

\begin{exercice}[Vrai ou Faux justifié : Caractéristiques du rapport]
Pour chaque affirmation, cochez \textbf{Vrai} ou \textbf{Faux} et justifiez brièvement votre réponse en une phrase.
\begin{enumerate}
    \item L'objet d'un rapport peut être formulé sous forme de question.
    \item Le rapport admet l'emploi du « je » pour marquer l'engagement personnel de l'auteur.
    \item Les annexes (photos, tableaux bruts, questionnaires) sont numérotées et citées dans le corps du texte.
    \item La tonalité tragique est recommandée pour rédiger le compte rendu d'une visite d'entreprise.
    \item La concordance des temps n'est pas obligatoire dans un rapport car c'est un texte administratif.
\end{enumerate}
\end{exercice}

\begin{exercice}[Définitions et associations : Terminologie technique]
Reliez chaque terme de la colonne de gauche à sa définition exacte dans la colonne de droite (écrivez les paires : exemple : 1--c).
\begin{center}
\begin{tabularx}{\linewidth}{|>{\bfseries}l|X|>{\bfseries}l|X|}\hline
\textbf{Terme} & \textbf{Définition} & \textbf{Terme} & \textbf{Définition} \\\hline
1. Objet & a. Partie qui expose les faits, l'analyse et les résultats. & 5. Annexe & e. Éléments joints en fin de rapport pour étayer le propos. \\\hline
2. Introduction & b. Formule finale adressée au destinataire. & 6. Recommandation & f. Proposition d'action concrète issue des conclusions. \\\hline
3. Corps & c. Phrase nominale résumant le sujet exact de la mission. & 7. Formule d'envoi & g. Première partie : contexte, objet, méthode, plan. \\\hline
4. Conclusion & d. Synthèse des constats et ouverture sur les suites à donner. & 8. Plan & h. Organisation logique des parties et sous-parties. \\\hline
\end{tabularx}
\end{center}
\end{exercice}

\begin{exercice}[Ordre opératoire : De la situation au rapport final]
Remettez les étapes suivantes dans l'ordre chronologique logique de la rédaction d'un rapport (numérotez de 1 à 7).
\begin{enumerate}
    \item Rédaction de l'introduction, du corps et de la conclusion.
    \item Analyse de la situation de communication (destinateur, destinataire, objet, finalité).
    \item Récolte et tri des informations (observations, entretiens, documents).
    \item Relecture, correction et amélioration (orthographe, cohérence, tonalité).
    \item Élaboration du plan détaillé (structure interne : parties, sous-parties).
    \item Mise en page finale (structure externe : page de garde, sommaire, annexes).
    \item Choix de la structure du plan (chronologique, thématique, analytique, etc.).
\end{enumerate}
\end{exercice}

\courssub{Je m'entraîne}

\begin{exercice}[Réordonnancement et complétion : Rapport d'accident]
Voici les paragraphes mélangés d'un rapport sur un accident de circulation survenu devant le lycée technique de Nkolbisson. L'objet et la conclusion manquent.
\begin{enumerate}
    \item \textbf{Remettez les paragraphes dans l'ordre logique} (numérotez-les de 1 à 5) pour constituer le \textbf{corps} du rapport.
    \item \textbf{Rédigez l'Objet} du rapport (une phrase nominale précise).
    \item \textbf{Rédigez la Conclusion} (synthèse + 2 recommandations + formule d'envoi adaptée).
\end{enumerate}
\textbf{Paragraphes à ordonner :}
\begin{itemize}
    \item[A.] \emph{Analyse des causes :} L'enquête révèle deux facteurs majeurs : l'excès de vitesse du véhicule (estimé à 60 km/h en zone 30) et l'absence de passage piéton signalisé à cet endroit précis.
    \item[B.] \emph{Constats matériels et humains :} Un élève de 1re TEE a été légèrement blessé au bras (évacué au dispensaire scolaire). Le pare-chocs avant du véhicule est arraché. La circulation a été perturbée pendant 45 minutes.
    \item[C.] \emph{Circonstances :} Le mardi 14 octobre 2025, vers 12h30, un véhicule de type Toyota Corolla, immatriculé LT-123-AB, a percuté un élève traversant la chaussée devant le portail principal du lycée.
    \item[D.] \emph{Témoignages :} Trois élèves témoins confirment que le conducteur n'a pas freiné avant l'impact. Le conducteur, M. Essomba, reconnaît une inattention due à un appel téléphonique.
    \item[E.] \emph{Mesures immédiates :} Le proviseur a fait appel aux services de santé scolaire et à la police. Le conducteur a été conduit au commissariat pour audition. La zone a été sécurisée par les agents de sécurité de l'établissement.
\end{itemize}
\end{exercice}

\begin{exercice}[Exploitation de données statistiques : Paragraphe « Constats et analyse »]
Le conseiller d'orientation du Lycée Technique de Bafoussam doit rédiger un rapport sur l'absentéisme en classe de Première Technique. Il dispose du tableau suivant (effectif total : 120 élèves).
\begin{center}
\begin{tabular}{|c|c|c|c|c|c|}\hline
\textbf{Mois} & \textbf{Sept.} & \textbf{Oct.} & \textbf{Nov.} & \textbf{Déc.} & \textbf{Janv.} \\\hline
\textbf{Nb d'absences totales} & 45 & 78 & 156 & 92 & 60 \\\hline
\textbf{Nb d'élèves concernés} & 12 & 18 & 45 & 22 & 15 \\\hline
\end{tabular}
\end{center}
\emph{Consigne :} À partir de ces données, rédigez le paragraphe \textbf{« Constats et analyse »} du corps du rapport (environ 150 mots). Vous calculerez les taux d'absentéisme pertinents, identifierez le pic, formulerez des hypothèses explicatives (contexte scolaire) et respecterez la tonalité objective du rapport.
\end{exercice}

\begin{exercice}[Identification des tonalités et réécriture objective]
Lisez les trois extraits ci-dessous.
\begin{enumerate}
    \item \textbf{Identifiez la tonalité} de chaque extrait (\cle{satirique} ou \cle{tragique}).
    \item \textbf{Justifiez} en relevant au moins deux procédés stylistiques par extrait (lexique, figures de style, ponctuation, registre de langue).
    \item \textbf{Transformez l'extrait satirique} en un paragraphe objectif adapté au \textbf{corps d'un rapport} administratif.
\end{enumerate}
\textbf{Extrait 1 :} « Ah ! La magnifique cantine flambant neuve ! Ses portes closes depuis six mois offrent un spectacle réconfortant : celui de l'herbe qui pousse fièrement dans la salle de repas. On dirait un jardin botanique gratuit, offert par la bienveillance de l'administration aux estomacs vides des élèves. Bravo pour cette « restauration » de haut vol ! »
\textbf{Extrait 2 :} « Le silence pesant de la salle de réanimation ne fut rompu que par le bip régulier du moniteur. Le jeune apprenti, plein de vie ce matin encore, gisait là, inerte, victime d'une machine défaillante. L'espoir s'éteignait goutte à goutte, comme le sang sur le béton froid de l'atelier. »
\textbf{Extrait 3 :} « Lors de la visite du 12 mars, l'équipe a constaté la conformité des extincteurs. Le registre de sécurité est tenu à jour. Aucune anomalie n'a été relevée sur les issues de secours. »
\end{exercice}

\begin{exercice}[Correction et amélioration : Rapport d'élève]
Voici un extrait de rapport rédigé par un élève de 1re TMA. Il contient \textbf{cinq erreurs ciblées} : (1) Objet absent, (2) Opinion personnelle non étayée, (3) Désordre logique, (4) Faute de concordance des temps, (5) Formule de conclusion inadaptée.
\begin{enumerate}
    \item \textbf{Identifiez et surlignez (ou citez)} les cinq erreurs dans le texte.
    \item \textbf{Proposez une version corrigée et améliorée} de l'extrait (introduction + début du corps).
\end{enumerate}
\textbf{Texte de l'élève :}
« \textbf{Introduction :} Je vais vous parler de mon stage. C'était super intéressant. \textbf{Objet :} Stage. J'ai fait mon stage à la SNEC à Douala du 1er au 30 juin. \textbf{Corps :} D'abord, j'ai été accueilli par le chef. Il est très sympa. Ensuite, j'ai vu comment on répare les transformateurs. C'était difficile mais j'ai aimé. Le chef m'a dit que j'étais le meilleur stagiaire qu'il a eu. Après, j'ai fait le câblage d'un tableau. Au début, j'ai visité les ateliers. \textbf{Conclusion :} Bisous à tout le monde, merci pour tout. Signé : Moumi. »
\end{exercice}

\courssub{Je me prépare au Bac}

\begin{exercice}[Sujet type Baccalauréat Technique : Rapport de stage BTP]
\textbf{Situation :} Vous êtes élève en \textbf{1re TBTP} (Bâtiment et Travaux Publics) au Lycée Technique de la Construction à Yaoundé. Vous avez effectué un \textbf{stage d'un mois (du 2 juillet au 2 août 2025)} dans l'entreprise \textbf{« BATICAM S.A.R.L. »}, chantier de construction d'un immeuble R+3 au quartier Mvog-Ada.
\textbf{Consigne :} Rédigez le \textbf{rapport de stage complet} à remettre à votre chef d'établissement (Proviseur). Votre rapport doit impérativement comporter :
\begin{enumerate}
    \item La \textbf{structure externe complète} (page de garde, sommaire, corps, annexes — présentez la maquette).
    \item Une \textbf{Introduction} standardisée (contexte, objet, lieu, dates, méthode, plan annoncé).
    \item Un \textbf{Corps structuré en 3 parties logiques} :
        \begin{itemize}
            \item Partie 1 : Présentation de l'entreprise et du chantier (organigramme, effectifs, nature de l'ouvrage).
            \item Partie 2 : Description chronologique des activités réalisées par semaine (fondations, élévation, coffrage/ferraillage/bétonnage, finitions).
            \item Partie 3 : Analyse technique (compétences acquises, difficultés rencontrées, respect des normes de sécurité).
        \end{itemize}
    \item Une \textbf{Conclusion} (bilan personnel, apports pour la formation, 2 recommandations pour l'entreprise/école, formule d'envoi).
\end{enumerate}
\textbf{Données fournies :} Encadreur : Ingénieur TCHOUA ; Tuteur école : M. FOPA ; Effectif chantier : 45 ouvriers ; Matériaux : Béton C25/30, Fer HA 500 ; Équipements : Grue à tour, bétonnière, vibreur.
\end{exercice}

\begin{exercice}[Rédaction d'introduction à partir de documents : Enquête insalubrité]
\textbf{Situation :} Vous êtes délégué à l'hygiène et à l'environnement de la Mairie d'Arrondissement de Yaoundé 2. Suite à une plainte collective des riverains, vous avez mené une enquête sur l'insalubrité du \textbf{Marché Mokolo} (quartier Mfoundi) les 10, 11 et 12 septembre 2025.
\textbf{Documents fournis :}
\begin{itemize}
    \item \textbf{Document 1 (Photographie - description) :} Vue aérienne partielle : tas d'ordures ménagères non collectés depuis 8 jours au niveau de l'entrée principale (voie d'accès bloquée), eaux usées stagnantes (coloration verdâtre, odeur nauséabonde) le long des étals de poissonnerie, présence de rongeurs en plein jour, poubelles débordantes et éventrées.
    \item \textbf{Document 2 (Témoignage 1 - Marchande tomates) :} « On paie la taxe journalière de 200 F mais les éboueurs ne passent plus. Les rats mangent nos marchandises la nuit. Mes enfants toussent tout le temps à cause de la poussière et de l'odeur. »
    \item \textbf{Document 3 (Témoignage 2 - Chauffeur taxi) :} « La route est impraticable. Les ordures montent jusqu'au milieu de la chaussée. On perd 30 min par trajet. Hier, un pneu a crevé sur un ferraille caché dans la boue. »
    \item \textbf{Document 4 (Témoignage 3 - Infirmier dispensaire marché) :} « Hausse de 40\% des consultations pour diarrhées aiguës et infections cutanées depuis 3 semaines. Corrélation directe avec l'absence de collecte et l'eau stagnante. »
\end{itemize}
\textbf{Consigne :} Rédigez \textbf{uniquement l'Introduction complète} du rapport d'enquête (Contexte/Déclencheur, Objet précis, Méthodologie, Plan annoncé). Respectez les conventions administratives (ton neutre, impersonnel, concordance des temps).
\end{exercice}

\begin{exercice}[Rapport d'activité : Analyse, planification et rédaction partielle]
\textbf{Situation :} Vous êtes le \textbf{Président du Club « Jeunes Écologistes »} de votre lycée technique (effectif : 60 adhérents). Le Proviseur vous demande le \textbf{rapport d'activités du 1er trimestre (septembre--décembre 2025)} pour le Conseil d'Administration.
\textbf{Activités réalisées :}
\begin{itemize}
    \item 15/09 : Assemblée générale de rentrée (45 présents) -- élection bureau, cotisation 1000 F.
    \item 05/10 : Journée « Lycée Propre » -- ramassage 50 kg plastiques, sensibilisation 200 élèves.
    \item 18/10 : Conférence « Changement climatique et BTP » avec Ingénieur NGONO (amphi plein, 120 participants).
    \item 12/11 : Pépinière : mise en sachets 500 plants (acacia, eucalyptus) -- partenariat MINEPDED.
    \item 03/12 : Sortie pédagogique Parc de la Mefou -- 30 adhérents, étude biodiversité.
    \item Gestion : Recettes 65 000 F (cotisations + subvention mairie 15 000 F) ; Dépenses 48 500 F (transport, plants, affiches, goûter) ; Solde : 16 500 F.
\end{itemize}
\textbf{Consignes :}
\begin{enumerate}
    \item \textbf{Analysez la situation de communication} (Destinateur, Destinataire, Objet, Finalité, Tonalité attendue).
    \item \textbf{Élaborez le plan détaillé} (structure interne) du rapport en 3 parties équilibrées avec sous-parties.
    \item \textbf{Rédigez intégralement le Corps du rapport} (les 3 parties développées) et la \textbf{Conclusion} (bilan moral/financier, perspectives 2e trimestre, recommandations, formule d'envoi).
\end{enumerate}
\end{exercice}

\newpage

\coursec{Corrigés des exercices}

\coursec{Leçon 11 : Production d'écrits utilitaires — Rédiger un rapport (Corrigés des exercices)}

\begin{corrige}[Exercice 1 — QCM : Structure et vocabulaire du rapport]
\begin{enumerate}
    \item \textbf{Réponse : d.} La structure externe (les éléments d'identification visibles) comprend : le titre, l'auteur (destinateur), la date, le destinataire, l'objet et la signature. La proposition \emph{b} mélange structure externe et organisation matérielle ; la proposition \emph{a} décrit la structure \textbf{interne} (introduction, corps, conclusion) ; la proposition \emph{c} décrit un contenu possible du corps.
    \item \textbf{Réponse : c.} L'analyse détaillée des résultats et des constatations appartient au \textbf{corps} du rapport, pas à l'introduction. Celle-ci contient uniquement : le contexte et le déclencheur (a), l'objet précis (b), la méthode et l'annonce du plan (d).
    \item \textbf{Réponse : b.} Le corps du rapport présente les faits avec \cle{objectivité}, selon un ordre chronologique ou logique. Le rapport exclut la subjectivité (a), ne se limite pas aux tableaux (c) et proscrit la tonalité satirique (d), incompatible avec le ton neutre et impersonnel exigé.
    \item \textbf{Réponse : b.} La conclusion d'un rapport technique s'achève par des \cle{recommandations} concrètes et opérationnelles (suites à donner). La formule de politesse (a) peut précéder la signature, mais la signature se place \emph{après} la conclusion (d est donc faux), et le résumé littéraire (c) n'a pas sa place dans un rapport.
\end{enumerate}
\end{corrige}

\begin{corrige}[Exercice 2 — Vrai ou Faux justifié]
\begin{enumerate}
    \item \textbf{Vrai.} L'objet peut être formulé sous forme interrogative (ex. : « Comment expliquer la baisse de fréquentation de la bibliothèque ? »), à condition de rester précis ; la forme la plus courante reste cependant la phrase nominale.
    \item \textbf{Faux.} Le rapport exige un ton \cle{neutre et impersonnel} : on emploie « nous », « l'équipe », « le soussigné » ou le passif, jamais un « je » marquant l'engagement personnel.
    \item \textbf{Vrai.} Les annexes sont numérotées (annexe 1, annexe 2\ldots) et doivent être \textbf{appelées} dans le corps du texte (ex. : « cf. annexe 2 »), sinon elles perdent leur valeur de preuve.
    \item \textbf{Faux.} Le compte rendu d'une visite d'entreprise exige la tonalité \textbf{objective} (constats factuels) ; la tonalité tragique relève du texte littéraire et dénaturerait le propos.
    \item \textbf{Faux.} La \cle{concordance des temps} est obligatoire dans tout texte rédigé, y compris administratif : un rapport relate des faits passés (passé composé, imparfait) et formule des recommandations (présent, futur) de manière cohérente.
\end{enumerate}
\end{corrige}

\begin{corrige}[Exercice 3 — Définitions et associations]
Les paires correctes sont :
\begin{center}
\begin{tabular}{|c|c|c|c|}\hline
1 -- c & 2 -- g & 3 -- a & 4 -- d \\\hline
5 -- e & 6 -- f & 7 -- b & 8 -- h \\\hline
\end{tabular}
\end{center}
\textbf{Justifications :} l'\textbf{objet} (1--c) est une phrase nominale résumant la mission ; l'\textbf{introduction} (2--g) présente contexte, objet, méthode et plan ; le \textbf{corps} (3--a) expose faits, analyses et résultats ; la \textbf{conclusion} (4--d) synthétise les constats et ouvre sur les suites ; l'\textbf{annexe} (5--e) regroupe les pièces jointes ; la \textbf{recommandation} (6--f) propose une action concrète ; la \textbf{formule d'envoi} (7--b) s'adresse au destinataire en fin de texte ; le \textbf{plan} (8--h) organise logiquement parties et sous-parties.
\end{corrige}

\begin{corrige}[Exercice 4 — Ordre opératoire]
L'ordre chronologique logique est le suivant :
\begin{enumerate}
    \item \textbf{1.} Analyse de la situation de communication (destinateur, destinataire, objet, finalité).
    \item \textbf{2.} Récolte et tri des informations (observations, entretiens, documents).
    \item \textbf{3.} Choix de la structure du plan (chronologique, thématique, analytique\ldots).
    \item \textbf{4.} Élaboration du plan détaillé (structure interne : parties, sous-parties).
    \item \textbf{5.} Rédaction de l'introduction, du corps et de la conclusion.
    \item \textbf{6.} Relecture, correction et amélioration (orthographe, cohérence, tonalité).
    \item \textbf{7.} Mise en page finale (structure externe : page de garde, sommaire, annexes).
\end{enumerate}
\textbf{Logique :} on ne peut choisir un plan qu'après avoir récolté et trié les informations ; la relecture précède toujours la mise en page définitive, qui est la dernière opération avant la remise du document.
\end{corrige}

\begin{corrige}[Exercice 5 — Rapport d'accident]
\textbf{1. Ordre logique des paragraphes (corps du rapport) :}
\begin{center}
\textbf{C (1) $\rightarrow$ B (2) $\rightarrow$ D (3) $\rightarrow$ A (4) $\rightarrow$ E (5)}
\end{center}
\textbf{Justification :} l'ordre est \cle{chronologique et logique} : on décrit d'abord les \emph{circonstances} de l'accident (C), puis les \emph{constats} matériels et humains immédiats (B), ensuite les \emph{témoignages} recueillis (D), puis l'\emph{analyse des causes} qui découle des témoignages et des constats (A), enfin les \emph{mesures immédiates} prises (E).

\textbf{2. Objet (proposition) :} « Rapport sur l'accident de circulation survenu le mardi 14 octobre 2025 devant le portail principal du Lycée Technique de Nkolbisson. »

\textbf{3. Conclusion (proposition) :}
« En synthèse, cet accident, qui a heureusement fait un blessé léger, résulte de la conjonction d'un excès de vitesse et de l'absence de passage piéton sécurisé devant l'établissement. Afin de prévenir tout renouvellement, il est recommandé : (1) de faire installer par les services municipaux un passage piéton signalisé et des ralentisseurs devant le portail principal ; (2) d'organiser une campagne de sensibilisation à la sécurité routière à l'intention des élèves et des usagers de la route. Nous restons à la disposition de Monsieur le Proviseur pour tout complément d'information et le prions d'agréer l'expression de notre respectueuse considération. »
\end{corrige}

\begin{corrige}[Exercice 6 — Paragraphe « Constats et analyse »]
\textbf{Calculs préalables (taux d'absentéisme = élèves concernés / effectif total $\times$ 100) :}
\begin{itemize}
    \item Septembre : $12/120 \times 100 = 10\,\%$ ; Octobre : $18/120 \times 100 = 15\,\%$ ;
    \item Novembre : $45/120 \times 100 = 37,5\,\%$ ; Décembre : $22/120 \times 100 \approx 18,3\,\%$ ; Janvier : $15/120 \times 100 = 12,5\,\%$.
    \item Pic : novembre, avec 156 absences et 45 élèves concernés, soit une multiplication par $3,75$ du nombre d'élèves concernés entre septembre et novembre ($45/12 = 3,75$).
\end{itemize}
\textbf{Paragraphe rédigé (proposition) :}
« L'examen des données du premier trimestre fait apparaître une évolution irrégulière de l'absentéisme en classe de Première Technique. Le taux d'élèves concernés passe de 10\,\% en septembre (12 élèves, 45 absences) à 15\,\% en octobre, avant d'atteindre un pic net en novembre : 156 absences totalisées par 45 élèves, soit 37,5\,\% de l'effectif. Ce chiffre est près de quatre fois supérieur à celui de la rentrée. Le phénomène reflue ensuite : 18,3\,\% en décembre (92 absences) et 12,5\,\% en janvier (60 absences). Le pic de novembre peut s'expliquer par la fatigue accumulée en milieu de trimestre, par la multiplication des évaluations de fin de séquence et par les pluies tardives qui compliquent les déplacements. La baisse de décembre coïncide avec l'approche des vacances et le resserrement du suivi des absences par l'administration. Ces constats appellent des mesures ciblées sur le mois de novembre. »
\end{corrige}

\begin{corrige}[Exercice 7 — Tonalités et réécriture]
\textbf{1. Identification :} Extrait 1 : tonalité \cle{satirique}. Extrait 2 : tonalité \cle{tragique}. Extrait 3 : tonalité \cle{objective} (ni satirique ni tragique : c'est le modèle du rapport).

\textbf{2. Justifications :}
\begin{itemize}
    \item \emph{Extrait 1 (satirique) :} ironie mordante (« magnifique cantine flambant neuve » alors qu'elle est fermée ; « Bravo ! ») ; antiphrase (« restauration de haut vol ») ; comparaison dévalorisante (« un jardin botanique gratuit ») ; lexique élogieux détourné (« bienveillance », « réconfortant ») pour dénoncer la négligence de l'administration.
    \item \emph{Extrait 2 (tragique) :} lexique de la mort et de la douleur (« gisait, inerte », « sang », « espoir s'éteignait ») ; opposition pathétique entre la vie et la mort (« plein de vie ce matin encore » / « inerte ») ; images sensorielles angoissantes (« silence pesant », « bip régulier ») ; gradation descendante (« l'espoir s'éteignait goutte à goutte »).
    \item \emph{Extrait 3 (objectif) :} lexique technique et neutre (« conformité », « registre », « anomalie ») ; phrases déclaratives courtes ; absence totale d'émotion et de jugement.
\end{itemize}
\textbf{3. Réécriture objective de l'extrait 1 (proposition) :}
« La nouvelle cantine de l'établissement, bien qu'achevée, demeure fermée depuis six mois. Lors de la visite, l'équipe a constaté que la végétation envahit la salle de restauration, faute d'entretien. Cette situation prive les élèves du service de restauration scolaire prévu. Il est recommandé de procéder à l'ouverture effective de la cantine et à l'entretien régulier des locaux. »
\end{corrige}

\begin{corrige}[Exercice 8 — Correction et amélioration]
\textbf{1. Les cinq erreurs :}
\begin{enumerate}
    \item \textbf{Objet absent ou vide :} « Objet : Stage » est imprécis ; il faut une phrase nominale complète (nature, lieu, période).
    \item \textbf{Opinion personnelle non étayée :} « C'était super intéressant », « Il est très sympa », « j'étais le meilleur stagiaire » : jugements subjectifs sans faits ni preuves.
    \item \textbf{Désordre logique :} « Au début, j'ai visité les ateliers » arrive \emph{après} le câblage d'un tableau ; la visite des ateliers doit précéder les tâches techniques.
    \item \textbf{Faute de concordance des temps :} « le meilleur stagiaire qu'il \textbf{a} eu » doit être corrigé en « qu'il \textbf{ait} eu » (subjonctif passé après la principale au passé « m'a dit » ; accord du participe passé « eu » au masculin singulier, COD « que » = « le meilleur stagiaire » placé avant).
    \item \textbf{Formule de conclusion inadaptée :} « Bisous à tout le monde » est familier et inacceptable dans un rapport adressé à un chef d'établissement.
\end{enumerate}
\textbf{2. Version corrigée (proposition) :}
« \textbf{Objet :} Rapport de stage d'observation et de pratique effectué à la SNEC de Douala du 1er au 30 juin 2025.
\textbf{Introduction :} Dans le cadre de la formation en 1re TMA, un stage d'un mois a été effectué à la SNEC de Douala, du 1er au 30 juin 2025. Ce stage avait pour objet la découverte des métiers de la maintenance des équipements électriques. La méthode adoptée a consisté en l'observation des ateliers, la participation aux tâches confiées et la prise de notes quotidienne. Le présent rapport présente d'abord le déroulement du stage, puis les acquis techniques.
\textbf{Corps (début) :} Dès le premier jour, le stagiaire a été accueilli par le chef d'atelier, qui a organisé une visite guidée des différents ateliers. Cette phase d'observation a permis de se familiariser avec l'organisation de l'entreprise. Ensuite, le stagiaire a assisté aux opérations de réparation des transformateurs, avant de participer, sous supervision, au câblage d'un tableau électrique. »
\end{corrige}

\begin{corrige}[Exercice 9 — Sujet type Bac : Rapport de stage BTP]
\textbf{Barème indicatif (sur 20) :} structure externe et présentation : 4 pts ; introduction conforme : 4 pts ; corps en 3 parties logiques et documentées : 8 pts ; conclusion (bilan, recommandations, formule d'envoi) : 3 pts ; correction de la langue : 1 pt.

\textbf{Éléments de correction :}
\begin{enumerate}
    \item \textbf{Structure externe (maquette) :} page de garde (nom de l'établissement, titre « Rapport de stage effectué à BATICAM S.A.R.L. », nom de l'élève, classe de 1re TBTP, année scolaire) ; sommaire ; corps paginé ; annexes (plan du chantier, photos, fiches de présence).
    \item \textbf{Introduction :} contexte (stage obligatoire dans la formation BTP) ; objet (phrase nominale : nature, entreprise, chantier R+3 à Mvog-Ada) ; dates (2 juillet -- 2 août 2025) ; encadrement (Ing. TCHOUA, tuteur école M. FOPA) ; méthode (observation, participation aux tâches, prise de notes) ; annonce du plan en 3 parties.
    \item \textbf{Corps :}
    \begin{itemize}
        \item \emph{Partie 1 :} présentation de BATICAM S.A.R.L., organigramme simplifié, effectif de 45 ouvriers, nature de l'ouvrage (immeuble R+3).
        \item \emph{Partie 2 :} description chronologique par semaine : semaine 1, terrassement et fondations ; semaine 2, élévation des murs du RDC ; semaine 3, coffrage, ferraillage (fer HA 500) et bétonnage (béton C25/30) de la dalle ; semaine 4, finitions et observations.
        \item \emph{Partie 3 :} compétences acquises (lecture de plans, dosage du béton, mise en œuvre du ferraillage) ; difficultés (intempéries, maniement du vibreur) ; respect des normes de sécurité (port du casque, balisage, grue à tour).
    \end{itemize}
    \item \textbf{Conclusion :} bilan personnel positif, apports pour la formation, deux recommandations (ex. : renforcer l'équipement de protection individuelle ; allonger la durée du stage ou multiplier les visites de chantier encadrées par l'école), formule d'envoi respectueuse au Proviseur.
\end{enumerate}
\textbf{Points de vigilance :} ton impersonnel (pas de « je ») ; concordance des temps (passé composé/imparfait pour les faits) ; vocabulaire technique exact ; ordre chronologique strict dans la partie 2 ; la signature se place \emph{après} la conclusion.
\end{corrige}

\begin{corrige}[Exercice 10 — Introduction du rapport d'enquête (insalubrité)]
\textbf{Barème indicatif (sur 10) :} contexte/déclencheur : 2 pts ; objet précis : 2 pts ; méthodologie : 2 pts ; plan annoncé : 2 pts ; ton neutre et langue : 2 pts.

\textbf{Introduction rédigée (proposition) :}
« Suite à une plainte collective des riverains du quartier Mfoundi, parvenue à la Mairie d'Arrondissement de Yaoundé 2, une mission d'enquête a été diligentée sur l'état d'insalubrité du Marché Mokolo. Le présent rapport a pour objet l'examen des conditions d'hygiène et d'assainissement dudit marché, ainsi que la formulation de mesures correctives. L'enquête s'est déroulée les 10, 11 et 12 septembre 2025. La méthodologie adoptée a combiné l'observation directe des lieux (documentée par des photographies), la collecte de témoignages auprès des commerçants, des usagers et du personnel de santé du dispensaire du marché, et l'examen des données sanitaires disponibles. Les constats effectués font état d'ordures ménagères non collectées depuis huit jours bloquant l'entrée principale, d'eaux usées stagnantes le long des étals de poissonnerie, de la présence de rongeurs en plein jour et d'une hausse de 40\,\% des consultations pour diarrhées aiguës et infections cutanées. Le présent rapport présentera d'abord les constats dressés sur le terrain, puis analysera les causes de la situation, avant de formuler des recommandations opérationnelles. »

\textbf{Points de vigilance :} ne pas développer les constats en détail dans l'introduction (ils appartiennent au corps) ; employer le passif ou « nous » ; respecter la concordance des temps ; annoncer un plan en deux ou trois parties maximum.
\end{corrige}

\begin{corrige}[Exercice 11 — Rapport d'activité du Club « Jeunes Écologistes »]
\textbf{Barème indicatif (sur 20) :} analyse de la situation : 4 pts ; plan détaillé : 4 pts ; corps rédigé : 8 pts ; conclusion : 3 pts ; langue : 1 pt.

\textbf{1. Analyse de la situation de communication :}
\begin{itemize}
    \item \textbf{Destinateur :} le Président du Club « Jeunes Écologistes » (au nom des 60 adhérents).
    \item \textbf{Destinataire :} le Proviseur, pour le Conseil d'Administration.
    \item \textbf{Objet :} rapport d'activités du 1er trimestre (septembre--décembre 2025) du Club.
    \item \textbf{Finalité :} rendre compte des actions menées et de la gestion financière, justifier l'utilisation des fonds, obtenir le maintien du soutien de l'administration.
    \item \textbf{Tonalité attendue :} objective, neutre, administrative (ton impersonnel, faits chiffrés).
\end{itemize}
\textbf{2. Plan détaillé (structure interne) :}
\begin{itemize}
    \item \emph{Introduction :} contexte, objet, période couverte, méthode, annonce du plan.
    \item \emph{Partie 1 : Activités de structuration et de sensibilisation} (1.1. Assemblée générale du 15/09 ; 1.2. Journée « Lycée Propre » du 05/10 ; 1.3. Conférence du 18/10).
    \item \emph{Partie 2 : Activités de production et de découverte} (2.1. Pépinière du 12/11 ; 2.2. Sortie pédagogique du 03/12).
    \item \emph{Partie 3 : Gestion financière et bilan} (3.1. Recettes ; 3.2. Dépenses ; 3.3. Solde et appréciation).
    \item \emph{Conclusion :} bilan moral et financier, perspectives, recommandations, formule d'envoi.
\end{itemize}
\textbf{3. Corps rédigé (proposition condensée) :}
« \textbf{Partie 1 :} Le trimestre s'est ouvert par l'assemblée générale de rentrée du 15 septembre, qui a réuni 45 adhérents sur 60 (taux de participation de 75\,\%), avec l'élection du bureau et la fixation de la cotisation à 1000 F. Le 5 octobre, la journée « Lycée Propre » a permis de ramasser 50 kg de déchets plastiques et de sensibiliser 200 élèves. Le 18 octobre, la conférence « Changement climatique et BTP », animée par l'Ingénieur NGONO, a accueilli 120 participants, l'amphithéâtre affichant complet.
\textbf{Partie 2 :} Le 12 novembre, en partenariat avec le MINEPDED, le club a mis en sachets 500 plants (acacia et eucalyptus) destinés au reboisement. Le 3 décembre, une sortie pédagogique au Parc de la Mefou a réuni 30 adhérents pour une étude de la biodiversité.
\textbf{Partie 3 :} Les recettes s'élèvent à 65\,000 F (cotisations : 50\,000 F ; subvention municipale : 15\,000 F). Les dépenses atteignent 48\,500 F (transport, plants, affiches, goûter). Le solde créditeur est de 16\,500 F, reporté sur le 2e trimestre. La gestion est donc saine et excédentaire. »
\textbf{Conclusion (proposition) :} « Le premier trimestre se solde par un bilan moral positif (cinq activités réalisées, forte mobilisation) et un bilan financier excédentaire de 16\,500 F. Pour le second trimestre, le club envisage la plantation effective des plants et une campagne de tri des déchets. Il est recommandé à l'administration de maintenir son appui logistique et de mettre à disposition un local de stockage du matériel. Nous prions Monsieur le Proviseur d'agréer l'expression de notre respectueuse considération. »

\textbf{Points de vigilance :} vérifier la cohérence des chiffres ($65\,000 - 48\,500 = 16\,500$ F) ; distinguer bilan moral et bilan financier ; les perspectives doivent rester réalistes et liées aux activités du club ; formule d'envoi adaptée à un supérieur hiérarchique.
\end{corrige}

\newpage

\coursec{Fiche bilan}

\begin{popfigure}
\begin{center}
\begin{tikzpicture}[mindmap, grow cyclic,
every node/.style={concept, text=white, font=\small\bfseries},
root concept/.append style={concept color=popdark, font=\bfseries},
level 1/.append style={level distance=3.2cm, sibling angle=51.5},
level 2/.append style={level distance=1.8cm, sibling angle=45, font=\scriptsize}]
\node[root concept]{Le rapport}
child[concept color=popblue]{node{Définition et types}
child{node[concept color=popblueL, text=popdark]{compte rendu}}
child{node[concept color=popblueL, text=popdark]{rapport de stage}}}
child[concept color=popgreen]{node{Structure externe}
child{node[concept color=popgreenL, text=popdark]{titre, date, lieu}}
child{node[concept color=popgreenL, text=popdark]{signature}}}
child[concept color=poporange]{node{Structure interne}
child{node[concept color=poporangeL, text=popdark]{introduction}}
child{node[concept color=poporangeL, text=popdark]{corps}}
child{node[concept color=poporangeL, text=popdark]{conclusion}}}
child[concept color=poppurple]{node{Exploiter les documents}
child{node[concept color=poppurpleL, text=popdark]{textes fonctionnels}}
child{node[concept color=poppurpleL, text=popdark]{textes iconiques}}}
child[concept color=poppink]{node{Tonalités}
child{node[concept color=poppinkL, text=popdark]{satirique}}
child{node[concept color=poppinkL, text=popdark]{tragique}}}
child[concept color=popgold]{node{De la situation au plan}
child{node[concept color=popgoldL, text=popdark]{analyser la situation}}
child{node[concept color=popgoldL, text=popdark]{rédiger et améliorer}}}
child[concept color=popmuted]{node{Compte rendu et remédiation}};
\end{tikzpicture}
\end{center}
\legende{Carte mentale de la leçon : les sept branches du rapport (définition, structures, documents, tonalités, méthode de rédaction, évaluation).}
\end{popfigure}

\begin{bilan}
\textbf{Définitions clés}
\begin{itemize}
\item Le \cle{rapport} : texte fonctionnel et utilitaire qui rend compte de manière objective, claire et ordonnée d'une mission, d'une enquête, d'une visite, d'un stage ou d'un événement, à destination d'un supérieur ou d'une institution.
\item \cle{Structure externe} : les éléments qui encadrent le texte (titre, lieu et date, nom du destinataire, nom et signature du rédacteur, éventuelles annexes).
\item \cle{Structure interne} : l'organisation du contenu en trois parties : introduction, corps, conclusion.
\item Tonalité \cle{satirique} : ton qui raille et critique les défauts ou les vices par l'ironie, l'exagération, la caricature.
\item Tonalité \cle{tragique} : ton qui suscite la pitié et la terreur face à un destin fatal et inévitable.
\end{itemize}

\textbf{Les trois parties du rapport à connaître par cœur}
\begin{center}
\begin{tabularx}{\linewidth}{|l|X|X|}
\hline
\rowcolor{popblueL} \textbf{Partie} & \textbf{Contenu} & \textbf{Questions auxquelles elle répond} \\
\hline
Introduction & contexte, circonstances, objet du rapport, annonce du plan & Qui ? Quand ? Où ? Pourquoi ? \\
\hline
Corps & faits observés, données, analyses, développement en paragraphes liés & Que s'est-il passé ? Qu'ai-je constaté ? \\
\hline
Conclusion & bilan, appréciation, suggestions ou recommandations & Qu'en conclure ? Que proposer ? \\
\hline
\end{tabularx}
\end{center}

\textbf{Méthodes express}
\begin{itemize}
\item \cle{Analyser la situation} : identifier le destinataire, l'objet, les faits et les documents utiles (textes fonctionnels : lettres, comptes rendus ; textes iconiques : images, tableaux, schémas).
\item \cle{Élaborer le plan} : classer les idées, ordonner les faits (ordre chronologique ou thématique), rédiger l'introduction et la conclusion en dernier si besoin.
\item \cle{Rédiger} : phrases claires, temps du passé (passé composé, imparfait), connecteurs logiques (d'abord, ensuite, enfin, par conséquent), ton neutre et objectif.
\item \cle{Améliorer} : confronter sa production à celle des autres, vérifier l'orthographe, la cohérence et le respect de la structure.
\end{itemize}
\end{bilan}

\begin{aretenir}[Checklist avant le Bac]
\begin{itemize}
\item[$\square$] Je sais définir le rapport et citer ses types (rapport de stage, de visite, d'enquête).
\item[$\square$] Je connais les éléments de la structure externe (titre, date, lieu, destinataire, signature).
\item[$\square$] Je sais que la structure interne comporte toujours introduction, corps et conclusion.
\item[$\square$] Je sais rédiger une introduction qui présente le contexte et l'objet du rapport.
\item[$\square$] Je sais organiser le corps en paragraphes liés par des connecteurs logiques.
\item[$\square$] Je sais rédiger une conclusion avec bilan et recommandations.
\item[$\square$] Je sais exploiter un texte fonctionnel et un texte iconique pour alimenter mon rapport.
\item[$\square$] Je distingue la tonalité satirique (ironie, critique) de la tonalité tragique (pitié, fatalité).
\item[$\square$] Je sais analyser une situation donnée pour élaborer le plan de mon rapport.
\item[$\square$] Je relis et j'améliore ma production (orthographe, cohérence, structure).
\item[$\square$] Je rédige de manière objective, claire et dans le respect du destinataire.
\end{itemize}
\end{aretenir}

\findecours

\end{document}
